%% 
%% Copyright 2019-2024 Elsevier Ltd
%% 
%% This file is part of the 'CAS Bundle'.
%% --------------------------------------
%% 
%% It may be distributed under the conditions of the LaTeX Project Public
%% License, either version 1.3c of this license or (at your option) any
%% later version.  The latest version of this license is in
%%    http://www.latex-project.org/lppl.txt
%% and version 1.3c or later is part of all distributions of LaTeX
%% version 1999/12/01 or later.
%% 
%% The list of all files belonging to the 'CAS Bundle' is
%% given in the file `manifest.txt'.
%% 
%% Template article for cas-sc documentclass for 
%% double column output.

\documentclass[a4paper,fleqn]{cas-sc}

% If the frontmatter runs over more than one page
% use the longmktitle option.

%\documentclass[a4paper,fleqn,longmktitle]{cas-sc}

%\usepackage[numbers]{natbib}
%\usepackage[authoryear]{natbib}
\usepackage{amsmath,amssymb,bm,mathtools,amsthm}
\usepackage{graphicx}
\usepackage{enumitem}
\usepackage{placeins}
\usepackage{hyperref}
 
\newtheorem{theorem}{Theorem}[section]
\newtheorem{lemma}[theorem]{Lemma}
\newtheorem{corollary}[theorem]{Corollary}
\newtheorem{proposition}[theorem]{Proposition}
\theoremstyle{remark}
\newtheorem{remark}[theorem]{Remark}
 
\newcommand{\eps}{\varepsilon}
\newcommand{\epsM}{\eps_{M}}
\newcommand{\epsB}{\eps_{B}}
\newcommand{\Bb}{\mathbf{B}}
\newcommand{\FF}{\mathbf{F}}
\newcommand{\GG}{\mathbf{G}}
\newcommand{\RR}{\mathbf{R}}
\newcommand{\LL}{\mathbf{L}}
\newcommand{\AAt}{\tilde{\mathbf{A}}}
\newcommand{\half}{\tfrac{1}{2}}
\newcommand{\mm}{\mathbf{m}}
\DeclareMathOperator{\minmod}{minmod}
\DeclareMathOperator{\median}{median}
% Floats. In cas-sc a figure/table without options may only go at the top
% of a page ([t]), and for tables "[htbp]" is silently ignored. Give every
% float its placement through the class key, e.g. \begin{figure}[pos=htbp];
% full-page figures use [pos=p].
\renewcommand{\topfraction}{0.9}
\renewcommand{\bottomfraction}{0.8}
\renewcommand{\textfraction}{0.1}
\renewcommand{\floatpagefraction}{0.6}
\renewcommand{\topfraction}{0.9}
\renewcommand{\bottomfraction}{0.9}
\renewcommand{\textfraction}{0.05}
\renewcommand{\floatpagefraction}{0.75}
\renewcommand{\floatpagefraction}{0.6}
\renewcommand{\textfraction}{0.05}
\renewcommand{\topfraction}{0.95}
%%%Author macros
\def\tsc#1{\csdef{#1}{\textsc{\lowercase{#1}}\xspace}}
\tsc{WGM}
\tsc{QE}
%%%

% Uncomment and use as if needed
%\newtheorem{theorem}{Theorem}
%\newtheorem{lemma}[theorem]{Lemma}
%\newdefinition{rmk}{Remark}
%\newproof{pf}{Proof}
%\newproof{pot}{Proof of Theorem \ref{thm}}

\begin{document}
\let\WriteBookmarks\relax
\def\floatpagepagefraction{1}
\def\textpagefraction{.001}

% Short title
\shorttitle{}    

% Short author
\shortauthors{}  

% Main title of the paper
\title [mode = title]{A high order shock capturing IMEX finite volume scheme with Mach- and Alfvén-adaptive flux splitting for ideal MHD}  

% Title footnote mark
% eg: \tnotemark[1]
\tnotemark[1] 

% Title footnote 1.
% eg: \tnotetext[1]{Title footnote text}
\tnotetext[1]{} 

% First author
%
% Options: Use if required
% eg: \author[1,3]{Author Name}[type=editor,
%       style=chinese,
%       auid=000,
%       bioid=1,
%       prefix=Sir,
%       orcid=0000-0000-0000-0000,
%       facebook=<facebook id>,
%       twitter=<twitter id>,
%       linkedin=<linkedin id>,
%       gplus=<gplus id>]

\author[1]{}%[<options>]

% Corresponding author indication
\cormark[1]

% Footnote of the first author
\fnmark[1]

% Email id of the first author
\ead{}

% URL of the first author
\ead[url]{}

% Credit authorship
% eg: \credit{Conceptualization of this study, Methodology, Software}
\credit{}

% Address/affiliation
\affiliation[1]{organization={},
            addressline={}, 
            city={},
%          citysep={}, % Uncomment if no comma needed between city and postcode
            postcode={}, 
            state={},
            country={}}

\author[2]{}%[]

% Footnote of the second author
\fnmark[2]

% Email id of the second author
\ead{}

% URL of the second author
\ead[url]{}

% Credit authorship
\credit{}

% Address/affiliation
\affiliation[2]{organization={},
            addressline={}, 
            city={},
%          citysep={}, % Uncomment if no comma needed between city and postcode
            postcode={}, 
            state={},
            country={}}

% Corresponding author text
\cortext[1]{Corresponding author}

% Footnote text
\fntext[1]{}

% For a title note without a number/mark
%\nonumnote{}

% Here goes the abstract
\begin{abstract}
We present an adaptive implicit--explicit (IMEX) finite volume scheme for the
ideal magnetohydrodynamics (MHD) equations that adapts to the local Mach and
Alfv\'en number regime. The flux is decomposed into advective,
acoustic and magnetic parts, and the fractions of the acoustic and of the
magnetic part that are treated explicitly are governed by two independent local
parameters, $\epsM$ and $\epsB$, driven by the acoustic and the Alfv\'en Mach
number respectively. Each stiffness is thus transferred to the implicit solver
independently: the channel that is not stiff is returned to an upwind Riemann
solver instead of being carried by the implicit operator, and the explicit time
step is freed from the sound-speed and from the Alfv\'en-speed constraint
separately, while in the supersonic and super-Alfv\'enic regime the scheme
reduces to a fully explicit Godunov-type method. A characteristic analysis gives
the Alfv\'en speeds from a quadratic in $\epsB$ alone and the magnetosonic speeds
from a quartic in both parameters, and shows the explicit sub-system to have real eigenvalues at every value of
the two parameters---which is what makes the explicit time step and the
wave-speed envelope of the Riemann solver well defined---and a complete set of
right eigenvectors whenever $\epsM>0$. The latter is what licenses the third-order characteristic WENO
reconstruction with a monotonicity-preserving limiter used for the explicit
terms, and it licenses only that, and not uniformly: the basis is evaluated
with $\epsM$ bounded from below by $10^{-2}$, while the flux is not.
Independently of it, and at every value of the parameters, an exact block
decomposition of the split system confines the singular limits of its
eigenvectors to low-dimensional sub-blocks: the Alfv\'en fields span a
two-dimensional invariant subspace whose dual functionals annihilate every
other wave family, so that their projection is available in closed form. On
this rests an HLLEM Riemann solver that resolves the contact and the rotational
discontinuities sharply even in the stiff limit, where the complete
characteristic basis is lost. On
the implicit side, the kinetic and magnetic energies are linearised consistently
with the partitioned IMEX Runge--Kutta scheme and an unstaggered,
potential-based constrained transport is written as a linear map of the
momentum; both are embedded in the pressure operator, so that each stage
solves for a magnetic predictor and then for the pressure, without an outer
iteration between the two, with a magnetic field that is divergence-free to
machine precision and with the total energy conserved without re-solving for
the pressure. Pressure positivity is enforced by a nested Newton iteration on
a second-order M-matrix operator, the fourth-order solution being recovered by
defect correction. Numerical experiments on standard benchmarks, from the
near-incompressible to the supersonic regime, confirm the design third order of
the scheme---approached from below at unit background density and exceeded in
the stiff regime---together with the asymptotic-preserving behaviour, a time
step larger than the explicit one by four orders of magnitude at the lowest
Mach number considered, positivity of density and pressure, and a magnetic
field divergence-free to round-off.
\end{abstract}

% Use if graphical abstract is present
%\begin{graphicalabstract}
%\includegraphics{}
%\end{graphicalabstract}

% Research highlights
\begin{highlights}
\item Flux splitting with two parameters driven by acoustic and Alfv\'en Mach numbers
\item Explicit sub-system proved hyperbolic at every parameter value, with complete eigenvectors whenever $\epsM>0$
\item Exact block decomposition gives the Alfv\'en projection in closed form, so an HLLEM solver resolves the rotational discontinuities even in the stiff limit
\item Constrained transport embedded in the pressure operator, with no outer iteration between the magnetic and the pressure sub-system
\item Each stiffness removed independently: the stiff channel leaves the time step, the non-stiff one keeps an upwind Riemann solver; third-order in space and time
\end{highlights}


% Keywords
% Each keyword is seperated by \sep
\begin{keywords}
Ideal magnetohydrodynamics \sep
All Mach and Alfv\'en numbers \sep
Semi-implicit IMEX finite volume schemes \sep
Adaptive flux splitting \sep
Divergence-free \sep
Third-order in space and time \sep
shock capturing
\end{keywords}

\maketitle

% Main text
\section{Introduction}\label{sec:intro}

The equations of ideal magnetohydrodynamics (MHD) describe an electrically
conducting fluid coupled to a magnetic field and form a first-order nonlinear system
of non-strictly hyperbolic conservation laws. They arise in astrophysical plasma
dynamics~\cite{GoedbloedPoedts2019,Priest2014} and in inertial and magnetic
confinement fusion~\cite{Jardin2010,AtzeniMeyer2004,Slutz2010,Sinars2020}, where a
magnetically dominated plasma may evolve rapidly from a near-equilibrium
configuration into a fully time-dependent compressible regime. What distinguishes
MHD from gas dynamics, for the purposes of numerical design, is that it is stiff in
\emph{two} independent ways: the acoustic Mach number $M$ controls the stiffness of the gas pressure, the
Alfv\'en Mach number $M_A$ that of the Maxwell stress (magnetic pressure and
tension). As
$M\to0$ the system approaches a weakly compressible regime governed by an elliptic
pressure balance; an analogous singular limit arises as $M_A\to0$, typical of
magnetic-pressure dominated plasmas. The two limits stem from different physical
mechanisms, are reached independently, and routinely coexist within a single
computation.

Density-based Godunov-type finite volume schemes~\cite{BrioWu1988,MiyoshiKusano2005}
face two difficulties in these limits. The time step is dictated by the fast
magnetosonic speed rather than by the material velocity, giving an $\mathcal O(1/M)$
or $\mathcal O(1/M_A)$ overhead beyond what the physically relevant transport
requires~\cite{Klein1995}; and the artificial dissipation of upwind fluxes scales
with that same speed, so that the asymptotic scaling $p=p_0+\mathcal O(M^{2})$ is
lost~\cite{GuillardViozat2017,Dellacherie2010}. Since the fast magnetosonic speed
degenerates to the Alfv\'en speed in the magnetic-pressure dominated regime, both
reappear as $M_A\to0$. Rescaling the flux
dissipation~\cite{barsukow2021truly,barsukow2023all} restores the correct scaling but
not the time step, which then calls for a fully implicit
integration~\cite{leidi2022finite,viallet2011towards} and hence for large
nonsymmetric nonlinear systems whose convergence is difficult to control.

Treating implicitly only the stiff terms, while keeping the nonlinear transport
explicit, avoids both difficulties. This is naturally realised within the class of
implicit--explicit (IMEX) Runge--Kutta integrators~\cite{RKIMEX,RKIMEX2} applied to a
flux-vector splitting: the implicit step reduces to linear or mildly nonlinear solves, and the time
step is constrained only by the eigenvalues of the explicit sub-system. Pressure-based
semi-implicit schemes of this kind were made conservative and shock-capturing by
Dumbser and Casulli~\cite{DumbserCasulli2016}, who combined them with the
Toro--V\'azquez splitting~\cite{TVSplit}, with high-order all-Mach extensions
following~\cite{BoscheriPareschi2021,TavelliDumbser2017,FVVEM}. For MHD the
development has proceeded by successively removing the stiff wave families from the
explicit step: Dumbser et al.~\cite{Dumbser2019} kept the pressure implicit and the
magnetic field explicit, so that the time step remains governed by the Alfv\'en
speed; removing that restriction requires the field to be implicit as well, which
Fambri~\cite{Fambri2021} achieved through a three-way splitting into advective,
acoustic and magnetic sub-systems, the structure also underlying the collocated
formulations of Dematt\'e et al.~\cite{Dematte2024} and of Boscheri and
Thomann~\cite{BoscheriThomann2024}. We adopt that three-way splitting; our point of
departure concerns not the splitting itself, but how its parts are assigned to the
explicit and the implicit sub-system.

In all of these schemes the two sub-systems coincide \emph{exactly} with the parts of
the splitting: the advective flux is explicit, the entire acoustic and magnetic
fluxes implicit. This is well suited to the stiff limits, and has two consequences
elsewhere. In supersonic and super-Alfv\'enic regions the pressure and magnetic
dynamics are still handled by the implicit solver, which relies on central-type
stencils rather than on upwind shock capturing, unless Rusanov-type dissipation is
explicitly reintroduced~\cite{DumbserCasulli2016,BoscheriThomann2024,MHDIMEX}. And, with those fluxes assigned
entirely to the implicit side, the explicit sub-system loses its spectral separation:
its eigenvalues collapse onto $\{0,u\}$ and the Jacobian admits no complete set of
eigenvectors, so it is only weakly hyperbolic. This precludes Riemann solvers built on the
eigenvectors of the explicit Jacobian---Roe~\cite{Roe1981}, and
HLLEM~\cite{einfeldt1991godunov,HLLEM} in its standard form, which projects the
jump onto the intermediate fields of that Jacobian---and leaves
scalar-dissipation fluxes of Rusanov type~\cite{Rusanov1962} as the usual
choice. Recovering those fields from a structure that survives the stiff limit,
rather than from the explicit Jacobian itself, is what Section~\ref{sec:hllem}
does. The consequence is
more severe than in gas dynamics: the degenerate spectrum retains no trace of the
rotational discontinuities either, so the Alfv\'en waves---which carry no acoustic
stiffness whatsoever---are smeared at the same rate as everything else.

We depart from this rigid identification. The fraction of the acoustic and of the
magnetic flux returned to the explicit sub-system is controlled by \emph{two
independent} adaptive parameters, $\varepsilon_M$ and $\varepsilon_B$, driven
respectively by the local acoustic and Alfv\'en Mach numbers. Letting a splitting
parameter respond to the Mach number has been explored for gas dynamics by Degond and
Tang~\cite{DegondTang2011} and by Iampietro and
collaborators~\cite{IampietroEtAl2018,IampietroEtAl2018b,IampietroEtAl2020}; a single
parameter, however, is necessarily sub-optimal for MHD whenever $M\neq M_A$, since
the two stiffness mechanisms are governed by distinct dimensionless numbers. It
either over-stabilises the non-stiff channel or leaves the stiff one partially
unresolved, and it is precisely the two \emph{mixed} regimes---a high-$\beta$
low-Mach plasma, whose Alfv\'en waves would be needlessly implicit, and a strongly
magnetised plasma at moderate Mach number, whose acoustic waves would be---that a
single parameter cannot represent. The dual-parameter splitting addresses the two stiffnesses simultaneously and
independently, and recovers the fully explicit Godunov discretisation where
neither channel is stiff. What it buys over a single parameter is not
symmetric in the two mixed regimes, and the asymmetry follows from the
definitions rather than from experiment. The conditions $\epsM\to0$ and
$\epsB\to1$ together force $c_B\lesssim\|\mathbf u\|\ll a$, hence
$\beta\gg1$: wherever the splitting returns the magnetic flux to the explicit
solver while the acoustic flux is implicit, the field is a weak perturbation of
the gas dynamics. The converse regime, $\epsM\to1$ with $\epsB\to0$, forces
$\beta\ll1$, and there the magnetic flux dominates the dynamics. The gain
accordingly sits in the second of the two.

Against a single parameter driven by the acoustic Mach number alone, which is
the direct transcription of the gas-dynamics
construction~\cite{DegondTang2011,IampietroEtAl2018}, a strongly magnetised
flow at moderate Mach number leaves the magnetic flux explicit and the time
step set by the Alfv\'en speed, which $\epsB$ removes---by orders of
magnitude, and over a large part of the domain: in the blast wave of
Section~\ref{sec:blast} the two parameters differ by more than one half over
about two fifths of the cells, with $\epsM\approx1$ and $\epsB\approx0$
there. Against a single parameter driven by whichever Mach number is the more
restrictive---the safe choice---the second parameter instead costs a little.
Both fluxes are implicit under the safe choice, so its admissible step is the
larger one: keeping the non-stiff flux explicit costs an explicit spectral
radius of at most $\|\mathbf u\|(1+\epsB/M_A)$ by~\eqref{eq:alfven_eig}, or
$\|\mathbf u\|(1+\epsM/M)$ in the mirror case. That factor is $1.13$ at the
strongest amplitude of the field-loop advection of Section~\ref{sec:fieldloop}
($M_A\approx7.9$) and $1.08$ in the Kelvin--Helmholtz test of
Section~\ref{sec:KH} ($M_A\approx11.8$); over the whole range
of~\eqref{eq:switch} it is bounded by $3.63$, attained just above the switch at
$M_A\approx0.38$, where the channel has already been returned to the explicit
solver while the Alfv\'en speed is still $2.7\|\mathbf u\|$. What is bought in
exchange is an upwind rather than a central discretisation of the non-stiff
flux. In two dimensions that exchange reaches the in-plane field only weakly,
since the in-plane field is advanced by the constrained transport of
Section~\ref{sec:CT} at either value of $\epsB$, which reaches it only through
the velocity; the measurement is reported in Section~\ref{sec:CT}.

Two structural results follow, and they serve different parts of the scheme.
The first concerns hyperbolicity, and it has two halves that are worth keeping
apart. The eight eigenvalues of the explicit Jacobian are real at every value
of both parameters, the stiff limit included; that is the weaker but
unconditional half, and it is what makes the explicit CFL
condition~\eqref{eq:cfl} and the wave-speed envelope of the Riemann solver well
defined, since both are built from the split spectrum. The stronger half is
conditional: whenever $\epsM>0$ the Jacobian also has a complete set of right
eigenvectors, whose closed form we derive, and this is what licenses the
characteristic reconstruction of the explicit terms. It is also all that the
complete basis licenses, and not uniformly: it degenerates as $\epsM\to0$, so
the reconstruction is carried out with the splitting parameter bounded away
from zero, while the flux itself is not.

The second result is what makes a wave-resolving Riemann solver possible at
\emph{every} value of the parameters, including the stiff limit in which the
complete basis is lost. The split system decomposes exactly: the divergence mode
is trivially decoupled, and the Alfv\'en fields span a two-dimensional invariant
subspace whose dual functionals annihilate every other wave family. Their
projection is therefore available in closed form at every value of the
parameters: the regularised pair is regular as $B_\perp\to0$, and at
$\epsB B_x=0$, where its normalisation degenerates, the block~\eqref{eq:A2} is
exactly diagonal and the pair is given by its decoupled limit---a degeneracy of
the parametrisation, not of the system. The two rotational discontinuities can
therefore be resolved where the degenerate spectrum of the standard splitting
retains no trace of them at all. We use this in an HLLEM flux of $\delta_k$-driven
asymptotic-preserving form, whose wave-speed envelope is built from the split
spectrum itself and whose entropy strength is taken in the classical form,
bounded uniformly in $\epsM$. Of the two singular limits of the closed-form
eigenvectors, only the remaining aligned sector requires a numerical inverse.

The order of accuracy is, by itself, a specification rather than a
contribution: high-order, asymptotic-preserving, divergence-free
discretisations of MHD have been built before, among them the finite-difference
WENO schemes with constrained transport of Chen, Wu and
Xiong~\cite{ChenWuXiong2023}. What has to be shown is that the splitting and
the two structural results above survive at that order. Third-order accuracy in
space on the explicit terms is obtained with a characteristic WENO3-Z
reconstruction~\cite{DonBorges2013, BORGES20083191} and a
monotonicity-preserving limiter~\cite{SureshHuynh1997}. For that limiter we
prove a result that is not a specification: under a non-degeneracy condition on
the critical points, it returns the unreconstructed WENO value unchanged on
every sufficiently fine mesh, so that the design order is retained at smooth
extrema wherever they fall relative to the mesh
(Proposition~\ref{prop:mp3}). Monotonicity-preserving limiters are usually
accepted with order loss at extrema; here they are not. The implicit terms are discretised
with a fourth-order finite volume operator family, and third-order accuracy in
time follows from a stiffly accurate IMEX Runge--Kutta
scheme~\cite{BoscheriPareschi2021}, the stiff accuracy being what guarantees
asymptotic consistency in both limits.

The second contribution concerns the implicit step. In the semi-implicit schemes
cited above the kinetic energy couples the pressure to the momentum
nonlinearly, which is resolved by a second pressure solve or by a Picard
iteration, and the solenoidal involution $\nabla\!\cdot\!\mathbf B=0$, which is
not inherited at the discrete level~\cite{Toth2000}, is restored by a separate
constrained-transport step~\cite{EvansHawley1988,Rossmanith2006}. Here the
kinetic and magnetic energies are linearised about the explicit argument of the
partitioned IMEX scheme, and the unstaggered, potential-based constrained
transport of Rossmanith~\cite{Rossmanith2006} is written as a linear map of the
momentum. Both enter the pressure operator, so that an implicit stage solves
for a magnetic predictor and then for the pressure, without iterating between
the two: the energy is balanced with a divergence-free field throughout the
solve, the total energy is conserved without a second solve, and the formal
order of the time integrator is retained. Positivity of the pressure is enforced with the
mechanism of~\cite{MHDIMEX}: the cell-local pressure relation is replaced by one
with a steep slope below a floor, which makes the pressure system nonlinear and
is solved by a nested Newton iteration. Because that mechanism requires an
operator with non-positive off-diagonal entries, which the fourth-order operator
does not have, it is applied to a second-order M-matrix operator and the
fourth-order solution is recovered by defect correction, with the constrained
transport, the induction of $B_z$ and the magnetic energy flux carried by the
defect. The nested Newton iteration therefore runs inside every defect step of
every computation reported here.

The remainder of this paper is organised as follows. Section~\ref{sec:governing}
introduces the governing equations, the three-way flux splitting and the
dual-parameter $\varepsilon$-split form, with a characteristic analysis of the
split system and its block decomposition. Section~\ref{sec:scheme} describes the
numerical method. Section~\ref{sec:result} assesses the scheme on benchmarks in
one and two space dimensions from the near-incompressible to the supersonic regime, and conclusions are drawn in Section~\ref{sec:Conclusion}.
% ===========================================================================
\section{System of governing equations}
\label{sec:governing}

In conservative form, the ideal MHD equations read
\begin{equation}
\label{eq:ideal_MHD}
\begin{cases}
\dfrac{\partial \rho}{\partial t}+ \nabla \cdot (\rho \mathbf{u}) = 0, \\[1.5ex]
\dfrac{\partial (\rho \mathbf{u})}{\partial t}+ \nabla \cdot \left( \rho \mathbf{u} \otimes \mathbf{u}
+ \left( p + \tfrac{1}{2}\|\Bb\|^2 \right)\mathbf{I}- \Bb \otimes \Bb \right) = \mathbf{0}, \\[1.5ex]
\dfrac{\partial E}{\partial t}+ \nabla \cdot \left( \left( E + p + \tfrac{1}{2}\|\Bb\|^2 \right)\mathbf{u}
- (\mathbf{u}\cdot\Bb)\,\Bb \right) = 0, \\[1.5ex]
\dfrac{\partial \Bb}{\partial t}+ \nabla \cdot \left( \Bb \otimes \mathbf{u}- \mathbf{u} \otimes \Bb \right) = \mathbf{0}.
\end{cases}
\end{equation}
Here $\rho$ denotes the fluid density, $\mathbf{u}=[u,v,w]^{\mathsf T}$ the velocity,
$\Bb=[B_x,B_y,B_z]^{\mathsf T}$ the magnetic field, and $E$ the total energy
per unit volume, the sum of internal, kinetic and magnetic contributions,
\begin{equation}
\label{eq:total_energy}
E=\rho e+k+m,\qquad
k=\half\rho\|\mathbf{u}\|^2,\qquad
m=\half\|\Bb\|^2,
\end{equation}
with $e$ the specific internal energy. The magnetic field is nondimensionalised
such that the vacuum magnetic permeability is unity.

In addition, the magnetic field must satisfy the divergence-free constraint
\begin{equation}
\label{eq:divB}
\nabla\cdot\Bb=0,
\end{equation}
which is a constraint of the system: if it holds at the initial time, the
evolution equations~\eqref{eq:ideal_MHD} preserve it for all later times at the
continuous level. At the discrete level this property is not automatically
inherited, and a dedicated treatment is required.

The system is closed by the ideal gas law,
\begin{equation}
\label{eq:eos}
e(p,\rho)=\frac{p}{(\gamma-1)\rho},
\end{equation}
with $\gamma$ the ratio of specific heats, so that the pressure depends linearly
on the internal energy. The specific enthalpy is $h=e+p/\rho$ and the sound speed
is
\begin{equation}
\label{eq:soundspeed}
a=\sqrt{\frac{\gamma p}{\rho}} .
\end{equation}

Letting $\mathbf{U}=(\rho,\rho\mathbf{u},E,\Bb)^{\mathsf T}$ denote the vector of conserved
quantities, the system~\eqref{eq:ideal_MHD} is written in compact form as
\begin{equation}
\label{eq:compact_MHD}
\partial_t\mathbf{U}+\nabla\cdot\boldsymbol{\mathcal F}(\mathbf{U})=\mathbf0,
\qquad \boldsymbol{\mathcal F}=(\FF,\GG).
\end{equation}
For one-dimensional flows aligned with the $x$-direction it reads
\begin{equation}
\label{eq:1DMHD}
\frac{\partial}{\partial t}
\begin{bmatrix}
\rho \\ \rho u \\ \rho v \\ \rho w \\ E \\ B_x \\ B_y \\ B_z
\end{bmatrix}
+
\frac{\partial}{\partial x}
\begin{bmatrix}
\rho u \\
\rho u^2+p+\half\|\Bb\|^2-B_x^2 \\
\rho uv-B_xB_y \\
\rho uw-B_xB_z \\
u\big(E+p+\half\|\Bb\|^2\big)-B_x(\mathbf{u}\cdot\Bb) \\
0 \\
uB_y-vB_x \\
uB_z-wB_x
\end{bmatrix}
=\mathbf 0 .
\end{equation}
This system has the eight eigenvalues
\begin{equation}
\label{eq:MHD_eigs}
\lambda_{1,8} = u \pm c_f, \qquad
\lambda_{2,7} = u \pm c_a, \qquad
\lambda_{3,6} = u \pm c_s, \qquad
\lambda_4 = u, \qquad
\lambda_5 = 0 ,
\end{equation}
with the Alfv\'en speed $c_a$ and the fast and slow magnetosonic speeds
$c_f$, $c_s$,
\begin{equation}
\label{eq:MHD_speeds}
c_a = \frac{|B_x|}{\sqrt{\rho}},\qquad
c_{f,s}^2 =\frac{1}{2}\left(c_B^2 + a^2\pm
\sqrt{\left(c_B^2 + a^2\right)^2-4 c_a^2 a^2}\right),
\qquad c_B=\frac{\|\Bb\|}{\sqrt{\rho}} .
\end{equation}

\paragraph{Two stiffness scales.}
The numerical stiffness of~\eqref{eq:ideal_MHD} is controlled by two
independent dimensionless numbers, the acoustic and the Alfv\'en Mach number,
\begin{equation}
\label{eq:two_mach}
M=\frac{\|\mathbf{u}\|}{a},\qquad
M_A=\frac{\|\mathbf{u}\|}{c_B} .
\end{equation}
When $M\ll1$ the spectral radius of the Jacobian is set by the fast magnetosonic
speed, forcing an $\mathcal O(1/M)$ time-step penalty relative to
material advection; when $M_A\ll1$ the Alfv\'en and magnetosonic speeds become
large, producing an \emph{independent} $\mathcal O(1/M_A)$ restriction.
 


\subsection{Three-way flux splitting}
\label{sec:3split}

The numerical method developed in this work is built upon the three-way flux splitting for MHD introduced by Fambri~\cite{Fambri2021} and later developed by Chen et al.~\cite{MHDIMEX}, in which the physical flux of the MHD system~\eqref{eq:ideal_MHD} is decomposed into advective, pressure and magnetic contributions,
\begin{equation}
\mathbf{F}(\mathbf{U}) = \mathbf{F}^c(\mathbf{U}) + \mathbf{F}^p(\mathbf{U}) + \mathbf{F}^B(\mathbf{U}),
\end{equation}
with
\begin{equation}
\mathbf{F}^c =
\begin{bmatrix}
\rho u \\
\rho u^2 \\
\rho uv \\
\rho uw \\
k u \\
0 \\
0 \\
0
\end{bmatrix},\quad
\mathbf{F}^p =
\begin{bmatrix}
0 \\
p \\
0 \\
0 \\
\rho h u \\
0 \\
0 \\
0
\end{bmatrix},\quad
\mathbf{F}^B =
\begin{bmatrix}
0 \\
\frac{1}{2}\|\mathbf{B}\|^2 - B_x^2 \\
-B_xB_y \\
-B_xB_z \\
\|\mathbf{B}\|^2u - B_x(\mathbf{u}\cdot\mathbf{B}) \\
0 \\
uB_y - vB_x \\
uB_z - wB_x
\end{bmatrix}.
\end{equation}


The three resulting sub-systems are referred to as the advection, pressure and magnetic sub-systems.

\paragraph{Advection sub-system.}
The convective flux $\mathbf{F}^c$ has eigenvalues
\begin{equation}
\label{eq:adv_eigs}
\lambda^c_{1,2,3,4}=0,\qquad \lambda^c_{5,6,7,8}=u,
\end{equation}
Thus, the sub-system is weakly hyperbolic, since one eigenvector associated with $\lambda^c$ = $u$ is missing. Its eigenstructure contains no trace of the acoustic, Alfvén or magnetosonic waves.

\paragraph{Pressure sub-system.}
The pressure flux $\mathbf{F}^p$ carries the eigenvalues
\begin{equation}
\label{eq:p_eigs}
\lambda^p_{1,2,3,4,5,6}=0,\qquad
\lambda^p_{7,8}=\frac{1}{2}\left(u\pm\sqrt{u^2+4a^2}\right).
\end{equation}
In the previous IMEX schemes for MHD~\cite{Fambri2021,MHDIMEX}
this sub-system is treated implicitly, which removes the acoustic stiffness
from the explicit time-step constraint.

\paragraph{Magnetic sub-system.}
The magnetic flux $\mathbf{F}^B$ carries the eigenvalues
\begin{equation}
\label{eq:B_eigs}
\lambda^B_{1,2,3,4}=0,\qquad
\lambda^B_{5,6}=\frac{1}{2}\left(u\pm\sqrt{u^2+4c_a^2}\right),\qquad
\lambda^B_{7,8}=\frac{1}{2}\left(u\pm\sqrt{u^2+4c_B^2}\right),
\end{equation}
This sub-system isolates the Alfvén and magnetic-tension stiffness.

In a standard IMEX framework \cite{MHDIMEX}, with the advective sub-system
treated explicitly and the pressure and magnetic sub-systems treated implicitly,
the weak hyperbolicity of $\mathbf{F}^c$ becomes a fundamental limitation:
without a complete set of eigenvectors no characteristic reconstruction can be
built for the explicit step, and no approximate Riemann solver that projects the
jump onto the eigenvectors of $\mathbf{F}^c$. A Rusanov-type flux is then the
usual choice, whose scalar dissipation acts on all components alike and smears
contact and shear waves. What the adaptive splitting of
Section~\ref{sec:eps_split} changes is that the intermediate fields keep a
trace in the explicit sub-system, and that trace is enough for a
wave-resolving flux even where the basis is still incomplete.
 
% ---------------------------------------------------------------------------
\subsection{Dual-parameter $\eps$-split form}
\label{sec:eps_split}

To overcome the deficiency identified at the end of Section~\ref{sec:3split},
a fraction of the acoustic and of the magnetic flux is returned to the
explicit sub-system, controlled by two independent parameters
$\epsM,\epsB\in[0,1]$:
\begin{equation}
\label{eq:eps_split}
\FF(\mathbf{U})=\tilde\FF^{c}(\mathbf{U})+\tilde\FF^{p}(\mathbf{U})+\tilde\FF^{B}(\mathbf{U}),
\end{equation}
where the modified convective flux $\tilde\FF^{c}$, the modified pressure flux
$\tilde\FF^{p}$ and the modified magnetic flux $\tilde\FF^{B}$ are defined as
\begin{equation}
\label{eq:new_split}
\begin{aligned}
&\tilde\FF^{c}=\FF^{c}(\mathbf{U})+\epsM\FF^{p}(\mathbf{U})+\epsB\FF^{B}(\mathbf{U})=
\begin{bmatrix}
\rho u\\[2pt]
\rho u^{2}+\epsM\,p+\epsB\!\left(\half\|\Bb\|^{2}-B_x^{2}\right)\\[2pt]
\rho uv-\epsB B_xB_y\\[2pt]
\rho uw-\epsB B_xB_z\\[2pt]
ku+\epsM\,\rho h u+\epsB\!\left(\|\Bb\|^{2}u-B_x(\mathbf{u}\cdot\Bb)\right)\\[2pt]
0\\[2pt]
\epsB\!\left(uB_y-vB_x\right)\\[2pt]
\epsB\!\left(uB_z-wB_x\right)
\end{bmatrix},
\\[10pt]
&\tilde\FF^{p}=(1-\epsM)\,\FF^{p}(\mathbf{U})=
\begin{bmatrix}
0\\[2pt]
(1-\epsM)\,p\\[2pt]
0\\[2pt]
0\\[2pt]
(1-\epsM)\,\rho h u\\[2pt]
0\\[2pt]
0\\[2pt]
0
\end{bmatrix},
% \\[10pt]
\tilde\FF^{B}=(1-\epsB)\,\FF^{B}(\mathbf{U})=
\begin{bmatrix}
0\\[2pt]
(1-\epsB)\!\left(\half\|\Bb\|^{2}-B_x^{2}\right)\\[2pt]
-(1-\epsB)\,B_xB_y\\[2pt]
-(1-\epsB)\,B_xB_z\\[2pt]
(1-\epsB)\!\left(\|\Bb\|^{2}u-B_x(\mathbf{u}\cdot\Bb)\right)\\[2pt]
0\\[2pt]
(1-\epsB)\!\left(uB_y-vB_x\right)\\[2pt]
(1-\epsB)\!\left(uB_z-wB_x\right)
\end{bmatrix},
\end{aligned}
\end{equation}


For every value of $\epsM$ and $\epsB$ the three fluxes sum to the full flux:
the parameters only redistribute the acoustic and the magnetic contributions
between the explicit and the implicit part, without altering the total. The advective sub-system is treated explicitly,
while the residual fluxes $\tilde\FF^{p}$ and $\tilde\FF^{B}$ are integrated
implicitly. The parameters $\epsM,\epsB$ are defined locally and adapt to the
flow conditions as described in Section~\ref{sec:eps_strategy}.
 


\subsubsection{Eigenstructure of the explicit sub-system}
\label{sec:eigenstructure}

We derive the eigenstructure of $\AAt=\partial\tilde\FF^{c}/\partial\mathbf{U}$    from the quasilinear form in primitive variables
$\mathbf Q=(\rho,u,v,w,p,B_x,B_y,B_z)^{\mathsf T}$. Linearising
\eqref{eq:new_split} about a state, using $a^2=\gamma p/\rho$ and the
divergence-free constraint~\eqref{eq:divB}, which in one dimension reads $\partial_xB_x=0$ and
therefore annihilates the $B_x$ column, the explicit
sub-system $\partial_t\mathbf U+\tilde{\mathbf A}\,\partial_x\mathbf U=0$
is equivalent to $\partial_t\mathbf Q+\tilde{\mathbf A}_p\,\partial_x\mathbf Q=0$ with 
\begin{equation}
\label{eq:prim_jac}
\AAt_{p}=
\begin{bmatrix}
u & \rho & 0 & 0 & 0 & 0 & 0 & 0\\[2pt]
0 & u & 0 & 0 & \tfrac{\epsM}{\rho} & 0 & \tfrac{\epsB B_y}{\rho} & \tfrac{\epsB B_z}{\rho}\\[2pt]
0 & 0 & u & 0 & 0 & 0 & -\tfrac{\epsB B_x}{\rho} & 0\\[2pt]
0 & 0 & 0 & u & 0 & 0 & 0 & -\tfrac{\epsB B_x}{\rho}\\[2pt]
0 & \epsM\rho a^2 & 0 & 0 & \epsM u & 0 & 0 & 0\\[2pt]
0 & 0 & 0 & 0 & 0 & 0 & 0 & 0\\[2pt]
0 & \epsB B_y & -\epsB B_x & 0 & 0 & 0 & \epsB u & 0\\[2pt]
0 & \epsB B_z & 0 & -\epsB B_x & 0 & 0 & 0 & \epsB u
\end{bmatrix}.
\end{equation}

The two scalings act on distinct couplings: $\epsM$ multiplies the
pressure--velocity acoustic block (rows/columns $u,p$), while $\epsB$ multiplies the velocity--field magnetic block (rows/columns
$u,v,w,B_y,B_z$). The advective
self-transport $u\,\partial_x$ of $\rho$ and of the transverse momenta is carried
by $\FF^c$ and is therefore unscaled; the transport of $p$ and of $\Bb$,
which lives in $\FF^p$ and $\FF^B$, is scaled by $\epsM$ and $\epsB$ respectively,
producing the diagonal entries $\epsM u$ and $\epsB u$.

\paragraph{Decoupling.}
The characteristic equation $\det(\AAt_p-\lambda\mathbf I)=0$ factorises into four
groups, exactly as in full MHD. The vanishing row for $B_x$ gives the divergence
mode $\lambda_5=0$. The density couples only to $u$, yielding the entropy wave
$\lambda_4=u$ with $\delta p=\delta\mathbf{u}=\delta\Bb=0$. The transverse components
split, under the rotation
\begin{equation}
\label{eq:beta}
\beta_y=\frac{B_y}{B_\perp},\quad\beta_z=\frac{B_z}{B_\perp},\quad
B_\perp=\sqrt{B_y^2+B_z^2},
\end{equation}
into a component perpendicular to $\Bb_\perp$ (Alfv\'en) and a component aligned
with $\Bb_\perp$ (magnetosonic).

\paragraph{Alfv\'en waves.}
The perpendicular combinations $v_\perp=\beta_z v-\beta_y w$,
$B_\perp^{\!*}=\beta_z B_y-\beta_y B_z$ satisfy a closed $2\times2$ system. The
$\partial_x u$ coupling drops, since $\beta_z B_y-\beta_y B_z=0$, leaving
\begin{equation}
(u-\lambda)\,v_\perp-\frac{\epsB B_x}{\rho}\,B_\perp^{\!*}=0,\qquad
-\epsB B_x\,v_\perp+(\epsB u-\lambda)\,B_\perp^{\!*}=0,
\end{equation}
whose determinant gives the Alfv\'en dispersion relation
\begin{equation}
\label{eq:alfven_disp}
(u-\lambda)(\epsB u-\lambda)-\epsB^2 c_a^2=0
\;\Longleftrightarrow\;
\lambda^2-(1+\epsB)u\,\lambda+\epsB u^2-\epsB^2 c_a^2=0,
\end{equation}
that is,
\begin{equation}
\label{eq:alfven_eig}
\tilde\lambda_A^{\pm}=\half\!\left[(1+\epsB)u\pm\Delta_A\right],\qquad
\Delta_A=\sqrt{(1-\epsB)^2u^2+4\epsB^2c_a^2}.
\end{equation}
The Alfv\'en speeds depend on $\epsB$ \emph{alone}: the magnetic parameter
controls them independently of the acoustic parameter.

\paragraph{Magnetosonic waves.}
The aligned combinations $v_\parallel=\beta_y v+\beta_z w$,
$B_\parallel=\beta_y B_y+\beta_z B_z=B_\perp$, together with $u$ and $p$, satisfy a
$4\times4$ system. Writing $\alpha=u-\lambda$, $\beta=\epsM u-\lambda$,
$\phi=\epsB u-\lambda$, the $u$--$p$ acoustic block and the
$v_\parallel$--$B_\parallel$ magnetic block read
\begin{equation}
\begin{bmatrix}\alpha & \epsM/\rho\\ \epsM\rho a^2 & \beta\end{bmatrix}
\!\begin{bmatrix}\delta u\\ \delta p\end{bmatrix}
+\begin{bmatrix}\epsB B_\perp/\rho\\ 0\end{bmatrix}\!\delta B_\parallel=0,
\qquad
\alpha\,\delta v_\parallel-\frac{\epsB B_x}{\rho}\delta B_\parallel=0,
\quad
\epsB B_\perp\,\delta u-\epsB B_x\,\delta v_\parallel+\phi\,\delta B_\parallel=0.
\end{equation}
Eliminating $\delta p$ from the acoustic block and $\delta v_\parallel$ from the
tension relation introduces the determinants
\begin{equation}
\label{eq:DsDA}
D_s=\alpha\beta-\epsM^2a^2,\qquad
D_A=\alpha\phi-\epsB^2c_a^2,
\end{equation}
and reduces the system to
\begin{equation}
\label{eq:ms_reduced}
D_s\,\delta u+\frac{\beta\,\epsB B_\perp}{\rho}\,\delta B_\parallel=0,
\qquad
\alpha\,\epsB B_\perp\,\delta u+D_A\,\delta B_\parallel=0,
\end{equation}
whose determinant vanishes precisely when
\begin{equation}
\label{eq:ms_disp}
D_s\,D_A=\epsB^2\,b_\perp^2\,\alpha\beta,\qquad b_\perp^2=\frac{B_y^2+B_z^2}{\rho},
\end{equation}
which is the magnetosonic quartic in $\lambda$. At $u=0$ one has
$\alpha=\beta=\phi=-\lambda$, $\alpha\beta=\lambda^2$, and~\eqref{eq:ms_disp}
reduces, using $c_a^2+b_\perp^2=c_B^2$, to the biquadratic
\begin{equation}
\label{eq:biquad}
\lambda^4-(\epsM^2a^2+\epsB^2c_B^2)\lambda^2+\epsM^2a^2\,\epsB^2c_a^2=0,
\end{equation}
with roots
\begin{equation}
\label{eq:ms_eig}
\tilde\lambda_{f,s}^{\pm}=\pm\frac{1}{\sqrt2}
\sqrt{\epsM^2a^2+\epsB^2c_B^2\pm\sqrt{(\epsM^2a^2+\epsB^2c_B^2)^2-4\epsM^2a^2\epsB^2c_a^2}}.
\end{equation}

\paragraph{Limits.}
\begin{itemize}
\item \textbf{As $\epsM,\epsB\to0$:} the Alfv\'en pair~\eqref{eq:alfven_eig}
tends to $\half(u\pm|u|)$ and the quartic~\eqref{eq:ms_disp} degenerates to
$\lambda^2(u-\lambda)^2=0$, so that the spectrum collapses to
\begin{equation}
\tilde\lambda^{c}\in\{0^{(4)},\,u^{(4)}\},
\end{equation}
recovering the weakly hyperbolic advective sub-system~\eqref{eq:adv_eigs} of the standard
three-way splitting.

\item \textbf{As $\epsM=\epsB\to1$:} the two relations recover the full MHD
spectrum $\{u\pm c_f,\;u\pm c_a,\;u\pm c_s,\;u,\;0\}$
of~\eqref{eq:MHD_eigs}.
\end{itemize}

\begin{theorem}[Hyperbolicity of the explicit sub-system]\label{thm:hyp}
Let $\rho>0$, $p>0$ and let $\varepsilon_M,\varepsilon_B\in[0,1]$ be frozen.
Then the matrix $\tilde{\mathbf A}_p$ of \eqref{eq:prim_jac} has the following properties.
\begin{enumerate}
\item[(i)] All eight eigenvalues are real.
\item[(ii)] If $\varepsilon_M>0$, $\tilde{\mathbf A}_p$ has a complete set of
real eigenvectors, for every $\varepsilon_B\in[0,1]$ and every state,
including $B_x=0$, $B_\perp=0$ and coinciding wave speeds.
\item[(iii)] If $\varepsilon_M=0$, $\tilde{\mathbf A}_p$ is defective, with
exactly one eigenvector missing for $\lambda=u$, unless $B_x=0$ and
$\varepsilon_BB_\perp\neq0$.
\end{enumerate}
\end{theorem}
 
\begin{proof}
The $B_x$ row and the $B_x$ column of \eqref{eq:prim_jac} vanish, so $\mathbf e_{B_x}$ is
an eigenvector with $\lambda=0$ and it suffices to consider the remaining
seven variables, ordered as $(\rho\,|\,\mathbf r)$ with
$\mathbf r=(u,v,w,p,B_y,B_z)^{\mathsf T}$. No equation other than the first
involves $\rho$, hence the corresponding $7\times7$ matrix is block triangular,
\begin{equation}\label{eq:A7}
\mathbf A_7=\begin{bmatrix} u & \rho\,\mathbf e_u^{\mathsf T}\\[2pt]
\mathbf 0 & \mathbf C\end{bmatrix},
\qquad
\mathbf C=\begin{bmatrix}
u & 0 & 0 & \varepsilon_M/\rho & \varepsilon_B B_y/\rho & \varepsilon_B B_z/\rho\\
0 & u & 0 & 0 & -\varepsilon_B B_x/\rho & 0\\
0 & 0 & u & 0 & 0 & -\varepsilon_B B_x/\rho\\
\varepsilon_M\rho a^2 & 0 & 0 & \varepsilon_M u & 0 & 0\\
\varepsilon_B B_y & -\varepsilon_B B_x & 0 & 0 & \varepsilon_B u & 0\\
\varepsilon_B B_z & 0 & -\varepsilon_B B_x & 0 & 0 & \varepsilon_B u
\end{bmatrix}.
\end{equation}
 
\emph{Step 1: symmetrisation.}
\emph{Step 1: symmetrisation.}
With the positive definite diagonal matrix
$\mathbf S=\operatorname{diag}\bigl(\rho,\rho,\rho,(\rho a^2)^{-1},1,1\bigr)$
one verifies directly that $\mathbf S\mathbf C$ is symmetric: the pairs of
off-diagonal entries become $(\varepsilon_M,\varepsilon_M)$,
$(\varepsilon_BB_{y,z},\varepsilon_BB_{y,z})$ and
$(-\varepsilon_BB_x,-\varepsilon_BB_x)$. Hence
$\mathbf S^{1/2}\mathbf C\,\mathbf S^{-1/2}
=\mathbf S^{-1/2}(\mathbf S\mathbf C)\,\mathbf S^{-1/2}$ is symmetric, so that
$\mathbf C$ has real eigenvalues and is diagonalisable over $\mathbb R$, for
all $\varepsilon_M,\varepsilon_B\in[0,1]$. The spectrum of $\mathbf A_7$ is
$\{u\}\cup\sigma(\mathbf C)$, which proves (i).
 
\emph{Step 2: a criterion for completeness.}
For $\mu\in\sigma(\mathbf C)$, $\mu\neq u$, every eigenvector $\mathbf r$ of
$\mathbf C$ lifts to the eigenvector
$\bigl(\rho\,r_u/(\mu-u),\,\mathbf r\bigr)$ of $\mathbf A_7$. For $\lambda=u$,
let $\mathcal N_u=\ker(\mathbf C-u\mathbf I)$ and $m=\dim\mathcal N_u\ge0$.
Since $\mathbf C$ is diagonalisable, the algebraic multiplicity of $u$ in
$\mathbf A_7$ is $1+m$, whereas
$(x,\mathbf r)\in\ker(\mathbf A_7-u\mathbf I)$ if and only if
$\mathbf r\in\mathcal N_u$ and $r_u=0$. Therefore $\mathbf A_7$ is
diagonalisable if and only if
\begin{equation}\label{eq:criterion}
r_u=0\quad\text{for every }\mathbf r\in\ker(\mathbf C-u\mathbf I),
\end{equation}
and otherwise exactly one eigenvector is missing, since $\{r_u=0\}$ then has
codimension one in $\mathcal N_u$.
 
\emph{Step 3: evaluation of the criterion.}
Writing $\mathbf r=(\delta u,\delta v,\delta w,\delta p,\delta B_y,\delta B_z)$,
the equations $(\mathbf C-u\mathbf I)\,\mathbf r=\mathbf 0$ read
\begin{align}
&\varepsilon_M\,\delta p+\varepsilon_B\,(B_y\,\delta B_y+B_z\,\delta B_z)=0,
\label{eq:ku}\\
&\varepsilon_BB_x\,\delta B_y=\varepsilon_BB_x\,\delta B_z=0,
\label{eq:kv}\\
&\varepsilon_M\rho a^2\,\delta u=(1-\varepsilon_M)\,u\,\delta p,
\label{eq:kp}\\
&\varepsilon_BB_y\,\delta u-\varepsilon_BB_x\,\delta v=(1-\varepsilon_B)\,u\,\delta B_y,
\qquad
\varepsilon_BB_z\,\delta u-\varepsilon_BB_x\,\delta w=(1-\varepsilon_B)\,u\,\delta B_z.
\label{eq:kB}
\end{align}
Suppose that \eqref{eq:criterion} fails and normalise $\delta u=1$.
 
\emph{Case $\varepsilon_M>0$.} By \eqref{eq:kp},
$(1-\varepsilon_M)\,u\,\delta p=\varepsilon_M\rho a^2>0$, so that
$\varepsilon_M<1$, $u\neq0$ and $\delta p\neq0$. If
$\varepsilon_BB_x\neq0$, \eqref{eq:kv} gives $\delta B_y=\delta B_z=0$ and
then \eqref{eq:ku} gives $\delta p=0$, a contradiction. If
$\varepsilon_BB_x=0$, \eqref{eq:kB} reduces to
$\varepsilon_BB_{y,z}=(1-\varepsilon_B)\,u\,\delta B_{y,z}$; multiplying
\eqref{eq:ku} by $(1-\varepsilon_M)(1-\varepsilon_B)\,u$ and inserting
\eqref{eq:kp} and \eqref{eq:kB} yields
\[
\varepsilon_M^2\,\rho a^2\,(1-\varepsilon_B)
+(1-\varepsilon_M)\,\varepsilon_B^2B_\perp^2=0 .
\]
Both terms are non-negative, so $\varepsilon_B=1$ and $B_\perp=0$; but then
\eqref{eq:ku} reads $\varepsilon_M\,\delta p=0$, again a contradiction. Hence
\eqref{eq:criterion} holds, which proves (ii).
 
\emph{Case $\varepsilon_M=0$.} If $\varepsilon_BB_x\neq0$, the vector
$\mathbf r=(1,\,B_y/B_x,\,B_z/B_x,\,0,\,0,\,0)$ satisfies
\eqref{eq:ku}--\eqref{eq:kB}; if $\varepsilon_B=0$, or if $B_x=0$ and
$B_\perp=0$, so does $\mathbf r=\mathbf e_u$. In each of these cases
\eqref{eq:criterion} is violated. In the remaining case, $B_x=0$ and
$\varepsilon_BB_\perp\neq0$, \eqref{eq:kB} gives
$\varepsilon_BB_{y,z}=(1-\varepsilon_B)\,u\,\delta B_{y,z}$, and multiplying
\eqref{eq:ku} by $(1-\varepsilon_B)\,u$ yields
$\varepsilon_B^2B_\perp^2=0$, a contradiction, so that \eqref{eq:criterion}
holds. This proves (iii).
\end{proof}
 
\begin{remark}[Non-uniform conditioning]\label{rem:nonuniform}
The eigenvector basis of Theorem~\ref{thm:hyp}(ii) is not uniformly
conditioned as $\varepsilon_M\to0$. For $\varepsilon_BB_x\neq0$ the entropy
strength of a jump $\Delta\mathbf Q$ in the split basis is
\begin{equation}\label{eq:went_split}
w_{\rm ent}^{\rm split}
=\Delta\rho-\frac{\Delta p}{\varepsilon_Ma^2}
-\kappa_M\Bigl(\Delta u+\frac{B_y\,\Delta v+B_z\,\Delta w}{B_x}\Bigr),
\qquad
\kappa_M=\frac{\rho\,(1-\varepsilon_M)\,u}{\varepsilon_M^2\,a^2},
\end{equation}
which reduces to the classical $\Delta\rho-\Delta p/a^2$ for
$\varepsilon_M=1$ and grows as $\varepsilon_M^{-2}$ otherwise, the trace of
the Jordan block of case (iii). In the characteristic reconstruction of Section~\ref{sec:reconstruction} the basis is therefore evaluated with $\varepsilon_M$ bounded from below by $10^{-2}$; by Theorem~\ref{thm:hyp} no bound on $\varepsilon_B$ is required.

\end{remark}

\paragraph{Right eigenvectors.}
The split right eigenvectors are closed-form deformations of the standard
Roe--Balsara basis~\cite{RoeBalsara1996}; we give them in the primitive
variables $\mathbf Q$, the conservative eigenvectors following by the
transformation $\mathbf M=\partial\mathbf U/\partial\mathbf Q$. The divergence
mode is $\mathbf e_{B_x}$. The entropy eigenvector is $\varepsilon$-independent
and, consistently with $\delta p=\delta\mathbf u=\delta\mathbf B=0$, is carried
entirely by the density, $\mathbf r_{\rm ent}=\mathbf e_\rho$. With
$\alpha_A^{\pm}=u-\tilde\lambda_A^{\pm}$, the Alfv\'en eigenvectors follow from
the perpendicular $2\times2$ block,
\begin{equation}\label{eq:R_alf}
\mathbf r_A^{\pm}=\Big(0,0,-\tfrac{\varepsilon_B B_xB_z}{\rho\,\alpha_A^{\pm}},
\tfrac{\varepsilon_B B_xB_y}{\rho\,\alpha_A^{\pm}},0,0,-B_z,B_y\Big)^{\mathsf T},
\end{equation}
and for each root $\lambda$ of the magnetosonic quartic~\eqref{eq:ms_disp},
\begin{equation}\label{eq:r_ms}
\mathbf r_{\rm ms}=\Bigl(-\frac{\rho}{\alpha},\;1,\;
-\frac{\varepsilon_B^2B_xB_y}{\rho D_A},\;
-\frac{\varepsilon_B^2B_xB_z}{\rho D_A},\;
-\frac{\varepsilon_M\rho a^2}{\beta},\;0,\;
-\frac{\varepsilon_BB_y\,\alpha}{D_A},\;
-\frac{\varepsilon_BB_z\,\alpha}{D_A}\Bigr)^{\mathsf T},
\end{equation}
with $\alpha$, $\beta$, $D_A$ from~\eqref{eq:DsDA} evaluated at that
$\lambda$. Theorem~\ref{thm:hyp}(ii) guarantees a complete set of
eigenvectors for $\varepsilon_M>0$;    the closed forms above, however, degenerate where $B_\perp$ vanishes and become singular where $\alpha$, $\beta$ or $D_A$ vanish, and in the
implementation the Alfv\'en pair is taken in the regularised form of
Corollary~\ref{cor:alfven} and the aligned sector is diagonalised
numerically, see Lemma~\ref{lem:blocks}.

\paragraph{Implicit sub-systems.}
Since $\tilde{\mathbf F}^p=(1-\epsM)\,\mathbf F^p$ and
$\tilde{\mathbf F}^B=(1-\epsB)\,\mathbf F^B$, the implicit sub-systems of the
$\varepsilon$-split form have, for frozen parameters, the eigenvalues
\eqref{eq:p_eigs} and~\eqref{eq:B_eigs} scaled by the complementary weights,
\begin{equation}
\label{eq:implicit_eigs}
\tilde\lambda^{p}_{7,8}=\frac{1-\epsM}{2}\!\left(u\pm\sqrt{u^2+4a^2}\right),
\qquad
\tilde\lambda^{B}_{5,6}=\frac{1-\epsB}{2}\!\left(u\pm\sqrt{u^2+4c_a^2}\right),
\qquad
\tilde\lambda^{B}_{7,8}=\frac{1-\epsB}{2}\!\left(u\pm\sqrt{u^2+4c_B^2}\right),
\end{equation}
the remaining ones being zero. Their implicit treatment thus removes the
acoustic and the Alfv\'en stiffness independently of one another, while the
explicit step is governed by $\tilde\lambda^{c}$ alone.

The closed-form eigenvectors fail in two limits: the Alfv\'en
pair~\eqref{eq:R_alf} degenerates as $B_\perp\to0$, and the entropy
projection is singular as $\epsM\to0$ (Remark~\ref{rem:nonuniform}). The
following decomposition confines each of them to a sub-block.

\begin{lemma}[Block decomposition]
\label{lem:blocks}
With $B_\perp$, $\beta_y$, $\beta_z$ as in~\eqref{eq:beta} (any fixed unit vector
when $B_\perp=0$), set
\[
\mathbf e_v=(0,0,\beta_z,-\beta_y,0,0,0,0)^{\mathsf T},\qquad
\mathbf e_B=(0,0,0,0,0,0,\beta_z,-\beta_y)^{\mathsf T}.
\]
Then $\AAt_{p}$ of~\eqref{eq:prim_jac} is block diagonal with respect to
\[
\mathbb R^{8}
=\operatorname{span}\{\mathbf e_{B_x}\}
\oplus\underbrace{\operatorname{span}\{\mathbf e_v,\mathbf e_B\}}_{S_A}
\oplus\underbrace{\operatorname{span}\{\mathbf e_\rho,\mathbf e_u,
 \mathbf e_{v_\parallel},\mathbf e_p,\mathbf e_{B_\parallel}\}}_{S_\parallel},
\]
and specifically
\begin{enumerate}
\item[\rm(i)] the $B_x$ row and the $B_x$ column of $\AAt_p$ vanish, so
$\Delta B_x$ is itself a characteristic variable with $\tilde\lambda=0$;
\item[\rm(ii)] $S_\parallel$ is invariant, the induced action being
\begin{equation}
\label{eq:A5}
\mathbf A_5=
\begin{bmatrix}
u & \rho & 0 & 0 & 0\\[2pt]
0 & u & 0 & \tfrac{\epsM}{\rho} & \tfrac{\epsB B_\perp}{\rho}\\[2pt]
0 & 0 & u & 0 & -\tfrac{\epsB B_x}{\rho}\\[2pt]
0 & \epsM\rho a^2 & 0 & \epsM u & 0\\[2pt]
0 & \epsB B_\perp & -\epsB B_x & 0 & \epsB u
\end{bmatrix},
\end{equation}
whose spectrum is the entropy speed $u$ together with the four roots
of~\eqref{eq:ms_disp};
\item[\rm(iii)] $S_A$ is invariant with induced action
\begin{equation}
\label{eq:A2}
\mathbf A_2=\begin{bmatrix} u & -\sigma_v\\ -\sigma_B & \epsB u\end{bmatrix},
\qquad \sigma_v=\frac{\epsB B_x}{\rho},\qquad \sigma_B=\epsB B_x,
\end{equation}
whose eigenvalues are the split Alfv\'en speeds~\eqref{eq:alfven_eig}, and the
functionals
$\ell_v(\delta\mathbf Q)=\beta_z\,\delta v-\beta_y\,\delta w$,
$\ell_B(\delta\mathbf Q)=\beta_z\,\delta B_y-\beta_y\,\delta B_z$
vanish on every eigenvector
of $\AAt_p$ not contained in $S_A$.
\end{enumerate}
\end{lemma}

\begin{proof}
(i) is read off~\eqref{eq:prim_jac}. For (iii), let $\delta\mathbf Q=c_1\mathbf e_v+c_2\mathbf e_B$. The $u$-row contributes $\epsB\rho^{-1}c_2(B_y\beta_z-B_z\beta_y)=0$; the $\rho$- and $p$-rows vanish because $\delta\mathbf Q$ has no $\rho$, $u$ or $p$ component; the $v$- and $w$-rows give
$(uc_1-\sigma_vc_2)(\beta_z,-\beta_y)$ and the $B_y$-, $B_z$-rows give
$(-\sigma_Bc_1+\epsB u\,c_2)(\beta_z,-\beta_y)$, both again in $S_A$. The
annihilation statement holds because the magnetosonic eigenvectors have
$(\delta v,\delta w)$ and $(\delta B_y,\delta B_z)$ parallel to
$(\beta_y,\beta_z)$, the entropy eigenvector is proportional to $\mathbf e_\rho$,
and the divergence eigenvector to $\mathbf e_{B_x}$. Statement (ii) is the same
computation for the complementary rotation.
\end{proof}

Since $\mathbf A_5$ and $\mathbf A_2$ are restrictions of $\tilde{\mathbf A}_p$ to
the invariant subspaces $S_\parallel$ and $S_A$, Theorem~\ref{thm:hyp}(i)
implies that the Alfv\'en quadratic~\eqref{eq:alfven_disp} and the
magnetosonic quartic~\eqref{eq:ms_disp} have only real roots, for every
value of $u$.

\begin{corollary}[Regularised split Alfv\'en eigenvectors]
\label{cor:alfven}
The eigenpairs of~\eqref{eq:A2} are
\[
\mathbf r_\pm=(\sigma_v,\ \alpha_A^{\pm}),\qquad
\boldsymbol\ell_\pm=(\sigma_B,\ \alpha_A^{\pm}),\qquad
\boldsymbol\ell_\pm\!\cdot\mathbf r_\pm=\epsB^{2}c_a^{2}+(\alpha_A^{\pm})^{2},
\]
so that
\begin{equation}
\label{eq:Ralf_reg}
\hat{\mathbf r}_A^{\pm}=\sigma_v\,\mathbf e_v+\alpha_A^{\pm}\,\mathbf e_B
\end{equation}
is~\eqref{eq:R_alf} rescaled by $-\alpha_A^{\pm}/B_\perp$ and is regular as
$B_\perp\to0$.
\end{corollary}

The normalisation $\epsB^{2}c_a^{2}+(\alpha_A^{\pm})^{2}$ vanishes only as
$\epsB B_x\to0$, and then only for the branch with $\alpha_A^{\pm}=0$. This is a
degeneracy of the parametrisation, not of the system: at $\epsB B_x=0$ the
block~\eqref{eq:A2} is diagonal, so $v_\perp$ and $B^{*}_\perp$ decouple with the
speeds $u$ and $\epsB u$ and the eigenvectors $\mathbf e_v$, $\mathbf e_B$. This
limit is used when $\epsB|B_x|/\sqrt\rho$ falls below $10^{-8}$ times the local
signal speed.
 

 
% ===========================================================================
\section{Semi-implicit numerical scheme}
\label{sec:scheme}

This section presents the semi-implicit numerical method proposed in this work. We
begin by establishing the computational framework and notation, followed by the
adaptive splitting strategy, the explicit treatment of the advective sub-system,
and the implicit solution of the magnetic and pressure sub-systems.

We restrict our presentation to two spatial dimensions on a uniform Cartesian mesh. The computational domain
$\Omega=[x_{\min},x_{\max}]\times[y_{\min},y_{\max}]$ is partitioned into
$N_x\times N_y$ rectangular control volumes with uniform spacing
\begin{equation}
\label{eq:mesh}
\Delta x=\frac{x_{\max}-x_{\min}}{N_x},\qquad
\Delta y=\frac{y_{\max}-y_{\min}}{N_y}.
\end{equation}
A control volume is indexed by the pair $(i,j)$, with its centre located at
$\mathbf x_{i,j}=(x_i,y_j)$. The midpoints of the cell interfaces in the $x$- and
$y$-directions are
\begin{equation}
\label{eq:face_centres}
\mathbf x_{i+\frac12,j}=\left(\frac{x_i+x_{i+1}}{2},\;y_j\right),\qquad
\mathbf x_{i,j+\frac12}=\left(x_i,\;\frac{y_j+y_{j+1}}{2}\right).
\end{equation}
Throughout this paper, superscripts denote discrete time levels and subscripts
identify spatial locations. The scheme is a finite volume method: the discrete
unknowns are the cell averages of the conservative variables, all of which,
including the magnetic field, are associated with the control volumes
(collocated arrangement). The divergence-free constraint is enforced by an
unstaggered constrained-transport step that evolves the cell averages of the
magnetic vector potential, see Section~\ref{sec:CT}.

The time interval $[0,t_{\mathrm{out}}]$ is discretised by time steps
$\Delta t=t^{n+1}-t^{n}$, always subject to a CFL-type stability condition based
\emph{only} on the characteristic speeds of the explicit sub-system
\eqref{eq:new_split}. Thus, with $|\tilde\lambda^{c}_{\max}|^{n}$ denoting the
largest eigenvalue of the split convective sub-system at time level $t^{n}$,
\begin{equation}
\label{eq:cfl}
\Delta t\le \mathrm{CFL}\,
\frac{\min(\Delta x,\Delta y)}{\max_\Omega|\tilde\lambda^{c}_{\max}|^{n}}.
\end{equation}
Condition~\eqref{eq:cfl} is based on the $\eps$-split convective spectrum rather
than on the full MHD spectrum~\eqref{eq:MHD_eigs}. In the low-Mach and low-Alfv\'en
regime $\epsM,\epsB\to0$ and $|\tilde\lambda^{c}_{\max}|\to\|\mathbf{u}\|$, so that the
time step is governed by the material velocity rather than by the sound speed or
the Alfv\'en speed; in the supersonic and super-Alfv\'enic regime
$\epsM,\epsB\to1$ and the full MHD spectrum is recovered. Because the two
parameters act independently, either constraint may be removed on its own: the
time step is freed from the acoustic restriction wherever $M\ll1$ and from the
Alfv\'enic restriction wherever $M_A\ll1$, regardless of the other.

Two remarks on what this does and does not gain. A single parameter driven by the
more restrictive of the two Mach numbers vanishes whenever either channel is
stiff, so it yields the same $\mathcal O(\|\mathbf u\|)$ explicit spectrum
as~\eqref{eq:cfl}; the second parameter then changes not the step but the
discretisation of the channel that is not stiff, which is returned to the upwind
solver of Section~\ref{sec:hllem} instead of being carried by the central-type
implicit operator. It is against a single parameter driven by the acoustic Mach
number alone that the step itself improves, and then by orders of magnitude,
since such a parameter leaves $\FF^{B}$ explicit in a strongly magnetised flow.

\subsection{Adaptive $\eps$-splitting strategy}
\label{sec:eps_strategy}

The two splitting parameters control the partition of the acoustic and of the
magnetic flux between the explicit and the implicit sub-systems, and their values
should reflect the local flow regime. We define them from the local acoustic and
Alfv\'en Mach numbers~\eqref{eq:two_mach} through a single switching
function. Following~\cite{NewAP}, for each cell $(i,j)$ we set
\begin{equation}
\label{eq:switch}
s(\mu)=\begin{cases}
1-\mu^{\alpha}, & 0\le\mu\le M_0,\\[4pt]
\exp\!\Big(1-\dfrac{1}{1-\big(\frac{\mu-M_0}{M_1-M_0}\big)^{2}}\Big)
\big[(1-M_0^{\alpha})-(1-M_1)^{\alpha}\big]+(1-M_1)^{\alpha}, & M_0<\mu<M_1,\\[6pt]
(1-\mu)^{\alpha}, & M_1\le\mu\le1,\\[4pt]
0,& \mu>1,
\end{cases}
\end{equation}
with $M_0=0.15$, $M_1=0.4$ and $\alpha=14$, and then define
\begin{equation}
\label{eq:eps_def}
\epsM{}_{,i,j}=1-s\big(M_{i,j}\big),
\qquad
\epsB{}_{,i,j}=1-s\big(M_{A,i,j}\big).
\end{equation}
The same switching function is applied to two different arguments: the
acoustic parameter is driven by $M$ and the magnetic one by $M_A$. By
construction $\epsM,\epsB\in[0,1]$, with $\varepsilon\to0$ in the
corresponding stiff limit (where nearly the entire flux of that channel is
delegated to the implicit solver) and $\varepsilon=1$ once the corresponding
Mach number reaches unity (where that channel is returned entirely to the
explicit sub-system), with a continuous monotone transition in between. By Theorem~\ref{thm:hyp} the explicit sub-system has real eigenvalues for all
values of the parameters and a complete set of eigenvectors whenever
$\epsM>0$, irrespective of $\epsB$. Since $\epsM=M^{\alpha}$ for $M\le M_0$, it
vanishes only where the velocity vanishes, but it falls below machine precision
already for $M\lesssim0.07$, so that low-Mach regions are numerically at the
limit $\epsM=0$ of Theorem~\ref{thm:hyp}(iii). The Riemann solver does not
require $\epsM>0$, since it uses only the regularised Alfv\'en projection and
the classical entropy strength (Section~\ref{sec:hllem}). The characteristic
basis does: its conditioning degrades as $\epsM\to0$
(Remark~\ref{rem:nonuniform}), and it is therefore evaluated with $\epsM$
bounded from below by $10^{-2}$, as described in
Section~\ref{sec:reconstruction}. No bound on $\epsB$ is needed.

Since the numerical flux is evaluated at cell interfaces, the interface splitting
parameters are defined as
\begin{equation}
\label{eq:eps_face}
\eps^{x}_{i+\frac12,j}=\max\!\left(\eps_{i,j},\;\eps_{i+1,j}\right),
\qquad
\eps^{y}_{i,j+\frac12}=\max\!\left(\eps_{i,j},\;\eps_{i,j+1}\right),
\end{equation}
for $\eps\in\{\epsM,\epsB\}$ alike. The $\max$ operator ensures that if either
adjacent cell requires a larger explicit contribution, the interface flux is
computed with the more explicit splitting, which enhances stability at transitions
between stiff and non-stiff regions.

The behaviour of each channel separately is summarised as follows.
\begin{itemize}
\item \textbf{Acoustically supersonic} ($M\ge1$): $\epsM=1$. The pressure flux is
entirely retained in the explicit sub-system, and the acoustic waves are handled
by the Riemann solver.
\item \textbf{Acoustically stiff} ($M\ll1$): $\epsM=M^{\alpha}\approx0$. Nearly
all of the acoustic flux is delegated to the implicit solver, and the explicit
time step ceases to feel the sound speed $a$.
\item \textbf{Magnetically super-Alfv\'enic} ($M_A\ge1$): $\epsB=1$, the magnetic
flux is fully explicit, and the Alfv\'en and magnetosonic waves are resolved by
the Riemann solver.
\item \textbf{Magnetically stiff} ($M_A\ll1$): $\epsB\approx0$, the magnetic
flux is delegated to the implicit solver, and the explicit time step ceases to feel the Alfv\'en speed $c_B$.
\end{itemize}

Because $\epsM$ and $\epsB$ are evaluated independently, the four combinations of
these regimes are all admissible, and it is precisely the two \emph{mixed} cases
that a single-parameter splitting cannot represent: a flow that is acoustically
stiff but magnetically non-stiff (a high-$\beta$ low-Mach plasma) would have its
Alfv\'en waves over-stabilised, while a flow that is magnetically stiff but
acoustically non-stiff (a strongly magnetised low-$\beta$ plasma at moderate Mach
number) would have its acoustic waves needlessly implicit.


\subsection{Advection sub-system}
\label{sec:hllem}

The advective sub-system is integrated explicitly using a standard finite volume
formulation. At each time step, the cell-averaged state $\mathbf{U}^{n}_{i,j}$ is advanced
to an intermediate value $\mathbf{U}^{*}_{i,j}$ via
\begin{equation}
\label{eq:adv_update}
\mathbf{U}^{*}_{i,j}=\mathbf{U}^{n}_{i,j}
-\frac{\Delta t}{\Delta x}\big(\tilde\FF^{\eps}_{i+\frac12,j}-\tilde\FF^{\eps}_{i-\frac12,j}\big)
-\frac{\Delta t}{\Delta y}\big(\tilde\GG^{\eps}_{i,j+\frac12}-\tilde\GG^{\eps}_{i,j-\frac12}\big),
\end{equation}
where $\tilde\FF^{\eps}$ and $\tilde\GG^{\eps}$ are numerical fluxes consistent
with the split convective flux $\tilde\FF^{c}$ of~\eqref{eq:new_split}. We
describe the construction in the $x$-direction; the $y$-direction flux follows by
the cyclic permutation $(x,y,z)\to(y,z,x)$. We note that $\rho^{n+1}=\rho^{*}$,
since the density flux involves neither the pressure nor the magnetic
contribution and is therefore already fully updated at this stage.

\paragraph{Standard IMEX baseline: scalar-dissipation fluxes.}
Existing semi-implicit three-way splittings~\cite{Fambri2021,MHDIMEX}
correspond to the fixed choice $\epsM=\epsB\equiv0$, for which the convective
sub-system is only weakly hyperbolic. With the whole magnetic flux on the
implicit side the Alfv\'en fields leave no trace in the explicit spectrum at
all, so there is nothing for an intermediate-wave correction to act on, and a
scalar dissipation of Rusanov type~\cite{Rusanov1962} is forced,
\begin{equation}
\label{eq:rusanov}
\tilde\FF^{\eps}_{i+\frac12,j}
=\half\!\left[\tilde\FF^{c}(\mathbf{U}^{L})+\tilde\FF^{c}(\mathbf{U}^{R})
-s^{\max}_{i+\frac12,j}\big(\mathbf{U}^{R}-\mathbf{U}^{L}\big)\right],
\qquad
s^{\max}_{i+\frac12,j}=\max_{\mathbf{U}\in\{\mathbf{U}^{L},\mathbf{U}^{R}\}}\big|\tilde\lambda^{c}_{\max}(\mathbf{U})\big| ,
\end{equation}
which smears the entropy, Alfv\'en and magnetosonic fields equally.

\paragraph{$\eps$-IMEX: AP-HLLEM flux.}
The HLLEM scheme of Einfeldt et al.~\cite{einfeldt1991godunov}, in the
formulation of Dumbser and Balsara~\cite{HLLEM}, needs only the linearly
degenerate fields and their projections, not a complete eigenvector basis, and
is therefore available at every value of the parameters---including the limit
$\epsM=0$ in which the convective sub-system is defective by
Theorem~\ref{thm:hyp}(iii). The block decomposition of
Lemma~\ref{lem:blocks} supplies the Alfv\'en pair in the regularised form of
Corollary~\ref{cor:alfven}, together with the transverse
functionals~\eqref{eq:alf_jumps} that isolate its content; the entropy field is
taken with the classical strength~\eqref{eq:went}, bounded uniformly in
$\epsM$. No bound on $\epsM$ is imposed in the flux: the bound of
Remark~\ref{rem:nonuniform} is confined to the characteristic reconstruction of
Section~\ref{sec:reconstruction}, which is what the complete basis of
Theorem~\ref{thm:hyp}(ii) licenses.
Starting from the two-wave HLLE solver, HLLEM adds an anti-diffusive correction
along the linearly degenerate intermediate fields, so that contact and rotational
discontinuities are resolved sharply while the robustness of HLLE is preserved.
The construction below is built so that the entire asymptotic-preserving scaling
resides in the wave coefficients $\delta_k$, while the wave structure is supplied
by the block decomposition of Section~\ref{sec:eigenstructure}.

The averaged state is the Cargo--Gallice MHD Roe mean~\cite{CargoGallice1997},
\begin{equation}
\label{eq:roe}
\tilde\rho=\sqrt{\rho_L\rho_R},\quad
\tilde\psi=\frac{\sqrt{\rho_L}\psi_L+\sqrt{\rho_R}\psi_R}{\sqrt{\rho_L}+\sqrt{\rho_R}}\ \ (\psi=u,v,w,H),\quad
\tilde B_i=\frac{\sqrt{\rho_R}B_{iL}+\sqrt{\rho_L}B_{iR}}{\sqrt{\rho_L}+\sqrt{\rho_R}}.
\end{equation}
where $H=(E+p+\tfrac12\|\mathbf B\|^2)/\rho$ is the total specific enthalpy, and the
Roe sound speed is
$\tilde a^2=(\gamma-1)\big(\tilde H-\tfrac12\|\tilde{\mathbf u}\|^2-\|\tilde{\mathbf B}\|^2/\tilde\rho\big)$. The wave-speed envelope is given by the Einfeldt estimates built from the extreme
\emph{split} eigenvalues of the left, right and Roe-average states. Denoting by
$\tilde\lambda^{\mp}_{f}(\mathbf{U})$ the smallest and largest roots of the magnetosonic
quartic~\eqref{eq:ms_disp} at a state $\mathbf{U}$,
\begin{equation}
\label{eq:einfeldt}
s_{L}=\min\Bigl(\min_{\mathbf{U}\in\{\mathbf{U}^{L},\tilde{\mathbf{U}},\mathbf{U}^{R}\}}\tilde\lambda^{-}_{f}(\mathbf{U}),\;-u^{\max}_{i+\frac12}\Bigr),
\qquad
s_{R}=\max\Bigl(\max_{\mathbf{U}\in\{\mathbf{U}^{L},\tilde{\mathbf{U}},\mathbf{U}^{R}\}}\tilde\lambda^{+}_{f}(\mathbf{U}),\;+u^{\max}_{i+\frac12}\Bigr),
\end{equation}
where
\begin{equation}
\label{eq:umaxface}
u^{\max}_{i+\frac12}
=\max_{\mathbf{U}\in\{\mathbf{U}^{L},\tilde{\mathbf{U}},\mathbf{U}^{R}\}}\bigl|u(\mathbf{U})\bigr|
\end{equation}
is the largest material speed among the same three states. Since $\tilde u$ is a
convex combination of $u^{L}$ and $u^{R}$, \eqref{eq:umaxface} equals
$\max(|u^{L}|,|u^{R}|)$, which is precisely the Rusanov speed used for the
convective sub-system in~\cite{MHDIMEX}: the bound introduces no speed that the
baseline scheme does not already use.

The bound closes the envelope in the stiff limit. As $\epsM,\epsB\to0$ the
magnetosonic quartic~\eqref{eq:ms_disp} tends to $\lambda^{2}(\lambda-u)^{2}$,
so its roots tend to $\{0,0,u,u\}$ and the inner extrema
of~\eqref{eq:einfeldt} tend to $\min(0,u)$ and $\max(0,u)$: wherever $u$ keeps
one sign across the interface one endpoint is exactly zero and
$s_{L}s_{R}=0$. Since~\eqref{eq:umaxface} is built from the states of that
interface alone and is itself a material speed, the bound leaves the asymptotic
scaling of the dissipation unchanged and contains no free parameter.
Only the magnetosonic roots enter~\eqref{eq:einfeldt}, and they are enough:
on the invariant blocks of Lemma~\ref{lem:blocks} the interval
$[\tilde\lambda^{-}_{f},\tilde\lambda^{+}_{f}]$ also contains the split
Alfv\'en speeds~\eqref{eq:alfven_eig}. We have verified this over the whole
admissible range of states and of $(\epsM,\epsB)\in[0,1]^2$; the two intervals
touch only in the degenerate limit $\epsB B_x\to0$, where the Alfv\'en pair
collapses onto $\{u,\epsB u\}$. The anti-diffusion coefficients
$\delta_k$ of~\eqref{eq:deltak} therefore lie in $[0,1]$ for all three linearly
degenerate fields without being clipped there; the shock sensor
of~\eqref{eq:flattener} is the only further modification they receive.

Writing the quartic~\eqref{eq:ms_disp} in monic form,
$\lambda^4+c_3\lambda^3+c_2\lambda^2+c_1\lambda+c_0=0$, all its roots are
real by Lemma~\ref{lem:blocks} and bounded in modulus by
\begin{equation}\label{eq:fujiwara}
R=2\max\Bigl(|c_3|,\;|c_2|^{1/2},\;|c_1|^{1/3},\;\bigl|\tfrac12c_0\bigr|^{1/4}\Bigr).
\end{equation}
Newton's method started from $R$ therefore decreases monotonically to
$\tilde\lambda_f^{+}$, and started from $-R$ increases monotonically to
$\tilde\lambda_f^{-}$, so no eigenvalue solver is required at the interface. The
iteration is stopped at a relative step of $10^{-13}$, after $200$ iterations,
or as soon as the derivative or the step changes sign, which cannot occur while
the iterate lies outside the root interval and therefore serves as a
monotonicity guard. A few tens of scalar iterations suffice; in the stiff limit
the two targeted roots are double, the convergence is linear, and the attainable
accuracy is that of a double root, which is immaterial there because the bound
in~\eqref{eq:einfeldt} replaces both speeds by $\mp u^{\max}_{i+\frac12}$. The
monotone iteration is also what makes the evaluation robust. The roots are real, but two of them nearly coincide over
a non-negligible part of the state space, and a general-purpose root finder
applied to the coefficients $(c_3,c_2,c_1,c_0)$ then returns spurious complex
pairs -- for randomly sampled admissible states, in a few per cent of them.
Newton started outside the root interval of a polynomial with real roots cannot
leave it, so the envelope is obtained without any such failure.
For $s_L\ge0$ or $s_R\le0$ the
flux reduces to the upwind value $\tilde\FF^{c}(\mathbf{U}^{L})$ or
$\tilde\FF^{c}(\mathbf{U}^{R})$; for $s_L<0<s_R$ the AP-HLLEM flux reads
\begin{equation}
\label{eq:hllem}
\tilde\FF^{\eps}_{i+\frac12,j}
=\FF^{\mathrm{HLL}}_{i+\frac12,j}
-\frac{s_Ls_R}{s_R-s_L}\sum_{k\in\{\mathrm{ent},A^{\pm}\}}\delta_k\,w_k\,\RR_k,
\end{equation}
with the standard HLL flux
\begin{equation}
\label{eq:hll}
\FF^{\mathrm{HLL}}_{i+\frac12,j}
=\frac{s_R\,\tilde\FF^{c}(\mathbf{U}^{L})-s_L\,\tilde\FF^{c}(\mathbf{U}^{R})
+s_Ls_R\big(\mathbf{U}^{R}-\mathbf{U}^{L}\big)}{s_R-s_L}.
\end{equation}
If $p^{L}\le0$, $p^{R}\le0$ or $\tilde a^{2}\le10^{-10}$ the anti-diffusive term
is omitted.

\paragraph{Wave strengths.}
The three linearly degenerate fields are the entropy wave at
$\tilde\lambda_{\mathrm{ent}}=\tilde u$ and the two Alfv\'en waves at
$\tilde\lambda_A^{\pm}$ of~\eqref{eq:alfven_eig}. Throughout,
$\Delta(\cdot)=(\cdot)^{R}-(\cdot)^{L}$ denotes the interface jump and tildes
denote Roe-average quantities.

For the Alfv\'en pair the projection is available in closed form: by
Lemma~\ref{lem:blocks}(iii) the transverse functionals
\begin{equation}
\label{eq:alf_jumps}
\Delta v_\perp=\beta_z\,\Delta v-\beta_y\,\Delta w,
\qquad
\Delta B_\perp^{*}=\beta_z\,\Delta B_y-\beta_y\,\Delta B_z
\end{equation}
isolate the Alfv\'en content of the interface jump, and
Corollary~\ref{cor:alfven} gives the strengths and the regularised
eigenvectors
\begin{equation}
\label{eq:walf}
w_A^{\pm}=\frac{\sigma_B\,\Delta v_\perp+\alpha_A^{\pm}\,\Delta B_\perp^{*}}
{\epsB^{2}c_a^{2}+(\alpha_A^{\pm})^{2}},
\qquad
\hat{\mathbf r}_A^{\pm}=\sigma_v\,\mathbf e_v+\alpha_A^{\pm}\,\mathbf e_B,
\end{equation}
with $\sigma_v$, $\sigma_B$ from~\eqref{eq:A2}. For
$\epsB|\tilde B_x|/\sqrt{\tilde\rho}\to0$ the pair is replaced by its
decoupled limit.

The entropy field is treated differently. Its exact strength in the split
basis, \eqref{eq:went_split}, is unbounded as $\epsM\to0$. In the complete
eigen-decomposition $\sum_kw_k\RR_k=\Delta\mathbf U$ this growth is cancelled
by the density components of the magnetosonic fields, but the anti-diffusive
term of~\eqref{eq:hllem} weights the fields with different $\delta_k$, so
that the cancellation is lost. The entropy strength is therefore taken in the
classical form, which is bounded for all $\epsM$ and coincides
with~\eqref{eq:went_split} across a contact discontinuity,
\begin{equation}
\label{eq:went}
w_{\mathrm{ent}}=\Delta\rho-\frac{\Delta p}{\tilde a^{2}},
\qquad
\RR_{\mathrm{ent}}=\bigl(1,\,\tilde u,\,\tilde v,\,\tilde w,\,
\half\|\tilde{\mathbf u}\|^{2},\,0,\,0,\,0\bigr)^{\mathsf T}.
\end{equation}

\paragraph{AP scaling in $\delta_k$.}
The anti-diffusion coefficients measure the separation of each linearly
degenerate speed from the wave-speed envelope, and are built from the
\emph{split} eigenvalues,
\begin{equation}
\label{eq:deltak}
\delta_k=1-\frac{\tilde\lambda_k^{-}}{s_L}-\frac{\tilde\lambda_k^{+}}{s_R},
\qquad
\tilde\lambda^{\pm}=\tfrac12\bigl(\tilde\lambda\pm|\tilde\lambda|\bigr),
\end{equation}
restricted to $[0,1]$. By construction $\delta_k\to0$ when the corresponding
speed approaches the envelope, as it does near strong shocks, and
$\delta_k\to1$ for a well-separated field, restoring the resolution of a
complete Riemann solver. In the stiff limit $\epsM,\epsB\to0$ the split spectrum collapses onto $\{0,u\}$,
so that $s_L,s_R=\mathcal O(|u|)$: the coefficient $s_Ls_R/(s_R-s_L)$
of~\eqref{eq:hllem}, and with it both the HLL dissipation and the anti-diffusive
term, scales with the material velocity rather than with the sound or Alfv\'en
speed. The bound of~\eqref{eq:einfeldt} preserves that scaling, being itself a
material speed of the interface, and it fixes the limit uniquely: with
$s_{L}=-u^{\max}_{i+\frac12}$ and $s_{R}=u^{\max}_{i+\frac12}$ the HLL
part~\eqref{eq:hll} becomes
\begin{equation}
\label{eq:aplimit}
\FF^{\mathrm{HLL}}_{i+\frac12,j}
=\half\!\left[\tilde\FF^{c}(\mathbf{U}^{L})+\tilde\FF^{c}(\mathbf{U}^{R})
-u^{\max}_{i+\frac12}\big(\mathbf{U}^{R}-\mathbf{U}^{L}\big)\right],
\end{equation}
which is exactly the scalar-dissipation flux~\eqref{eq:rusanov} that weak
hyperbolicity forces on the standard IMEX baseline, with the same speed: in this
limit $|\tilde\lambda^{c}_{\max}|\to|u|$, so $s^{\max}$ of~\eqref{eq:rusanov}
and $u^{\max}_{i+\frac12}$ of~\eqref{eq:umaxface} coincide. The AP-HLLEM flux is
therefore that baseline minus the anti-diffusive sum of~\eqref{eq:hllem}, whose
entropy coefficient there is
$\delta_{\mathrm{ent}}=1-|u|/u^{\max}_{i+\frac12}$. Since the bound is taken
over the states of that interface alone, this coefficient vanishes on a uniform
interface and is positive only across a genuine velocity jump: the
anti-diffusion switches itself off in the quiet part of the flow and acts where
the linearly degenerate fields carry structure. A bound taken over the whole
domain instead would drive $\delta_{\mathrm{ent}}\to1$ wherever
$|u|\ll\max_\Omega|u|$ and so remove the dissipation from the three linearly
degenerate fields in the quiet ambient of a flow with a large dynamic range. The
asymptotic-preserving property is thus continuous in $\epsM,\epsB$ and needs no
solver switching: the $\eps$-split does not discard the Rusanov
flux~\eqref{eq:rusanov}, it recovers it as the stiff limit and sharpens the
three linearly degenerate fields on top of it.

\paragraph{Shock stabilisation of the anti-diffusion.}
On an interface whose normal is transverse to a strong shock the normal velocity
is small, so $\tilde\lambda_{\mathrm{ent}}=\tilde u$ lies far from both ends of
$[s_L,s_R]$ and~\eqref{eq:deltak} gives $\delta_{\mathrm{ent}}\to1$, leaving the
entropy anti-diffusion along the front undamped; the two Alfv\'en coefficients
behave likewise. This is the grid-aligned shock instability of HLLEM, which the
linear perturbation analysis of~\cite{SimonMandal2019} attributes to the
activation of the anti-diffusive terms on the interfaces transverse to the
shock. All three coefficients are therefore damped by a dimensionless sensor
built from the pressure, following the pressure criterion that accompanies the
divergence flattener of Balsara~\cite{Balsara2012},
\begin{equation}
\label{eq:flattener}
\delta_k \longleftarrow \bigl(1-\sigma_{i+\frac12,j}\bigr)\,\delta_k,
\qquad
\sigma_{i+\frac12,j}=\max\bigl(s^{p}_{i,j},\,s^{p}_{i+1,j}\bigr),
\end{equation}
\begin{equation}
\label{eq:sp}
s^{p}_{i,j}=\min\!\left(1,\;\frac{1}{\varsigma}
\max_{(k,l)\in\mathcal N(i,j)}
\frac{\bigl|p_{k,l}-p_{i,j}\bigr|}{\min\bigl(p_{k,l},p_{i,j}\bigr)}\right),
\qquad \varsigma=\tfrac12,
\end{equation}
where $\mathcal N(i,j)$ are the four face neighbours of the cell.

A contact discontinuity and a rotational discontinuity are both
pressure-continuous, so $s^{p}$ does not respond to either, however large the
jump in density or in the transverse field; it responds to a shock, across which
the pressure jump is the defining feature. The parameter $\varsigma$ is a
saturation scale rather than a threshold: the sensor rises linearly from zero
and reaches one at a relative pressure jump of $50\%$, so there is no dead zone
and no discontinuous switch, and where it saturates the flux is the plain HLL
flux~\eqref{eq:hll}. The design order is unaffected, since in smooth flow
$\sigma=\mathcal O(h)$ multiplies an anti-diffusive term whose own jump is
$\mathcal O(h^{3})$ after the reconstruction of
Section~\ref{sec:reconstruction}.


\subsection{Characteristic WENO3-Z reconstruction}
\label{sec:reconstruction}

Third-order accuracy of the explicit terms is obtained by a characteristic
WENO3-Z reconstruction~\cite{DonBorges2013} with a monotonicity-preserving
limiter, performed on the cell averages $\mathbf q_i$ of the primitive variables.
These are computed from the conservative cell averages through point values
(deconvolution, pointwise conversion, convolution) with the operators of
Section~\ref{sec:fv_ops}. We describe the left state at the face $i+\tfrac12$;
the right state follows by symmetry and the $y$-direction by the permutation
$(x,y,z)\to(y,z,x)$.

The jumps of the primitive variables are projected onto the split characteristic
basis, reconstructed, and mapped back,
\begin{equation}
\label{eq:incr_proj}
\mathbf q^{L}_{i+\frac12}=\mathbf q_i+\RR\,\mathrm{incr},
\qquad
\mathrm{incr}_k=\half\big(w_0\,\Delta_-c_k+w_1\,\Delta_+c_k\big),
\qquad
\Delta_\pm c=\LL\,\Delta_\pm\mathbf q,
\end{equation}
where $\RR$ and $\LL$ are the split eigenvector operators at the arithmetic mean
of the two adjacent states and at the interface parameters~\eqref{eq:eps_face},
with $\epsM$ bounded from below by $10^{-2}$ in the basis only. The block
decomposition of Lemma~\ref{lem:blocks} reduces the inversion to the
five-dimensional aligned sector and removes the $B_\perp\to0$ degeneracy.
With ideal weights
$d_0=\tfrac13$, $d_1=\tfrac23$ and $\beta_0=(\Delta_-c_k)^2$,
$\beta_1=(\Delta_+c_k)^2$, the weights are
\begin{equation}
\label{eq:wenoz}
\tau=|\beta_0-\beta_1|,\qquad
\omega_j=d_j\!\left[1+\left(\frac{\tau}{\beta_j+\epsilon_{w,k}}\right)^{\!2}\right],\qquad
w_j=\frac{\omega_j}{\omega_0+\omega_1},
\qquad
\epsilon_{w,k}=(\Delta x\,S_k)^{2},
\end{equation}
with $S_k$ the magnitude of the $k$-th characteristic variable,
$S_k^{2}=\sum_m(\LL_{km}\max_\Omega|q_m|)^2$. The mesh-dependent regularisation
keeps $w_j=d_j+\mathcal O(h^2)$ also in the neighbourhood of critical points,
where a constant $\epsilon_w$ reduces the three-point scheme to second order.

The face values are then limited componentwise by the monotonicity-preserving
bounds of Suresh and Huynh~\cite{SureshHuynh1997}, applied here to the
third-order value and restricted to the four-point stencil
$\{q_{i-1},q_{i},q_{i+1},q_{i+2}\}$ of the reconstruction: the curvature measure $d_M$
is the minmod of the two available second differences, and it is used in both
the median and the large-curvature bounds. We refer to this variant as MP3. With
$\Delta q_i=q_i-q_{i-1}$ and
$q^{\mathrm{MP}}=q_i+\minmod(q_{i+1}-q_i,\,2\Delta q_i)$, the WENO value
$q^{\mathrm W}$ is accepted if
$(q^{\mathrm W}-q_i)(q^{\mathrm W}-q^{\mathrm{MP}})\le0$; otherwise
\begin{equation}
\label{eq:mp3_cand}
d_M=\minmod\!\big(q_{i-1}-2q_i+q_{i+1},\;q_i-2q_{i+1}+q_{i+2}\big),
\end{equation}
\begin{equation}
\label{eq:mp3_three}
q^{\mathrm{UL}}=q_i+2\Delta q_i,\qquad
q^{\mathrm{MD}}=\half(q_i+q_{i+1})-\half d_M,\qquad
q^{\mathrm{LC}}=q_i+\half\Delta q_i+\tfrac43 d_M,
\end{equation}
\begin{equation}
\label{eq:mp3_bounds}
\begin{aligned}
q^{\min}&=\max\!\big(\min(q_i,q_{i+1},q^{\mathrm{MD}}),\,\min(q_i,q^{\mathrm{UL}},q^{\mathrm{LC}})\big),\\
q^{\max}&=\min\!\big(\max(q_i,q_{i+1},q^{\mathrm{MD}}),\,\max(q_i,q^{\mathrm{UL}},q^{\mathrm{LC}})\big),
\end{aligned}
\qquad
q^{\mathrm{final}}=\median\!\big(q^{\mathrm W},q^{\min},q^{\max}\big).
\end{equation}

\begin{proposition}\label{prop:mp3}
Let $q\in C^3(J)$ on a compact interval $J$ and assume that $q$ has no
degenerate critical point, $|q'(x)|+|q''(x)|>0$ on $J$. Let $q_i$ denote the
cell averages of $q$ and let $q^{\rm W}$ satisfy
$|q^{\rm W}-q(x_{i+1/2})|\le K_{\rm W}h^3$. Then there exists $h_0>0$,
depending only on $\|q\|_{C^3(J)}$, on $\min_J(|q'|+|q''|)$ and on
$K_{\rm W}$, such that for $h<h_0$ the limiter
\eqref{eq:mp3_cand}--\eqref{eq:mp3_bounds} returns $q^{\rm W}$ unchanged in every cell whose
stencil $\{x_{i-1},\dots,x_{i+2}\}$ lies in $J$. In particular the design
order is retained at smooth extrema, wherever they are located relative to
the mesh.
\end{proposition}
 
\begin{proof}
By compactness there is $c_0>0$ such that, at every $x\in J$,
$|q'(x)|\ge c_0$ or $|q''(x)|\ge c_0$. All remainders below are bounded by
$Kh^3$ (or $Kh^2$ where stated), with $K$ depending only on
$\|q\|_{C^3(J)}$ and $K_{\rm W}$. Since the limiter commutes with
$q\mapsto-q$, signs may be fixed without loss of generality. We write
$q'=q'(x_i)$ and $q''=q''(x_i)$.
 
\emph{Case A: $|q'|\ge c_0$; let $q'>0$.} Taylor expansion gives
$q_{i+1}-q_i=hq'+\mathcal O(h^2)$, $2\Delta q_i=2hq'+\mathcal O(h^2)$ and
$q^{\rm W}-q_i=\tfrac12hq'+\mathcal O(h^2)$. Hence
$q^{\rm MP}-q_i=hq'+\mathcal O(h^2)$ and, for $h$ small,
$0<q^{\rm W}-q_i<q^{\rm MP}-q_i$, so that
$(q^{\rm W}-q_i)(q^{\rm W}-q^{\rm MP})<0$ and $q^{\rm W}$ is accepted.
 
\emph{Case B: $|q''|\ge c_0$; let $q''>0$.} Set
\[
\kappa=h^2q''\ \ge\ c_0\,h^2,\qquad
\vartheta=\frac{q'}{h\,q''}\in\mathbb R ,
\]
so that $\vartheta h$ is, to leading order, the signed distance from the
critical point to $x_i$. No bound on $\vartheta$ is assumed. For cell
averages,
$q_{i+k}-q_i=\kappa\bigl(k\vartheta+\tfrac12k^2\bigr)+\mathcal O(h^3)$, and
$q(x_{i+1/2})-q_i
=\kappa\bigl(\tfrac12\vartheta+\tfrac18-\tfrac1{24}\bigr)+\mathcal O(h^3)$.
Therefore, each up to $\mathcal O(h^3)$,
\begin{gather*}
q_{i+1}-q_i=\kappa\bigl(\vartheta+\tfrac12\bigr),\qquad
\Delta q_i=\kappa\bigl(\vartheta-\tfrac12\bigr),\qquad
d_M=\kappa,\qquad
q^{\rm W}-q_i=\kappa\bigl(\tfrac12\vartheta+\tfrac1{12}\bigr),\\[2pt]
q^{\rm UL}-q_i=\kappa\,(2\vartheta-1),\qquad
q^{\rm MD}-q_i=\kappa\bigl(\tfrac12\vartheta-\tfrac14\bigr),\qquad
q^{\rm LC}-q_i=\kappa\bigl(\tfrac12\vartheta+\tfrac{13}{12}\bigr);
\end{gather*}
for $d_M$ this uses that both second differences in \eqref{eq:mp3_cand} equal
$\kappa+\mathcal O(h^3)$ and are therefore positive for $h$ small.
Consequently, for every $\vartheta$,
\begin{equation}\label{eq:mp_ident}
q^{\rm W}=q^{\rm MD}+\tfrac13\kappa=q^{\rm LC}-\kappa
\qquad\text{up to }\mathcal O(h^3).
\end{equation}
 
\emph{Lower bound.} By \eqref{eq:mp_ident},
$\min(q_i,q_{i+1},q^{\rm MD})\le q^{\rm MD}=q^{\rm W}-\tfrac13\kappa$.
Moreover
\[
\min(q_i,q^{\rm UL},q^{\rm LC})\le\min(q_i,q^{\rm UL})
=q^{\rm W}-\kappa\,\max\bigl(\tfrac12\vartheta+\tfrac1{12},\,
\tfrac{13}{12}-\tfrac32\vartheta\bigr)\le q^{\rm W}-\tfrac13\kappa ,
\]
the maximum of the two affine functions being smallest at
$\vartheta=\tfrac12$, where it equals $\tfrac13$. Hence
$q^{\min}\le q^{\rm W}-\tfrac13\kappa$.
 
\emph{Upper bound.} By \eqref{eq:mp_ident},
$\max(q_i,q^{\rm UL},q^{\rm LC})\ge q^{\rm LC}=q^{\rm W}+\kappa$, and
\[
\max(q_i,q_{i+1},q^{\rm MD})\ge\max(q_i,q_{i+1})
=q^{\rm W}+\kappa\,\max\bigl(-\tfrac12\vartheta-\tfrac1{12},\,
\tfrac12\vartheta+\tfrac5{12}\bigr)\ge q^{\rm W}+\tfrac16\kappa ,
\]
the maximum being smallest at $\vartheta=-\tfrac12$, where it equals
$\tfrac16$. Hence $q^{\max}\ge q^{\rm W}+\tfrac16\kappa$.
 
Including the remainders, and since $\min$ and $\max$ are Lipschitz,
\[
q^{\min}\le q^{\rm W}-\tfrac13\kappa+Kh^3,\qquad
q^{\max}\ge q^{\rm W}+\tfrac16\kappa-Kh^3 ,
\]
and $\kappa\ge c_0h^2$ gives $q^{\min}<q^{\rm W}<q^{\max}$ for
$h<c_0/(6K)$. Therefore
$\operatorname{median}(q^{\rm W},q^{\min},q^{\max})=q^{\rm W}$, whether or
not the acceptance test is passed.
\end{proof}
 
% (former Remark 3.3 deleted; one sentence of running text instead)
\noindent The result relies on the specific coefficients $2$ and $\tfrac43$
in \eqref{eq:mp3_three} and on the use of the same curvature measure $d_M$ in
$q^{\rm MD}$ and in $q^{\rm LC}$.


\paragraph{Admissibility.}
\label{sec:zslimiter}
In strongly magnetised regions the internal energy is a small difference of
large kinetic and magnetic contributions, and a small overshoot of the
reconstruction can produce a face state with negative pressure, for which the
Roe average of Section~\ref{sec:hllem} is not defined. Following the first
stage of the Zhang--Shu procedure~\cite{ZhangShu2010}, each limited face value
$\mathbf q^{\mathrm{rec}}$ is therefore contracted towards the cell average,
\begin{equation}
\label{eq:thetascale}
\tilde{\mathbf q}=\bar{\mathbf q}_{i,j}
+\theta\big(\mathbf q^{\mathrm{rec}}-\bar{\mathbf q}_{i,j}\big),
\qquad
\theta=\min_{x\in\{\rho,p\}}
\min\!\Big(1,\ \frac{\bar x_{i,j}-x_{\min}}{\bar x_{i,j}-x^{\mathrm{rec}}}\Big),
\end{equation}
with $x_{\min}\in\{\rho_{\min},p_{\min}\}$, $\rho_{\min}=p_{\min}=10^{-13}$,
$\theta=1$ where $x^{\mathrm{rec}}\ge x_{\min}$ and $\theta=0$ where
$\bar x_{i,j}\le x_{\min}$. Since the reconstruction is performed in primitive
variables, $\rho$ and $p$ are affine in $\theta$, so that
$\rho,p\ge(\rho_{\min},p_{\min})$ holds for $\tilde{\mathbf q}$ whenever it
holds for $\bar{\mathbf q}_{i,j}$. This floor only keeps the face states
admissible for the Riemann solver and is unrelated to the pressure floor
$p_\varepsilon$ of Section~\ref{sec:nestednewton}.
 


\subsection{Implicit magnetic and acoustic sub-systems}
\label{sec:implicit}
After the explicit advective update the density $\rho^{n+1}=\rho^{*}$ is known, and
the residual fluxes $(1-\epsB)\FF^{B}$ and $(1-\epsM)\FF^{p}$ are integrated
implicitly, the magnetic field first and the pressure second. Each sub-system is
reduced to a single linear system by eliminating the momentum: the first one
yields a predictor of the magnetic field, the second one the pressure together
with the divergence-free magnetic field. This section presents both in
semi-discrete form; the discrete operators are given in
Section~\ref{sec:highorder}.

\subsubsection{Magnetic sub-system}
\label{sec:magnetic}
Following the elimination strategy of~\cite{BoscheriThomann2024}, the momentum
equation is substituted into the induction equation, with the gas pressure
gradient taken at the old time level. The quadratic Maxwell stress is linearised about $\mathbf B^{n}$: the magnetic
pressure by its tangent,
$\tfrac12\|\mathbf B\|^{2}\approx\mathbf B^{n}\!\cdot\!\mathbf B-\tfrac12\|\mathbf B^{n}\|^{2}$,
and the tension as $\mathbf B\otimes\mathbf B^{n}$. In this section a star
denotes the state after the explicit update~\eqref{eq:adv_update}, which
already contains the explicit shares $\epsM\mathbf F^p$ and
$\epsB\mathbf F^B$, and the superscript $n$ the old time level, about which
the nonlinear terms are linearised. The momentum entering the induction
equation is thus
\begin{equation}
\label{eq:mag_mom}
(\rho\mathbf u)^{\dagger}=(\rho\mathbf u)^{*}
-\Delta t\,\nabla\!\cdot\!\Bigl[(1-\epsB)\Bigl(
\bigl(\mathbf B^{n}\!\cdot\!\tilde{\mathbf B}-\tfrac12\|\mathbf B^{n}\|^{2}\bigr)\mathbf I
-\tilde{\mathbf B}\otimes\mathbf B^{n}\Bigr)\Bigr]
-\Delta t\,\nabla\bigl[(1-\epsM)\,p^{n}\bigr],
\end{equation}
and the field is advanced by
\begin{equation}
\label{eq:induction}
\tilde{\mathbf B}=\mathbf B^{*}
-\Delta t\,\nabla\!\cdot\!\Bigl[(1-\epsB)\bigl(
\mathbf B^{n}\otimes\mathbf u^{\dagger}-\mathbf u^{\dagger}\otimes\mathbf B^{n}\bigr)\Bigr],
\qquad
\mathbf u^{\dagger}=\frac{(\rho\mathbf u)^{\dagger}}{\rho^{n+1}} .
\end{equation}
Inserting \eqref{eq:mag_mom} into \eqref{eq:induction} eliminates the velocity and
gives a linear system for $\tilde{\mathbf B}$ alone,
\begin{equation}
\label{eq:mag_sys}
\tilde{\mathbf B}-\Delta t^{2}\,\mathcal D_B^{n}\bigl[\tilde{\mathbf B}\bigr]
-\mathcal S_h\bigl[\tilde{\mathbf B}\bigr]=\mathbf b ,
\end{equation}
in which $\mathcal D_B^{n}$ is the composition of the linearised stress divergence in~\eqref{eq:mag_mom}, the division by
$\rho^{n+1}$ and the induction operator in~\eqref{eq:induction},
$\mathbf b$ collects the terms that do not depend on $\tilde{\mathbf B}$, and
$\mathcal S_h$ is the stabilisation term defined below. The system is solved
once with GMRES.

The solution $\tilde{\mathbf B}$ is a predictor: it removes the Alfv\'enic
stiffness, but it is not the magnetic field of the new time level, which is
obtained by constrained transport in the pressure sub-system. Its only use is
the implicit share of the Lorentz force, evaluated with the full quadratic
stress,
\begin{equation}
\label{eq:mag_inter}
(\rho\mathbf u)^{**}=(\rho\mathbf u)^{*}
-\Delta t\,\nabla\!\cdot\!\bigl[(1-\epsB)\,\mathbf F^{B}_{\rho\mathbf{u}}(\tilde{\mathbf B})\bigr].
\end{equation}


\paragraph{Shock-activated stabilisation.}
The central discretisation of $\mathcal D_{B}^{n}$ carries no numerical
dissipation. Following~\cite{MHDIMEX}, a Laplacian is added to the
predictor,
\begin{equation}
\label{eq:mag_stab}
\mathcal S_{h}\big[\tilde\Bb\big]
=\Delta t\sum_{d\in\{x,y\}}\partial_{d}\Big(\chi\,(1-\epsB)\,
\big|\lambda^{B}_{d}\big|_{\max}\,\Delta d\;\partial_{d}\tilde\Bb\Big),
\end{equation}
% --- chi, directly after (71) ---
where $\chi\in[0,1]$ is the shock indicator (flattener) of
Balsara~\cite{Balsara2012}, evaluated on the explicit stage state,
\begin{equation}\label{eq:chi}
\tilde\chi_{i,j}=\min\Bigl[1,\,\max\Bigl(0,\,
-\frac{\min(\Delta x,\Delta y)\,(\nabla\!\cdot\!\mathbf u)_{i,j}+\kappa_1\,c_{i,j}^{\min}}
{\kappa_1\,c_{i,j}^{\min}}\Bigr)\Bigr],
\qquad
\chi_{i,j}=\max_{|k|,|l|\le1}\tilde\chi_{i+k,j+l},
\end{equation}
with central differences for the divergence, $c^{\min}_{i,j}$ the minimum of
$\sqrt{a^2+c_B^2}$ over the $3\times3$ neighbourhood and $\kappa_1=0.4$, and
$|\lambda^B_d|_{\max}=\max_\Omega\tfrac12\big(|u_d|+\sqrt{u_d^2+4c_B^2}\big)$ the
largest magnetic eigenvalue of~\eqref{eq:B_eigs} in direction $d$, evaluated at the
beginning of the time step.

\subsubsection{Constrained transport}
\label{sec:CT}
The solenoidal constraint is enforced by the unstaggered, potential-based
constrained transport of Rossmanith~\cite{Rossmanith2006}. In two dimensions only
the $z$-component $A$ of the vector potential is needed; with
$B_x=\partial_yA$, $B_y=-\partial_xA$ the induction equation for the in-plane
field reduces to the advection equation
\begin{equation}
\label{eq:Aadvect}
\frac{\partial A}{\partial t}+\mathbf{u}\cdot\nabla A=0 ,
\end{equation}
which is advanced in this non-conservative form. The cell averages of $A$ are
updated with the transport term evaluated on the old potential and with the new
velocity,
\begin{equation}
\label{eq:Aupdate}
A^{n+1}_{i,j}=A^{n}_{i,j}-\Delta t\,\Big[\bar u\,G_x+\bar v\,G_y
+\tfrac{\Delta x^{2}}{12}\big(\partial_x\bar u\,\partial_xG_x+\partial_x\bar v\,\partial_xG_y\big)
+\tfrac{\Delta y^{2}}{12}\big(\partial_y\bar u\,\partial_yG_x+\partial_y\bar v\,\partial_yG_y\big)\Big]_{i,j},
\end{equation}
where $\bar{\mathbf{u}}$ is the cell average of the velocity
$(\rho\mathbf{u})^{n+1}/\rho^{n+1}$, $G_x$ and $G_y$ are the upwind derivatives of
$A^{n}$ computed with the fifth-order Hamilton--Jacobi WENO
stencils~\cite{JiangPeng2000} and Z-type weights, and the last two terms are the
covariance that makes the product of two cell averages a fourth-order
approximation of the cell average of the product; the slopes of $G$ are evaluated with
limited differences. The in-plane field follows from the fourth-order discrete
curl
\begin{equation}
\label{eq:curl4}
(B_x)_{i,j}=\frac{-A_{i,j+2}+8A_{i,j+1}-8A_{i,j-1}+A_{i,j-2}}{12\,\Delta y},
\qquad
(B_y)_{i,j}=-\frac{-A_{i+2,j}+8A_{i+1,j}-8A_{i-1,j}+A_{i-2,j}}{12\,\Delta x},
\end{equation}
which satisfies the matching fourth-order central divergence identically. The
initial field is defined through~\eqref{eq:curl4} from the cell averages of the
initial potential. Since~\eqref{eq:Aadvect} is the full induction equation, the
field obtained in this way replaces the in-plane components produced by the
explicit update, and no weighting by $\epsB$ is involved.

This last point is a consistency property of the splitting rather than an
omission, and it is worth making explicit. The only characteristic speed
of~\eqref{eq:Aadvect} is the material velocity. In the potential formulation an
in-plane Alfv\'en wave is a coupled oscillation of $(\rho\mathbf u)$ and $A$ in
which the restoring force sits in the Lorentz term of the momentum equation and
not in the transport of $A$; the Alfv\'enic stiffness of that pair is therefore
carried entirely by $\mathbf F^{B}_{\rho\mathbf{u}}$, which is precisely what
$(1-\epsB)$ sends to the implicit side in~\eqref{eq:mag_mom}
and~\eqref{eq:mag_inter}. Equation~\eqref{eq:Aadvect} may consequently be
advanced explicitly at every value of the parameters, at no cost beyond the
material CFL condition that~\eqref{eq:cfl} already imposes, and the splitting
weights exactly those components in which the Alfv\'en stiffness appears. The
same observation is what makes the potential form preferable to a staggered one
in the present setting: a staggered constrained transport advances the in-plane
field from a face-centred electric field, which would have to be apportioned
between the explicit and implicit shares of the induction flux consistently with
the partitioned tableau, whereas~\eqref{eq:Aadvect} is the physical transport
equation and carries all shares at once.

One consequence has to be stated plainly, because it bounds what the Riemann
solver of Section~\ref{sec:hllem} can be claimed to do. The anti-diffusion
of~\eqref{eq:hllem} acts on the components that the finite volume operator
advances: in two dimensions the momentum, the energy and $B_z$, and in one
dimension the transverse field as well, since the in-plane induction is then
itself a finite volume flux. The in-plane field of a two-dimensional
computation is advanced by~\eqref{eq:Aadvect} instead, and inherits the
anti-diffusion only through the velocity that transports $A$. That path is
real but weak. In the field-loop advection of
Section~\ref{sec:fieldloop} at $A_0=1$ on $128\times64$ cells, a case in which
the splitting realises $\epsM=0$ and $\epsB=1$ simultaneously, switching off
every $\delta_k$ changes $\Bb_\perp$ by $1.2\times10^{-6}$ in relative maximum
norm, against a discretisation error of $8.8\times10^{-3}$ in the same norm,
and leaves the retained magnetic energy and the peak of $\|\Bb_\perp\|$
unchanged to six digits. The sharp resolution of rotational discontinuities
demonstrated in Section~\ref{sec:MHD_1d} is accordingly a one-dimensional
statement about the field; in two dimensions it is a statement about the
momentum, the energy and $B_z$.

For given $A^{n}$ and a given upwind direction, \eqref{eq:Aupdate}
and~\eqref{eq:curl4} define the new in-plane field as an \emph{affine} function
of the velocity. We make this exact by selecting the upwind side of $G_x$, $G_y$
with the velocity of the old time level and by using central differences for the
slopes of $\bar{\mathbf{u}}$; the nonlinear weights then act on $A^{n}$ only. We
denote the resulting map by
\begin{equation}
\label{eq:CTmap}
\Bb_\perp^{n+1}=\mathcal C\big[(\rho\mathbf{u})^{n+1}\big].
\end{equation}

\subsubsection{Pressure sub-system}
\label{sec:pressure}
The pressure sub-system couples the acoustic pressure gradient, the energy
equation and the magnetic field of the new time level. The momentum is
\begin{equation}
\label{eq:p_mom}
(\rho\mathbf{u})^{n+1}=(\rho\mathbf{u})^{\ast\ast}-\Delta t\,\nabla\big[(1-\epsM)\,p^{n+1}\big],
\end{equation}
the in-plane field is given by~\eqref{eq:CTmap}, and the out-of-plane field and
the total energy receive the implicit share of the induction and of the magnetic
energy flux, both evaluated with the new velocity and the old field,
\begin{align}
\label{eq:Bz}
B_z^{n+1}&=B_z^{\ast}-\Delta t\,\nabla\!\cdot\!\big[(1-\epsB)\big(\mathbf{u}^{n+1}B_z^{n}-w^{n+1}\Bb^{n}\big)\big],\\
\label{eq:Umag}
E^{\ast\ast}&=E^{\ast}-\Delta t\,\nabla\!\cdot\!\big[(1-\epsB)\big(\|\Bb^{n}\|^{2}\mathbf{u}^{n+1}
-(\mathbf{u}^{n+1}\!\cdot\Bb^{n})\,\Bb^{n}\big)\big],
\end{align}
where a star denotes the value after the explicit update. With
$\rho e=p/(\gamma-1)$ the energy equation reads
\begin{equation}
\label{eq:energy_imp}
\frac{p^{n+1}}{\gamma-1}+k^{n+1}+m^{n+1}
=E^{\ast\ast}-\Delta t\,\nabla\!\cdot\!\big[(1-\epsM)\,h^{n}(\rho\mathbf{u})^{n+1}\big],
\end{equation}
in which the kinetic and magnetic energies are linearised about the old time
level,
\begin{equation}
\label{eq:klin}
k^{n+1}=\frac{(\rho\mathbf{u})^{n}\!\cdot\!(\rho\mathbf{u})^{n+1}-\half\|(\rho\mathbf{u})^{n}\|^{2}}{\rho^{n+1}},
\qquad
m^{n+1}=\Bb^{n}\!\cdot\!\Bb^{n+1}-\half\|\Bb^{n}\|^{2}.
\end{equation}
By~\eqref{eq:p_mom} the momentum is an affine function of $p^{n+1}$, and so are,
through~\eqref{eq:CTmap}, \eqref{eq:Bz}, \eqref{eq:Umag} and~\eqref{eq:klin},
the magnetic field, $E^{\ast\ast}$, $k^{n+1}$ and $m^{n+1}$. Writing
$\Phi[\mathbf m]=E^{\ast\ast}[\mathbf m]-k[\mathbf m]-m[\mathbf m]$ for the
dependence of these terms on the momentum, \eqref{eq:energy_imp} becomes a single
linear equation for the pressure,
\begin{equation}
\label{eq:pwave}
\frac{p^{n+1}}{\gamma-1}
-\Delta t^{2}\,\nabla\!\cdot\!\Big[(1-\epsM)\,h^{n}\,\nabla\big((1-\epsM)\,p^{n+1}\big)\Big]
-\Phi\Big[(\rho\mathbf{u})^{\ast\ast}-\Delta t\,\nabla\big((1-\epsM)p^{n+1}\big)\Big]
=-\Delta t\,\nabla\!\cdot\!\big[(1-\epsM)\,h^{n}(\rho\mathbf{u})^{\ast\ast}\big].
\end{equation}
The constrained transport, the induction of $B_z$ and the magnetic energy flux
are thus part of the pressure operator: every evaluation of the left-hand side
computes the divergence-free field associated with the trial pressure, so that
the energy is balanced with a solenoidal field throughout the solve. The
operator is not symmetric, because of the first-order terms contributed by
$\Phi$, and~\eqref{eq:pwave} is solved as described in Section \ref{sec:nestednewton}.

The central discretisation of the elliptic term in~\eqref{eq:pwave} carries no
numerical dissipation, exactly as that of $\mathcal D^{n}_{B}$ did, so the same
shock-activated stabilisation is applied to it. Following~\cite{MHDIMEX},
\eqref{eq:pwave} is augmented on its left-hand side by
$-\mathcal S^{p}_{h}[p^{n+1}]$, with
\begin{equation}
\label{eq:p_stab}
\mathcal S^{p}_{h}\big[p^{n+1}\big]
=\Delta t\sum_{d\in\{x,y\}}\partial_{d}\Big(
\mathbf V'(p^{n})\,\chi\,(1-\epsM)\,\big|u_{d}\big|_{\max}\,\Delta d\;
\partial_{d}p^{n+1}\Big),
\qquad
\big|u_{d}\big|_{\max}=\max_{\Omega}|u_{d}| ,
\end{equation}
where $\chi$ is the indicator~\eqref{eq:chi} already used
in~\eqref{eq:mag_stab}, and $\mathbf V'(p)=1/(\gamma-1)$ for an ideal gas. Three
remarks. The coefficient is the largest material speed of that direction over
the mesh, as in~\eqref{eq:mag_stab}, rather than the local one; it is evaluated
at the beginning of the time step. The weight $(1-\epsM)$ confines the term to
the share of the pressure flux that is treated implicitly, so it vanishes with
the implicit sub-system itself as $\epsM\to1$. Keeping $\mathbf V'$ inside the
divergence preserves the strict flux form of the dissipation, so that the term
enters the coupling matrix of Section~\ref{sec:nestednewton} as a conservative
contribution and does not disturb its M-matrix structure. Because the Laplacian
is discretised implicitly it imposes no time-step restriction beyond~\eqref{eq:cfl}.

Once $p^{n+1}$ is known, the
momentum follows from~\eqref{eq:p_mom}, the magnetic field
from~\eqref{eq:CTmap} and~\eqref{eq:Bz}, and the total energy from
\begin{equation}
\label{eq:Unew}
E^{n+1}=\frac{p^{n+1}}{\gamma-1}+k^{n+1}+m^{n+1}.
\end{equation}

\paragraph{Conservation.}
The new total energy~\eqref{eq:Unew} is \emph{defined} as the left-hand side
of the solved equation~\eqref{eq:energy_imp}, with the same linearised energies
$k^{n+1}$ and $m^{n+1}$. Hence
\begin{equation}
\label{eq:cons_energy}
E^{n+1}=E^{\ast\ast}-\Delta t\,\nabla\!\cdot\!\big[(1-\epsM)\,h^{n}(\rho\mathbf{u})^{n+1}\big],
\end{equation}
and since $E^{\ast\ast}$ is itself obtained from $E^{n}$ by flux differences, the
total energy is conserved. A second pressure solve would be necessary only if the
energy were re-assembled from the quadratic energies
$\|(\rho\mathbf{u})^{n+1}\|^{2}/(2\rho^{n+1})$ and $\half\|\Bb^{n+1}\|^{2}$, which
are not the quantities appearing in the solved equation. Mass and momentum are
updated in flux form by~\eqref{eq:step_adv}, \eqref{eq:mag_inter}
and~\eqref{eq:p_mom}, and are therefore conserved as well.

An implicit step thus solves for the predictor $\tilde\Bb$ and then for the
pressure, without any iteration between the two sub-systems and without a second
pressure solve. The pressure system is itself solved by the defect correction
of Section~\ref{sec:nestednewton}, in which the constrained transport, the
induction of $B_z$ and the magnetic energy flux are carried by the defect rather
than by the inverted operator; what is avoided here is the outer iteration
between the magnetic field and the pressure, not iteration as such. Every linearisation
above---in~\eqref{eq:mag_mom}, in the transport of $A$ and in~\eqref{eq:klin}---is
of the form $H(\mathbf y^{n},\mathbf y^{n+1})$ with
$H(\mathbf y,\mathbf y)$ equal to the original nonlinear term, which is the
structure required by the partitioned IMEX Runge--Kutta schemes of
Section~\ref{sec:time}; no accuracy is lost with respect to an iterated
treatment. The predictor $\tilde\Bb$ and the field $\Bb^{n+1}$ are two consistent
discretisations of the same induction equation: they differ through the
pressure level in the advecting velocity and, in two dimensions, through the
discretisation of the in-plane field, which for $\Bb^{n+1}$ is the constrained
transport of Section~\ref{sec:CT}. Both differences are of truncation order,
so that the Lorentz force in~\eqref{eq:mag_inter} is consistent to the order of
the scheme.


\subsubsection{Enforcement of pressure positivity}
\label{sec:nestednewton}

After spatial discretisation the pressure equation~\eqref{eq:pwave} has the form
$\underline{\mathbf V}(\underline{\mathbf p})
+\underline{\underline{\mathbf T}}\,\underline{\mathbf p}=\underline{\mathbf b}$,
with the cell-local relation $V_i(p)=p/(\gamma-1)$. Positivity of the pressure is
enforced with the mechanism of~\cite{MHDIMEX}: the relation
$\underline{\mathbf V}$ is replaced by a modified relation
$\underline{\mathbf V}_{\mathrm{mod}}$ with a steep slope below a floor
$p_\varepsilon$, and the resulting system is solved by the nested Newton method.
That mechanism requires a matrix with non-positive off-diagonal entries, which
the fourth-order operator $\underline{\underline{\mathbf T}}$ is not. We therefore apply it to a second-order operator $\underline{\mathbf T}_2$
with that property: the compact five-point discretisation of the acoustic
term of~\eqref{eq:pwave}, with upwind differences for the convective term
$-\Delta t\,\mathbf u^{n}\!\cdot\!\nabla p$ arising from the linearised
kinetic energy. The constrained transport, the induction of $B_z$ and the
magnetic energy flux are not part of $\underline{\mathbf T}_2$; they enter
only through the defect. The fourth-order solution is recovered by defect
correction,
\begin{equation}
\label{eq:defect}
\underline{\mathbf V}_{\mathrm{mod}}\big(\underline{\mathbf p}^{(k+1)}\big)
+\underline{\underline{\mathbf T}}_{2}\,\underline{\mathbf p}^{(k+1)}
=\underline{\mathbf b}
+\big(\underline{\underline{\mathbf T}}_{2}-\underline{\underline{\mathbf T}}\big)\,
\underline{\mathbf p}^{(k)},
\qquad k=0,1,\dots
\end{equation}
Every iterate obeys the lower bound $\theta\,p_\varepsilon$, with
$\theta=0.9$ and $p_\varepsilon=10^{-7}$ in all computations, and the converged pressure satisfies the fourth-order equation wherever the floor is not active. The total energy is
assembled from~\eqref{eq:Unew} with the physical relation $V(p)$.

% \begin{remark}[Scope]
% \label{rem:posscope}
% In~\cite{MHDIMEX} the same nested Newton machinery serves two purposes at once:
% it accommodates a general nonlinear equation of state, for which
% $\underline{\mathbf V}$ is genuinely nonlinear, and it enforces positivity. Here
% the thermodynamic closure is the ideal gas law, so $\underline{\mathbf V}$ is
% linear and the elliptic system would be linear; the iteration is present
% \emph{solely} because of the positivity modification~\eqref{eq:Vmod}. The extension
% to a general equation of state requires no further machinery and follows
% directly~\cite{MHDIMEX}.
% \end{remark}

% \begin{remark}[Operator-family pairing]
% \label{rem:composed}
% The elliptic operator of~\eqref{eq:pwave} is not assembled independently: it is
% \emph{derived} by inserting the momentum update~\eqref{eq:p_mom} into the energy
% balance~\eqref{eq:energy_imp}. At the discrete level this means that the three
% ingredients---the elliptic left-hand side, the enthalpy-divergence right-hand side,
% and the momentum update---must be drawn from one and the same operator family, with
% $\epsM$ evaluated at faces throughout; the fourth-order family used in this work is
% given in Section~\ref{sec:highorder}. Violating the pairing breaks the discrete
% energy--momentum consistency even when every ingredient is individually of the
% nominal order. In our experiments, replacing the divergence-form momentum update by
% an independent fourth-order cell-centred pressure gradient produced a systematic
% spurious compression plateau ($\rho\approx1.23$) precisely in the $\epsM$
% transition zone of the Brio--Wu rarefaction, where the
% $\mathcal O(\Delta t\,\Delta x^{2})$ mismatch between the two gradient
% discretisations accumulates. Componentwise convergence tests cannot detect this
% failure mode, since they probe each formula in isolation rather than the
% consistency of the pair. The same principle fixes the compact three-point stencil at
% second order: with spatially varying $\epsM$ the compact form is no longer the
% composition of the first-order building blocks, and the composed operator must be
% used instead.
% \end{remark}

% \subsection{Enforcement of the divergence-free constraint}
% \label{sec:CT}

% Preserving $\nabla\!\cdot\!\mathbf B=0$ at the discrete level is essential:
% divergence errors generate spurious forces and unphysical states that are
% particularly damaging in the low-Mach/low-Alfv\'en regimes targeted here. In
% this work we implement two fourth-order adaptations of the constrained-transport
% (CT) methodology: a staggered adaptation in the spirit of Gardiner and
% Stone~\cite{gardiner2005}, and an unstaggered, potential-based adaptation of the
% technique introduced by Rossmanith~\cite{rossmanith2006MPACT}. Both share the
% same fourth-order electric-field construction, which we describe first.

% \paragraph{Fourth-order interpolation operators.}
% Let $\mathcal I^x_4$ denote the fourth-order face-interpolation operator,
% \begin{equation}
% \label{eq:I4}
% \mathcal I^x_4[q]_{i+\frac12,j}
%  =\tfrac{1}{16}\left(-q_{i-1,j}+9\,q_{i,j}+9\,q_{i+1,j}-q_{i+2,j}\right),
% \end{equation}
% with $\mathcal I^y_4$ defined analogously. For the dissipative part we use
% \emph{biased} fourth-order interface values, e.g.\ in $x$,
% \begin{equation}
% \label{eq:biased4}
% q^{L}_{i+\frac12,j}=\tfrac{1}{16}\!\left(q_{i-2,j}-5q_{i-1,j}+15q_{i,j}+5q_{i+1,j}\right),
% \qquad
% q^{R}_{i+\frac12,j}=\tfrac{1}{16}\!\left(q_{i+3,j}-5q_{i+2,j}+15q_{i+1,j}+5q_{i,j}\right),
% \end{equation}
% so that the jump $q^R-q^L$ is $\mathcal O(\Delta x^4)$ for smooth data but
% $\mathcal O(1)$ on grid-scale (odd--even) modes.

% \paragraph{Fourth-order face electric field.}
% With the cell-centred electric field
% $\widetilde E_{z,i,j}=v_{i,j}(B_x)_{i,j}-u_{i,j}(B_y)_{i,j}$, the face values
% read
% \begin{equation}
% \label{eq:Ezface}
% (E_z)_{i+\frac12,j}=\mathcal I^x_4[\widetilde E_z]_{i+\frac12,j}
%  +\tfrac12\,s^{x}_{i+\frac12,j}\!\left(B_y^{R}-B_y^{L}\right),
% \qquad
% (E_z)_{i,j+\frac12}=\mathcal I^y_4[\widetilde E_z]_{i,j+\frac12}
%  -\tfrac12\,s^{y}_{i,j+\frac12}\!\left(B_x^{R}-B_x^{L}\right),
% \end{equation}
% where the transverse-field jumps are evaluated with~\eqref{eq:biased4} and the
% dissipation speeds are
% \begin{equation}
% \label{eq:ctspeed}
% s^{x}_{i+\frac12,j}=\max_{k\in\{i,i+1\}}
% \left(|u_{k,j}|+\frac{|(B_x)_{k,j}|}{\sqrt{\rho_{k,j}}}\right),
% \end{equation}
% and analogously in $y$. Two design choices deserve comment. First, the
% dissipation speed~\eqref{eq:ctspeed} includes the normal Alfv\'en speed and not
% merely the material velocity: the characteristic speeds of the induction
% equation are $v_n\pm c_a$, so a purely material coefficient vanishes in regions
% where $\mathbf v=0$ but $\mathbf B\neq 0$, leaving the face-field update as an
% undamped central difference whose odd--even modes grow slowly into large-scale
% spurious structures. Second, because the dissipation acts on the biased
% fourth-order jumps~\eqref{eq:biased4}, it is $\mathcal O(\Delta x^4)$ on smooth
% solutions---the electric field construction is genuinely fourth-order
% accurate---while retaining an $\mathcal O(1)$ damping rate on grid-scale
% decoupling modes, the same mechanism by which the WENO-based Riemann solver of
% Section~\ref{sec:hllem} combines high order with robustness.

% The corner electric field is obtained by fourth-order transverse interpolation
% of the face values, averaged over the two available constructions,
% \begin{equation}
% \label{eq:Ezcorner}
% (E_z)_{i+\frac12,j+\frac12}
%  =\tfrac12\Big(\mathcal I^y_4\big[(E_z)_{i+\frac12,\cdot}\big]_{j+\frac12}
%               +\mathcal I^x_4\big[(E_z)_{\cdot,j+\frac12}\big]_{i+\frac12}\Big),
% \end{equation}
% which supersedes the second-order gradient-correction of Gardiner and
% Stone~\cite{gardiner2005}.

% \subsubsection{Staggered adaptation of the constrained-transport method}
% \label{sec:CTstag}
% Staggered magnetic-field components $(B_x)_{i+\frac12,j}$ and
% $(B_y)_{i,j+\frac12}$ are stored at the cell faces, initialised there at $t=0$,
% and evolved by a discrete Faraday law using the corner electric
% field~\eqref{eq:Ezcorner},
% \begin{equation}
% \label{eq:faraday}
% \begin{aligned}
% (B_x)^{n+1}_{i+\frac12,j}&=(B_x)^{n}_{i+\frac12,j}
%  +\frac{\Delta t}{\Delta y}\Big[(E_z)_{i+\frac12,j+\frac12}-(E_z)_{i+\frac12,j-\frac12}\Big],\\
% (B_y)^{n+1}_{i,j+\frac12}&=(B_y)^{n}_{i,j+\frac12}
%  -\frac{\Delta t}{\Delta x}\Big[(E_z)_{i+\frac12,j+\frac12}-(E_z)_{i-\frac12,j+\frac12}\Big].
% \end{aligned}
% \end{equation}
% Equation~\eqref{eq:faraday} is a discrete Stokes identity: the staggered
% divergence
% \begin{equation}
% \label{eq:divstag}
% [\nabla\!\cdot\!\mathbf B]_{i,j}
% =\frac{(B_x)_{i+\frac12,j}-(B_x)_{i-\frac12,j}}{\Delta x}
% +\frac{(B_y)_{i,j+\frac12}-(B_y)_{i,j-\frac12}}{\Delta y}=0
% \end{equation}
% is preserved \emph{exactly}, independently of the accuracy of the electric
% field, so the fourth-order construction~\eqref{eq:Ezface}--\eqref{eq:Ezcorner}
% raises the accuracy of the field evolution without touching the constraint.
% The magnetic field returned by the implicit magnetic sub-system of
% Section~\ref{sec:implicit} serves only as a predictor for the momentum update
% and is then replaced at cell centres by the fourth-order face-to-centre
% average of the staggered field,
% \begin{equation}
% \label{eq:face2cell}
% (B_x)^{n+1}_{i,j}=\tfrac{1}{16}\Big(\!-(B_x)_{i-\frac32,j}+9(B_x)_{i-\frac12,j}
%  +9(B_x)_{i+\frac12,j}-(B_x)_{i+\frac32,j}\Big),
% \end{equation}
% and analogously for $B_y$; the two-point average of~\cite{gardiner2005} would
% cap the two-dimensional accuracy of the magnetic field at second order.

% \subsubsection{Unstaggered adaptation of the constrained-transport method}
% \label{sec:CTunstag}
% The staggered formulation requires a dual representation of the magnetic field
% at faces and centres. To avoid it we also implement a fully cell-centred,
% potential-based adaptation following Rossmanith~\cite{rossmanith2006MPACT}: in
% two dimensions only the $z$-component $A$ of the vector potential is tracked,
% stored at cell centres, and never overwritten by the predictor field. Using
% $\partial_t A=-E_z$ (equivalent to the advection form
% $\partial_t A+\mathbf{u}\!\cdot\!\nabla A=0$ since $E_z=-(\mathbf{u}\times\mathbf B)_z$ and
% $\mathbf B=\nabla\times A\hat{\mathbf z}$), the potential is advanced with the
% cell-centred electric field obtained from the corner
% values~\eqref{eq:Ezcorner} by tensor-product fourth-order interpolation,
% \begin{equation}
% \label{eq:Aupdate}
% A^{n+1}_{i,j}=A^{n}_{i,j}-\Delta t\,
% \mathcal I^x_4\mathcal I^y_4\big[(E_z)_{\cdot+\frac12,\cdot+\frac12}\big]_{i,j}
% +\Delta t\,\mathcal D^x\!\big[\tfrac12 s^x (A^R\!-\!A^L)\big]_{i,j}
% +\Delta t\,\mathcal D^y\!\big[\tfrac12 s^y (A^R\!-\!A^L)\big]_{i,j},
% \end{equation}
% so that the potential inherits the fourth-order accuracy and the built-in
% grid-scale damping of the electric-field construction. The magnetic field is
% recovered by the \emph{fourth-order} discrete curl
% \begin{equation}
% \label{eq:curl4}
% (B_x)^{n+1}_{i,j}=\frac{-A^{n+1}_{i,j+2}+8A^{n+1}_{i,j+1}-8A^{n+1}_{i,j-1}+A^{n+1}_{i,j-2}}{12\,\Delta y},
% \qquad
% (B_y)^{n+1}_{i,j}=-\frac{-A^{n+1}_{i+2,j}+8A^{n+1}_{i+1,j}-8A^{n+1}_{i-1,j}+A^{n+1}_{i-2,j}}{12\,\Delta x},
% \end{equation}
% which satisfies the \emph{matching} fourth-order central divergence identically,
% \begin{equation}
% \label{eq:div4}
% [\nabla\!\cdot\!\mathbf B]^{(4)}_{i,j}
% =\mathcal D^x_4[B_x]_{i,j}+\mathcal D^y_4[B_y]_{i,j}=0,
% \end{equation}
% with $\mathcal D_4$ the five-point fourth-order central difference, since the
% mixed fourth-order differences of $A$ commute and cancel exactly. We stress
% that the divergence-free property is tied to the pair of operators: measuring
% the field~\eqref{eq:curl4} with a second-order divergence returns the
% $\mathcal O(\Delta x^2)$ truncation difference between the two difference
% operators, which reflects the measuring stick rather than the field. If the
% initial data are given in terms of $\mathbf B$ alone, the initial potential is
% obtained by solving a Poisson problem as in~\cite{rossmanith2006MPACT}, and the
% initial cell-centred field is re-initialised through~\eqref{eq:curl4} so that
% the discrete constraint holds from $t=0$. Boundary conditions for $A$ follow
% those of the magnetic field: periodic where $\mathbf B$ is periodic, and
% first-order extrapolated where a zeroth-order extrapolated field is used.



% ============================================================================
%  Sec.\ref{sec:CT} -- revised against the standard-IMEX version.
%
%  IMPORTED from standard-IMEX:
%   (a) the literature survey (Powell / Brackbill-Barnes / Dedner / Evans-Hawley);
%   (b) the "two aspects worth stating explicitly" framing -- the involution is
%       enforced after all sub-systems, the cell-centred field is a predictor,
%       and rho^{n+1}, v^{n+1} are already available so E sits at the new level;
%   (c) the initialisation requirement for the staggered field;
%   (d) the Toth cell-centred divergence eq:CT-div-toth.
%
%  STRENGTHENED (these are the parts that differ from standard-IMEX):
%   (e) the flux-dual paragraph now states EXPLICITLY that the standard-IMEX
%       simplification -- induction rows of the advective and pressure fluxes
%       vanish identically -- FAILS under the eps-splitting.  That is the new
%       difficulty and it deserves to be named as such;
%   (f) the post-CT energy re-assembly is promoted from a trailing sentence to
%       a Remark stating a structural requirement.
%
%  LABELS REPAIRED:  eq:splitflux -> eq:Fc_explicit ; eq:magsys -> eq:mag_sys ;
%                    sec:flattener -> sec:stabilisation ; eq:stagger4 -> the
%                    operator family of sec:highorder.
%  CITATION KEYS: standard-IMEX uses GardinerStone2005 / Rossmanith2006 /
%                 Dematte2024 ; this draft uses gardiner2005 / rossmanith2006MPACT.
%                 Unify before submission -- below the standard-IMEX keys are used.
% ============================================================================


% \subsection{Enforcement of the divergence-free constraint}
% \label{sec:CT}
 
% As anticipated in Section~\ref{sec:governing}, the solenoidal condition
% $\nabla\!\cdot\!\Bb=0$ is an involution of the ideal MHD system: if it holds for
% the initial data, the evolution preserves it for all later times at the continuous
% level. At the discrete level this property is not inherited automatically, and the
% divergence errors generated by a generic discretisation are known to produce
% spurious oscillations, incorrect shock speeds and unphysical
% states~\cite{BrackbillBarnes1980,Toth2000}---all of which are particularly damaging
% in the low-Mach and low-Alfv\'en regimes targeted here, where the magnetic pressure
% $\half\|\Bb\|^{2}$ can exceed the gas pressure by orders of magnitude, so that a
% small error in $\Bb$ translates into a large spurious force. Several strategies
% have been devised to control or eliminate these errors. The eight-wave formulation
% of Powell et al.~\cite{Powell1999} augments the system with source terms
% proportional to the divergence error, at the price of losing strict conservation.
% In the projection method of Brackbill and Barnes~\cite{BrackbillBarnes1980} the
% magnetic field produced by the base scheme is projected onto the space of
% solenoidal fields through an auxiliary Poisson problem. Hyperbolic divergence
% cleaning~\cite{Dedner2002} couples the constraint to the conservation laws through
% a generalised Lagrange multiplier, which damps and advects the divergence error
% towards the boundaries. The most widely adopted alternative is the constrained
% transport (CT) approach originally introduced by Evans and
% Hawley~\cite{EvansHawley1988}, in which the magnetic field is represented on the
% cell interfaces and the electric field at the cell corners; a discrete Faraday law
% then preserves a face-based divergence operator down to machine accuracy. We refer
% to the review of T\'oth~\cite{Toth2000} and
% to~\cite{BalsaraSpicer1999,GardinerStone2005,Rossmanith2006} for a detailed
% comparison of these families.
 
% In the present work we implement two constrained transport adaptations: a staggered
% one following Gardiner and Stone~\cite{GardinerStone2005}, used as the default, and
% an unstaggered potential-based one following Rossmanith~\cite{Rossmanith2006}. Two
% aspects of their use within the present algorithm are worth stating explicitly.
% First, since $\nabla\!\cdot\!\Bb=0$ is an involution rather than a member of any of
% the three flux sub-systems of Section~\ref{sec:3split}, it is enforced after all
% sub-systems have been solved, and the correction acts on the magnetic field alone:
% the cell-centred field $\tilde\Bb^{n+1}$ delivered by the implicit magnetic
% sub-system of Section~\ref{sec:magnetic} is regarded as a \emph{predictor}, used to
% advance the momentum and the energy, and is subsequently replaced. Second, because
% the correction is applied at the end of the step, the density $\rho^{n+1}$ and the
% velocity $\mathbf{u}^{n+1}$ resulting from the pressure sub-system are already available,
% so that the transport step is driven by data at the new time level. The two
% adaptations differ in how the induction is discretised: the staggered variant
% performs a discrete Stokes update on the cell faces, driven by an electromotive
% force which, as shown next, must be taken as the dual of the numerical fluxes; the
% unstaggered variant advances the vector potential by the equivalent scalar
% advection equation, for which that construction does not arise.
 
% \subsubsection{Staggered adaptation}
% \label{sec:CTstag}
 
% \paragraph{Why the electric field must be the flux dual.}
% In a fully explicit Godunov scheme the induction rows of the numerical flux
% \emph{are} the electromotive force, and a CT step built from them reproduces
% exactly the magnetic field the scheme would have produced, while additionally
% enforcing~\eqref{eq:divstag}; this is the duality exploited in~\cite{Dematte2024}.
% In the single-parameter three-way splitting the same simplification is available
% for a different reason: the induction rows of the advective and of the pressure
% flux vanish identically, so the entire numerical flux of the magnetic field is
% carried by the implicit magnetic sub-system and no contribution has to be
% accumulated from elsewhere.
 
% \emph{Under the $\eps$-splitting this is no longer true.} The explicit sub-system
% $\tilde\FF^{c}$ of~\eqref{eq:Fc_explicit} contains the share $\epsB\FF^{B}$, whose
% induction rows are non-zero, so that $\Bb$ is now advanced in three separate
% places:
% \begin{enumerate}
% \item explicitly, through the share $\epsB\FF^{B}$ carried by $\tilde\FF^{c}$ and
%       evaluated by the AP-HLLEM solver of Section~\ref{sec:hllem};
% \item implicitly, through the right-hand side $\Bb^{\ast}$ of the magnetic wave
%       system~\eqref{eq:mag_sys}, weighted by $(1-\epsB)$ at each face;
% \item implicitly, through the operator itself---the magnetic wave term, which
%       carries $(1-\epsB)^{2}$ because the momentum update and the induction
%       equation each contribute one factor, and the shock-activated Laplacian of
%       Section~\ref{sec:stabilisation}, which carries $(1-\epsB)$.
% \end{enumerate}
% A CT step driven by any proper subset of these overwrites $\Bb$ with a field
% inconsistent with the momentum and the energy that were built from the predictor.
% We therefore assemble the face electromotive force from the \emph{total}
% increment. Writing $\mathcal F^{\,\mathrm{exp}}$, $\mathcal F^{\,\mathrm{rhs}}$ and
% $\mathcal F^{\,\mathrm{imp}}$ for the three face fluxes of the induction rows, and
% using the sign convention $E_z=vB_x-uB_y$, for which the $x$-flux of $B_y$ equals
% $-E_z$ and the $y$-flux of $B_x$ equals $+E_z$,
% \begin{equation}
% \label{eq:emfdual}
% (E_z)_{i+\frac12,j}=-\Big(\mathcal F^{\,\mathrm{exp}}
% +\mathcal F^{\,\mathrm{rhs}}+\mathcal F^{\,\mathrm{imp}}\Big)_{i+\frac12,j},
% \qquad
% (E_z)_{i,j+\frac12}=+\Big(\mathcal G^{\,\mathrm{exp}}
% +\mathcal G^{\,\mathrm{rhs}}+\mathcal G^{\,\mathrm{imp}}\Big)_{i,j+\frac12}.
% \end{equation}
% No dissipation is added to~\eqref{eq:emfdual}: the upwind dissipation of the
% explicit share already resides in the HLLEM flux, and that of the implicit share in
% the Laplacian of Section~\ref{sec:stabilisation}. In particular the concern that
% motivates a separate dissipation speed in classical CT---that a purely material
% coefficient vanishes where $\mathbf{u}=\mathbf 0$ but $\Bb\neq\mathbf 0$, leaving the
% face-field update as an undamped central difference---does not arise here, because
% both the Einfeldt speeds of the HLLEM solver and the coefficient of the implicit
% Laplacian contain the Alfv\'en speed.
 
% \begin{remark}[Three fluxes, three divergence operators]
% \label{rem:ctoperators}
% The three shares in~\eqref{eq:emfdual} are \emph{not} differenced by the same
% operator in the scheme itself. The explicit share is a finite-volume flux, for which
% the two-point difference is exact; the right-hand side $\Bb^{\ast}$ uses the
% fourth-order staggered divergence $\mathcal D_4$ of Section~\ref{sec:highorder},
% since it carries a single power of $\Delta t$ and therefore requires
% $\mathcal O(\Delta x^{3})$ or better; the operator share uses the compact two-point
% form, which suffices because it carries $\Delta t^{2}$. A single corner
% electromotive force cannot carry three different difference operators. Two
% consistent resolutions are available: perform one Stokes update per operator, each
% of which preserves~\eqref{eq:divstag} separately so that their sum does as well; or
% accept a single two-point update and regard CT as \emph{replacing} rather than
% reproducing the predictor, in which case the two discretisations of the
% $\Bb^{\ast}$ share differ by $\mathcal O(\Delta x^{4})$ on smooth data. We adopt the
% second option, whose structural consequence is stated in Remark~\ref{rem:ctenergy}.
% Note that the explicit share must in either case be differenced with the two-point
% operator: replacing it by $\mathcal D_4$ would introduce an
% $\mathcal O(\Delta t\,\Delta x^{2})=\mathcal O(\Delta x^{3})$ error and cap the
% scheme at second order.
% \end{remark}
 
% \paragraph{Corner electric field.}
% The corner value is obtained by fourth-order transverse interpolation of the face
% values~\eqref{eq:emfdual}, averaged over the two available constructions,
% \begin{equation}
% \label{eq:Ezcorner}
% (E_z)_{i+\frac12,j+\frac12}
%  =\half\Big(\mathcal I^y_4\big[(E_z)_{i+\frac12,\cdot}\big]_{j+\frac12}
%               +\mathcal I^x_4\big[(E_z)_{\cdot,j+\frac12}\big]_{i+\frac12}\Big),
% \qquad
% \mathcal I^x_4[q]_{i+\frac12}=\tfrac{1}{16}\big(-q_{i-1}+9q_{i}+9q_{i+1}-q_{i+2}\big),
% \end{equation}
% which replaces the gradient-corrected four-point average of Gardiner and
% Stone~\cite{GardinerStone2005}, for the reason given in Remark~\ref{rem:ctaccuracy}.
 
% \paragraph{Faraday update.}
% Staggered components $(B_x)_{i+\frac12,j}$ and $(B_y)_{i,j+\frac12}$ are stored at
% the cell faces and evolved from the stage base state by a discrete Faraday law
% using the corner field~\eqref{eq:Ezcorner},
% \begin{equation}
% \label{eq:faraday}
% \begin{aligned}
% (B_x)^{n+1}_{i+\frac12,j}&=(B_x)^{s}_{i+\frac12,j}
%  -\frac{\Delta t}{\Delta y}\Big[(E_z)_{i+\frac12,j+\frac12}-(E_z)_{i+\frac12,j-\frac12}\Big],\\
% (B_y)^{n+1}_{i,j+\frac12}&=(B_y)^{s}_{i,j+\frac12}
%  +\frac{\Delta t}{\Delta x}\Big[(E_z)_{i+\frac12,j+\frac12}-(E_z)_{i-\frac12,j+\frac12}\Big],
% \end{aligned}
% \end{equation}
% consistent with $\partial_t B_x=-\partial_y E_z$ and $\partial_t B_y=+\partial_x E_z$
% for $E_z=vB_x-uB_y$. Equation~\eqref{eq:faraday} is a discrete Stokes identity: the
% staggered divergence
% \begin{equation}
% \label{eq:divstag}
% [\nabla\!\cdot\!\Bb]_{i,j}
% =\frac{(B_x)_{i+\frac12,j}-(B_x)_{i-\frac12,j}}{\Delta x}
% +\frac{(B_y)_{i,j+\frac12}-(B_y)_{i,j-\frac12}}{\Delta y}
% \end{equation}
% is preserved exactly, at machine accuracy and at every time level, provided it
% vanishes initially, since the four corner contributions entering the two increments
% cancel in pairs. It is therefore essential that the staggered field be initialised
% compatibly with~\eqref{eq:divstag}, which is achieved by evaluating the initial
% condition for each component at the corresponding face centre at $t=0$. The
% staggered components are the primary representation of the magnetic field: they are
% stored and evolved for the entire simulation and are never reset from cell-centred
% data, since doing so would reintroduce divergence errors.
 
% \begin{remark}[Where the accuracy resides]
% \label{rem:ctaccuracy}
% The two-point difference in~\eqref{eq:faraday} is not an approximation: for the
% face-averaged $B_x$ it is the exact line integral of Faraday's law over the face.
% Raising its order would therefore buy no accuracy, and would destroy the cancellation
% in~\eqref{eq:divstag}, which holds only when the difference operator of the Faraday
% update and that of the divergence coincide. All of the accuracy of the CT step
% resides in the corner construction~\eqref{eq:Ezcorner}. Measuring the error of
% $\mathrm dB_x/\mathrm dt$ against exact corner values on a smooth manufactured
% electric field confirms this: the four-point interpolation~\eqref{eq:Ezcorner}
% converges at order $4.00$, whereas the gradient-corrected corner average of Gardiner
% and Stone~\cite{GardinerStone2005} converges at order $2.00$ and would cap the scheme
% at second order.
% \end{remark}
 
% \begin{remark}[Flux-consistency diagnostic]
% \label{rem:ctselfcheck}
% Because~\eqref{eq:emfdual} is assembled from the same face fluxes that the scheme
% uses, the construction admits a direct verification: differencing the collected
% fluxes with the operators listed in Remark~\ref{rem:ctoperators} must reproduce the
% cell-centred $\Bb$ that the predictor produced, up to the components of the implicit
% Laplacian that are not expressible as an electromotive force (the $x$-diffusion of
% $B_x$ and the $y$-diffusion of $B_y$). A residual of the order of the explicit share
% instead of that floor indicates a missing or mis-weighted term, and we use this as a
% standing check.
% \end{remark}
 
% Once the staggered components have been evolved through~\eqref{eq:faraday}, the
% cell-centred field is recovered by the fourth-order face-to-centre average
% \begin{equation}
% \label{eq:face2cell}
% (B_x)^{n+1}_{i,j}=\tfrac{1}{16}\Big(\!-(B_x)_{i-\frac32,j}+9(B_x)_{i-\frac12,j}
%  +9(B_x)_{i+\frac12,j}-(B_x)_{i+\frac32,j}\Big),
% \end{equation}
% and analogously for $B_y$, which overwrites the predictor; the two-point average
% of~\cite{GardinerStone2005} would cap the two-dimensional accuracy of the magnetic
% field at second order. In two space dimensions the out-of-plane component is left
% unchanged, $(B_z)^{n+1}_{i,j}=(\tilde B_z)^{n+1}_{i,j}$. As discussed by
% Rossmanith~\cite{Rossmanith2006} and T\'oth~\cite{Toth2000}, the cell-centred field
% so obtained satisfies, in addition to the face-based constraint~\eqref{eq:divstag},
% the corner-based discrete divergence-free condition
% \begin{equation}
% \label{eq:divtoth}
% [\nabla\!\cdot\!\Bb]_{i+\frac12,j+\frac12}
% =\frac{(B_x)_{i+1,j}+(B_x)_{i+1,j+1}-(B_x)_{i,j}-(B_x)_{i,j+1}}{2\Delta x}
% +\frac{(B_y)_{i,j+1}+(B_y)_{i+1,j+1}-(B_y)_{i,j}-(B_y)_{i+1,j}}{2\Delta y}=0 .
% \end{equation}
 
% \begin{remark}[The energy must be re-assembled after the transport step]
% \label{rem:ctenergy}
% Because Remark~\ref{rem:ctoperators} adopts the second option---a single two-point
% Stokes update, with the transport step \emph{replacing} rather than reproducing the
% predictor---the re-assembly of the total energy is not a matter of ordering but a
% structural requirement. The energy delivered by the implicit stage is built from the
% pressure returned by the solve of Section~\ref{sec:pressure} together with the
% predictor field $\tilde\Bb^{n+1}$; the transport step then overwrites $\Bb$
% through~\eqref{eq:face2cell}, and every correction $\delta\Bb$ shifts the internal
% energy by $-\half\,\delta(\|\Bb\|^{2})$ without the pressure being informed. In the
% low-$\beta$ interior of the blast-wave problem of Section~\ref{sec:result}, where
% $\half\|\Bb\|^{2}\sim10^{3}p$, this alone is enough to render the internal energy
% negative. We therefore re-assemble
% \begin{equation}
% \label{eq:energy_reassembly}
% U^{n+1}_{i,j}=\frac{p^{n+1}_{i,j}}{\gamma-1}
% +\frac{\|(\rho\mathbf{u})^{n+1}_{i,j}\|^{2}}{2\rho^{n+1}_{i,j}}
% +\half\big\|\Bb^{n+1}_{i,j}\big\|^{2}
% \end{equation}
% from the \emph{same} pressure and the post-transport divergence-free field, so that
% $\rho e=p^{n+1}/(\gamma-1)>0$ by construction. With the re-assembly omitted we
% observe internal energies of order $-10^{-2}$ within the first few steps, amplifying
% from stage to stage; with it in place the same computation runs to the final time
% with a strictly positive internal energy throughout. The same applies to the
% unstaggered variant below, which likewise replaces the predictor field.
% \end{remark}


% \subsection{Enforcement of the divergence-free constraint}
% \label{sec:CT}

% As anticipated in Section~\ref{sec:governing}, the solenoidal condition
% $\nabla\!\cdot\!\Bb=0$ is an involution of the ideal MHD system: if it holds for
% the initial data, the evolution preserves it for all later times at the continuous
% level. At the discrete level this property is not inherited automatically, and the
% divergence errors generated by a generic discretisation are known to produce
% spurious oscillations, incorrect shock speeds and unphysical
% states~\cite{BrackbillBarnes1980,Toth2000}---all of which are particularly damaging
% in the low-Mach and low-Alfv\'en regimes targeted here, where the magnetic pressure
% $\frac{1}{2}\|\Bb\|^{2}$ can exceed the gas pressure by orders of magnitude, so that
% a small error in $\Bb$ translates into a large spurious force. Several strategies
% have been devised to control or eliminate these errors. The eight-wave formulation
% of Powell et al.~\cite{Powell1999} augments the system with source terms
% proportional to the divergence error, at the price of losing strict conservation.
% In the projection method of Brackbill and Barnes~\cite{BrackbillBarnes1980} the
% magnetic field produced by the base scheme is projected onto the space of
% solenoidal fields through an auxiliary Poisson problem. Hyperbolic divergence
% cleaning~\cite{Dedner2002} couples the constraint to the conservation laws through
% a generalised Lagrange multiplier, which damps and advects the divergence error
% towards the boundaries. Among these alternatives, the constrained transport (CT)
% approach originally introduced by Evans and Hawley~\cite{EvansHawley1988} has
% emerged as one of the most widely adopted. In the present work we implement the
% unstaggered potential-based variant of CT following
% Rossmanith~\cite{Rossmanith2006}. The magnetic field is represented through a
% vector potential $\mathbf{A}$, such that $\Bb = \nabla\times\mathbf{A}$, which
% guarantees $\nabla\!\cdot\!\Bb=0$ identically at the discrete level. The induction
% equation is then replaced by an evolution equation for $\mathbf{A}$, discretised on
% the cell corners, so that no additional projection or cleaning step is required.
% Because the correction is applied at the end of the step, the density
% $\rho^{n+1}$ and the velocity $\mathbf{u}^{n+1}$ resulting from the pressure
% sub-system are already available, so that the transport step is driven by data at
% the new time level.


% \subsubsection{Unstaggered potential-based adaptation}
% \label{sec:CTunstag}
 
% The staggered formulation requires a dual representation of the magnetic field at
% faces and centres, which is cumbersome under adaptive mesh refinement and awkward to
% extend to multi-material settings~\cite{Rossmanith2006}. To avoid it we also
% implement a fully cell-centred, potential-based variant: in two dimensions only the
% $z$-component $A$ of the vector potential is tracked, stored at cell centres, and
% never overwritten by the predictor field. With $B_x=\partial_yA$ and
% $B_y=-\partial_xA$, the induction equation reduces to the scalar advection equation
% \begin{equation}
% \label{eq:Aadvect}
% \frac{\partial A}{\partial t}
% +u\,\frac{\partial A}{\partial x}
% +v\,\frac{\partial A}{\partial y}=0 ,
% \end{equation}
% which is nothing other than $\partial_tA=-E_z$ rewritten, since
% $E_z=vB_x-uB_y=u\,\partial_xA+v\,\partial_yA$.
 
% We advance~\eqref{eq:Aadvect} \emph{directly}, in advective form, with an upwind
% high-resolution method---the face values of $A$ are reconstructed by the same
% WENO3-Z procedure used in Section~\ref{sec:reconstruction} and selected by the sign
% of the face velocity---rather than by assembling an electric field and interpolating
% it to the cell centre. The distinction is not cosmetic. The face-to-corner-to-cell
% chain of the staggered construction is central at every stage, so its transfer
% function vanishes at the Nyquist frequency: no upwind dissipation would reach $A$,
% and its grid-scale modes would be undamped. Rossmanith advances the potential with a
% wave-propagation method for precisely this reason. In advective form the upwinding,
% and the limiter, act on $A$ itself. We stress that~\eqref{eq:Aadvect} is
% \emph{not} written in conservation form: $A$ is not a density, and the two forms
% differ by $A\,\nabla\!\cdot\!\mathbf{u}$, which does not vanish in compressible MHD.
 
% The magnetic field is then recovered by the fourth-order discrete curl
% \begin{equation}
% \label{eq:curl4}
% (B_x)^{n+1}_{i,j}=\frac{-A^{n+1}_{i,j+2}+8A^{n+1}_{i,j+1}-8A^{n+1}_{i,j-1}+A^{n+1}_{i,j-2}}{12\,\Delta y},
% \qquad
% (B_y)^{n+1}_{i,j}=-\frac{-A^{n+1}_{i+2,j}+8A^{n+1}_{i+1,j}-8A^{n+1}_{i-1,j}+A^{n+1}_{i-2,j}}{12\,\Delta x},
% \end{equation}
% which satisfies the \emph{matching} fourth-order central divergence identically,
% since the mixed fourth-order differences of $A$ commute and cancel exactly. As in the
% staggered case the divergence-free property is tied to the pair of operators:
% measuring~\eqref{eq:curl4} with a second-order divergence returns the
% $\mathcal O(\Delta x^{2})$ truncation difference between the two operators, which
% reflects the measuring stick rather than the field.
 
% If the initial data are given in terms of $\Bb$ alone, the initial potential is
% obtained by solving a Poisson problem as in~\cite{Rossmanith2006}. In either case the
% initial cell-centred field is re-initialised through~\eqref{eq:curl4} at $t=0$; were
% it left at its analytic value it would not satisfy the fourth-order divergence that
% the curl satisfies identically, and the mismatch would persist as a fixed
% $\mathcal O(\Delta x^{4})$ divergence floor for the whole computation. Boundary
% conditions for $A$ follow those of the magnetic field, with the caveat noted
% in~\cite{Rossmanith2006} that a zeroth-order extrapolation of $\Bb$ requires a
% first-order extrapolation of $A$.
 
% \begin{remark}[The unstaggered variant needs no flux-dual bookkeeping]
% \label{rem:unstag_advantage}
% Under the $\eps$-splitting this variant has a structural advantage over the
% staggered one. The construction~\eqref{eq:emfdual}, together with the
% share-dependent choice of difference operator discussed in
% Remark~\ref{rem:ctoperators}, is needed because the staggered Stokes update must be
% driven by an electromotive force consistent with the three separate numerical
% contributions by which $\Bb$ is advanced. Equation~\eqref{eq:Aadvect}, by contrast,
% is the \emph{physical} induction equation, and $\mathbf{u}^{n+1}$ is the velocity of the
% final state, obtained after the explicit advective update and both implicit
% sub-systems. All three shares are therefore carried automatically, without any of
% them having to be identified or weighted individually.
% \end{remark}






% ============================================================================
%  Sec.\ref{sec:highorder} -- rewritten following the structure of
%  \cite{BoscheriPareschi2021}: compact form -> high order in time -> high order
%  in space.  Corrections and additions relative to the previous draft:
%
%   FIXED   the third-order pair was quoted with delta = 1 - 1/sqrt(2) = 0.2929;
%           the implemented value is delta = 0.435866.  Both tableaux are now
%           given in full and match the implementation coefficient by
%           coefficient.
%   ADDED   the second-order pair actually implemented is the three-stage
%           SA SSP(3,3,2) of \cite{BoscheriPareschi2021}, not only LSDIRK2.
%   ADDED   the defect-correction solution of the fourth-order elliptic system,
%           absent from the draft although it is what makes the fourth-order
%           operator compatible with the positivity mechanism.
%   QUALIFIED  the spatial implicit discretisation is a staggered face family,
%           not the cell-centred finite differences of \cite{BoscheriPareschi2021}.
% ============================================================================

\subsection{High-order discretisation}
\label{sec:highorder}

The semi-discrete scheme of the preceding sections is made third-order accurate
in time by a semi-implicit IMEX Runge--Kutta
method~\cite{RKIMEX,BoscheriPareschi2021} and in space by the
reconstruction of Section~\ref{sec:reconstruction} for the explicit terms and by a
fourth-order finite volume operator family for the implicit ones.

\subsubsection{Time discretisation}
\label{sec:time}
Following~\cite{BoscheriPareschi2021}, the spatially discretised system is written
as $\partial_t\mathbf{U}=\mathcal H(\mathbf{U}_E,\mathbf{U}_I)$, where the first
argument is treated explicitly and the second implicitly. In the notation of
Sections~\ref{sec:hllem}--\ref{sec:implicit}, every quantity carrying the
superscript $n$ belongs to $\mathbf{U}_E$ and every quantity carrying $n+1$ to
$\mathbf{U}_I$: the explicit argument enters the advective flux, the explicit
shares $\epsM\FF^{p}$, $\epsB\FF^{B}$, the splitting parameters, and all the
coefficients and linearisation points of the implicit sub-systems (the enthalpy
$h^{n}$, the field $\Bb^{n}$ and the momentum $(\rho\mathbf{u})^{n}$ in the linearised
terms, the potential $A^{n}$ in the transport term), while the implicit argument
enters linearly. By construction $\mathcal H(\mathbf{U},\mathbf{U})$ is a
consistent discretisation of the flux divergence of the full system.

The system is integrated with an $s$-stage IMEX Runge--Kutta scheme with explicit
and implicit tableaux $(\tilde{\mathbf A},\tilde{\mathbf b})$ and
$(\mathbf A,\mathbf b)$, $\tilde{\mathbf b}=\mathbf b$. With
$\mathbf{U}_E^{1}=\mathbf{U}^{n}$,
\begin{align}
\label{eq:stage_E}
\mathbf{U}_E^{i}&=\mathbf{U}^{n}+\Delta t\sum_{j=1}^{i-1}\tilde a_{ij}\,\mathbf k_j,
\qquad
\tilde{\mathbf{U}}_I^{i}=\mathbf{U}^{n}+\Delta t\sum_{j=1}^{i-1}a_{ij}\,\mathbf k_j,\\
\label{eq:stage_k}
\mathbf k_i&=\mathcal H\big(\mathbf{U}_E^{i},\ \tilde{\mathbf{U}}_I^{i}+\Delta t\,a_{ii}\,\mathbf k_i\big),
\qquad 1\le i\le s,
\end{align}
and $\mathbf{U}^{n+1}=\mathbf{U}^{n}+\Delta t\sum_ib_i\mathbf k_i$. The implicit
relation~\eqref{eq:stage_k} is the step described in
Section~\ref{sec:implicit}, with $\Delta t\,a_{ii}$ in place of $\Delta t$,
$\tilde{\mathbf{U}}_I^{i}$ as the state to be updated and $\mathbf{U}_E^{i}$ as
the old time level: the magnetic predictor is solved for first and the pressure
second. The vector potential is advanced with the same tableaux. We use the
stiffly accurate third-order pair SA~DIRK$(3,4,3)$ of~\cite{BoscheriPareschi2021}, since stiff accuracy is required for consistency with the limit model as
$\epsM,\epsB\to0$,
\begin{equation}
\label{eq:sadirk343}
\begin{array}{c|cccc}
0        & 0        & 0         & 0         & 0\\
\delta   & \delta   & 0         & 0         & 0\\
c_3      & \tilde a_{31} & \tilde a_{32} & 0  & 0\\
1        & \tilde a_{41} & 1/2  & \tilde a_{43} & 0\\\hline
         & 0        & b_2  & b_3 & \delta
\end{array}
\qquad
\begin{array}{c|cccc}
\delta   & \delta & 0        & 0         & 0\\
\delta   & 0      & \delta   & 0         & 0\\
c_3      & 0      & c_3-\delta & \delta    & 0\\
1        & 0      & b_2 & b_3 & \delta\\\hline
         & 0      & b_2 & b_3 & \delta
\end{array},
\end{equation}
where $\delta\approx0.4358665215$ is the middle root of
$x^{3}-3x^{2}+\tfrac32x-\tfrac16=0$, $c_3=(1+\delta)/2$,
$b_2=-(6\delta^{2}-16\delta+1)/4$, $b_3=(6\delta^{2}-20\delta+5)/4$, and
$\tilde a_{31}\approx1.4377449738$, $\tilde a_{32}=c_3-\tilde a_{31}$,
$\tilde a_{41}\approx0.9169932984$, $\tilde a_{43}=\tfrac12-\tilde a_{41}$.



\subsubsection{Finite volume operators}
\label{sec:fv_ops}
All discrete unknowns are cell averages. Cell averages $\bar q$ and point values
$q$ at the cell centre are related, to fourth order, by
\begin{equation}
\label{eq:convdeconv}
q_{i,j}=\bar q_{i,j}-\tfrac1{24}\big(\delta_x^{2}+\delta_y^{2}\big)\bar q_{i,j},
\qquad
\bar q_{i,j}=q_{i,j}+\tfrac1{24}\big(\delta_x^{2}+\delta_y^{2}\big)q_{i,j},
\end{equation}
with $\delta_x^{2}q_{i,j}=q_{i+1,j}-2q_{i,j}+q_{i-1,j}$; the same relations with
the transverse second difference alone convert between face averages and
face-midpoint values. Every nonlinear function of the unknowns is evaluated on
point values and converted back. The implicit sub-systems are discretised with
the face average and the face-normal gradient of a cell-averaged quantity,
\begin{equation}
\label{eq:family4}
\mathcal I[\bar q]_{i+\frac12}=\tfrac{1}{12}\big(-\bar q_{i-1}+7\bar q_i+7\bar q_{i+1}-\bar q_{i+2}\big),
\qquad
\mathcal G[\bar q]_{i+\frac12}=\frac{\bar q_{i-1}-15\bar q_i+15\bar q_{i+1}-\bar q_{i+2}}{12\,\Delta x},
\end{equation}
both fourth-order accurate for the average over the face, and with the exact
divergence $\mathcal D[F]_i=(F_{i+\frac12}-F_{i-\frac12})/\Delta x$ of face-averaged
fluxes. The face average of a product is obtained from the product of the face
averages plus the transverse covariance,
\begin{equation}
\label{eq:covariance}
\overline{fg}=\bar f\,\bar g+\tfrac{\Delta y^{2}}{12}\,\partial_y\bar f\,\partial_y\bar g
+\mathcal O(\Delta y^{4})
\end{equation}
on an $x$-face. The elliptic term of~\eqref{eq:pwave} is the composition
\begin{equation}
\label{eq:composed4}
\mathcal D^{x}\Big[(1-\eps^{x}_{M})\Big((1-\eps^{x}_{M})\,
\overline{\mathcal I^{x}[h^{n}]\;\mathcal G^{x}[p]}
\;-\;\mathcal I^{x}[h^{n}]\;\mathcal I^{x}[p]\;\partial_x\eps^{x}_{M}\Big)\Big]+(x\to y),
\end{equation}
where $\partial_x\eps^{x}_{M}$ is the fourth-order central difference of the face values,
$(\partial_x\eps^{x}_{M})_{i+\frac12}=\big(-\eps^{x}_{M,i+\frac52}+8\eps^{x}_{M,i+\frac32}
-8\eps^{x}_{M,i-\frac12}+\eps^{x}_{M,i-\frac32}\big)/(12\Delta x)$.
with the splitting parameter taken at the faces; the same operators discretise
the enthalpy-flux divergence on the right-hand side, the momentum
update~\eqref{eq:p_mom}, the magnetic system~\eqref{eq:mag_sys} and the
terms~\eqref{eq:Bz},~\eqref{eq:Umag}.

The explicit terms are third-order accurate through the reconstruction of
Section~\ref{sec:reconstruction} at the stage state $\mathbf{U}_E^{i}$; the
numerical fluxes are evaluated at the face midpoints and converted to face
averages by~\eqref{eq:convdeconv}.
 
% \paragraph{Cell averages versus point values.}
% A finite volume scheme evolves cell averages, whereas the reconstruction returns
% face values and the Riemann solver returns a flux at a single point of the face.
% Every conversion between the two representations is only second-order accurate
% unless it is corrected, and each such conversion therefore caps the scheme at
% second order in two space dimensions, however accurate the reconstruction. Two
% appear in the present algorithm. The initial condition is sampled at the cell
% centre, whereas $\bar q_{i,j}$ must be the average of $q$ over the cell; and the
% differenced quantity in~\eqref{eq:adv_update} must be the face-averaged flux
% $\bar{\mathcal F}_{i+\frac12,j}
% =\Delta y^{-1}\!\int \mathcal F(x_{i+\frac12},y)\,\mathrm dy$, whereas the AP-HLLEM
% solver returns $\mathcal F$ at the face midpoint. Both are corrected by the
% convolution pair
% \begin{equation}
% \label{eq:convolution}
% \bar q = q+\tfrac{1}{24}\big(\delta^2_x+\delta^2_y\big)q+\mathcal O(h^{4}),
% \qquad
% q = \bar q-\tfrac{1}{24}\big(\delta^2_x+\delta^2_y\big)\bar q+\mathcal O(h^{4}),
% \end{equation}
% with $\delta^2$ the undivided second difference, applied in both directions to the
% initial data and in the transverse direction alone to the face fluxes. In one space
% dimension the transverse difference vanishes identically, so the face correction is
% inactive and the scheme attains its design order without it; this is why a
% one-dimensional convergence study does not reveal the effect. Because $\delta^2$ is
% $\mathcal O(1)$ at a discontinuity, both corrections are applied only to smooth
% data and are disabled in the shock-dominated tests of
% Section~\ref{sec:result}, where the formal order is in any case one.


 
\subsection{Fully discrete scheme in two dimensions}
\label{sec:fully_discrete}

We use the operators of Section~\ref{sec:fv_ops}: the face average
$\mathcal I^{x}$, the face-normal gradient $\mathcal G^{x}$ and the divergence
$\mathcal D^{x}[F]_{i,j}=(F_{i+\frac12,j}-F_{i-\frac12,j})/\Delta x$, with their
$y$-analogues, and write $\overline{f\,g}$ for the face average of a
product~\eqref{eq:covariance}. One stage of the IMEX Runge--Kutta scheme of
Section~\ref{sec:time}, with $\Delta t$ standing for the stage time step
$\Delta t\,a_{ii}$, the superscript $n$ for the explicit stage state
$\mathbf{U}_E^{i}$ and a star for the state obtained from the implicit base
state $\tilde{\mathbf{U}}_I^{i}$ by the explicit update, proceeds as follows.

\begin{enumerate}[leftmargin=2.2em]

\item \textbf{Splitting parameters.}
Compute $\epsM,\epsB$ from~\eqref{eq:eps_def} in every cell and the interface
values $\eps^{x}_{i+\frac12,j}$, $\eps^{y}_{i,j+\frac12}$ from~\eqref{eq:eps_face}.

\item \textbf{Explicit update.}
With the face states of Section~\ref{sec:reconstruction} and the AP-HLLEM
fluxes~\eqref{eq:hllem}, converted to face averages,
\begin{equation}\label{eq:step_adv}
\mathbf{U}^{*}_{i,j}=\tilde{\mathbf{U}}_{I,i,j}
-\Delta t\,\big(\mathcal D^{x}[\tilde\FF^{\eps}]+\mathcal D^{y}[\tilde\GG^{\eps}]\big)_{i,j},
\end{equation}
which gives the density $\rho^{n+1}=\rho^{*}$.

\item \textbf{Magnetic predictor.}
Solve the linear system~\eqref{eq:mag_sys} for $\tilde\Bb$, with the
stress divergence of~\eqref{eq:mag_mom} and the induction operator
of~\eqref{eq:induction} discretised as $\mathcal D[(1-\epsB)\,\mathcal I[\,\cdot\,]]$
and with the pressure gradient of~\eqref{eq:mag_mom} as in~\eqref{eq:step_mom}
below.

\item \textbf{Lorentz force.}
\begin{equation}\label{eq:step_mag_mom}
(\rho\mathbf{u})^{\ast\ast}_{i,j}=(\rho\mathbf{u})^{\ast}_{i,j}
-\Delta t\,\Big(\mathcal D^{x}\big[(1-\eps^{x}_{B})\,\mathcal I^{x}[\FF^{B}_{\rho\mathbf{u}}(\tilde\Bb)]\big]
+\mathcal D^{y}\big[(1-\eps^{y}_{B})\,\mathcal I^{y}[\GG^{B}_{\rho\mathbf{u}}(\tilde\Bb)]\big]\Big)_{i,j},
\qquad
\FF^{B}_{\rho\mathbf{u}}(\Bb)=\half\|\Bb\|^{2}\mathbf e_x-B_x\Bb .
\end{equation}

\item \textbf{Pressure solve.}
For a trial pressure $p$, the momentum is
\begin{equation}\label{eq:step_mom}
(\rho u)_{i,j}=(\rho u)^{\ast\ast}_{i,j}
-\Delta t\,\mathcal D^{x}\big[(1-\eps^{x}_{M})\,\mathcal I^{x}[p]\big]_{i,j},
\qquad
(\rho v)_{i,j}=(\rho v)^{\ast\ast}_{i,j}
-\Delta t\,\mathcal D^{y}\big[(1-\eps^{y}_{M})\,\mathcal I^{y}[p]\big]_{i,j},
\end{equation}
the in-plane field follows from~\eqref{eq:Aupdate} and~\eqref{eq:curl4}, and
\begin{align}
\label{eq:step_Bz}
B_{z,i,j}&=B^{\ast}_{z,i,j}
-\Delta t\,\mathcal D^{x}\big[(1-\eps^{x}_{B})\,\mathcal I^{x}[u\,B_z^{n}-w\,B_x^{n}]\big]_{i,j}
-\Delta t\,\mathcal D^{y}\big[(1-\eps^{y}_{B})\,\mathcal I^{y}[v\,B_z^{n}-w\,B_y^{n}]\big]_{i,j},\\
\label{eq:step_Umag}
E^{\ast\ast}_{i,j}&=E^{\ast}_{i,j}
-\Delta t\,\mathcal D^{x}\big[(1-\eps^{x}_{B})\,\mathcal I^{x}[\|\Bb^{n}\|^{2}u-(\mathbf{u}\!\cdot\!\Bb^{n})B_x^{n}]\big]_{i,j}
-\Delta t\,\mathcal D^{y}\big[(1-\eps^{y}_{B})\,\mathcal I^{y}[\|\Bb^{n}\|^{2}v-(\mathbf{u}\!\cdot\!\Bb^{n})B_y^{n}]\big]_{i,j},
\end{align}
with $\mathbf{u}=(\rho\mathbf{u})/\rho^{n+1}$ and all products formed on point values. The
pressure $p^{n+1}$ is the solution of
\begin{equation}\label{eq:step_p1}
\frac{p_{i,j}}{\gamma-1}+k_{i,j} + m_{i,j}
+\Delta t\,\Big(\mathcal D^{x}\big[(1-\eps^{x}_{M})\,\overline{\mathcal I^{x}[h^{n}]\;(\rho u)_{f}}\big]
+\mathcal D^{y}\big[(1-\eps^{y}_{M})\,\overline{\mathcal I^{y}[h^{n}]\;(\rho v)_{f}}\big]\Big)_{i,j}
=E^{\ast\ast}_{i,j},
\end{equation}
where $k$ and $m$ are the cell average of~\eqref{eq:klin} and the face momentum is
\begin{equation}\label{eq:step_mf}
(\rho u)_{f}=\mathcal I^{x}[(\rho u)^{\ast\ast}]
-\Delta t\,\Big((1-\eps^{x}_{M})\,\mathcal G^{x}[p]-\mathcal I^{x}[p]\,\partial_x\eps^{x}_{M}\Big),
\end{equation}
that is, the face average of the momentum~\eqref{eq:step_mom} with the pressure
gradient evaluated at the face. Equation~\eqref{eq:step_p1} is linear in $p$ and
is solved by the defect correction~\eqref{eq:defect}.

\item \textbf{Update.}
With the converged pressure, $(\rho\mathbf{u})^{n+1}$ is given by~\eqref{eq:step_mom},
$\Bb^{n+1}$ by~\eqref{eq:Aupdate}, \eqref{eq:curl4} and~\eqref{eq:step_Bz}, and the
total energy by
\begin{equation}
\label{eq:step_energy}
E^{n+1}_{i,j}=\frac{p^{n+1}_{i,j}}{\gamma-1}
+\frac{(\rho\mathbf{u})^{n}_{i,j}\!\cdot\!(\rho\mathbf{u})^{n+1}_{i,j}-\half\|(\rho\mathbf{u})^{n}_{i,j}\|^{2}}{\rho^{n+1}_{i,j}}
+\Bb^{n}_{i,j}\!\cdot\!\Bb^{n+1}_{i,j}-\half\|\Bb^{n}_{i,j}\|^{2},
\end{equation}
that is, with the kinetic and magnetic energies linearised as
in~\eqref{eq:klin}. This is the energy for which the pressure
system~\eqref{eq:energy_imp} was derived, so that the update is conservative.

\end{enumerate}

Each stage thus solves for the magnetic predictor and then for the pressure,
with no iteration between the sub-systems. The scheme is
third-order accurate in space and time, the magnetic field is divergence-free
with respect to the fourth-order central divergence, and mass, momentum and
total energy are conserved, the latter exactly wherever the pressure floor is
inactive.




 
\section{Numerical Results}\label{sec:result}

In this section the proposed scheme is validated on a set of standard test
problems. We first measure the experimental order of convergence on a smooth
solution, decreasing the background density over five orders of magnitude so
that the magnetic sub-system becomes progressively stiffer. One-dimensional
Riemann problems then exercise the HLLEM flux and the characteristic
reconstruction on the individual wave families. The remaining tests are
two-dimensional: a stationary vortex, three benchmarks with shocks, and two
problems dominated by the advection and the shear of the magnetic field.

Unless otherwise stated, the ideal gas equation of state is used with
$\gamma=5/3$, and the time step is given by~\eqref{eq:cfl} with
$\mathrm{CFL}=0.8$.
In two space dimensions the magnetic field is advanced by the unstaggered
constrained transport of Section~\ref{sec:CT},
and the discrete divergence remains at the level of round-off in all tests; in
one space dimension $B_x$ is constant and constrained transport is not needed.

\subsection{Accuracy of the numerical scheme}
\label{sec:accuracy}
The order of accuracy of the scheme is verified on the smooth magnetised
vortex of~\cite{accuracyvortex, BoscheriThomann2024}, set on the periodic-free
square $\Omega=[-5,5]^2$. The initial condition consists of a homogeneous state
$\mathbf{q}_0=(\rho_0,\mathbf{u}_0,p_0,\mathbf{0})$ to which a localised,
rotationally symmetric perturbation $\delta\mathbf{q}$ is added; the density is
left unperturbed. Let $r=\sqrt{x^2+y^2}$ and $\theta$ denote polar coordinates
centred at the origin, and introduce the Gaussian profile
%
\begin{equation}
\varphi(r)=\exp\!\left(\frac{1-r^{2}}{2}\right).
\label{eq:vortex_profile}
\end{equation}
%
Both the velocity and the magnetic field are purely azimuthal, with
%
\begin{equation}
\delta v_\theta=\frac{\tilde{v}}{2\pi}\,r\,\varphi(r),
\qquad
\delta B_\theta=\frac{\tilde{B}}{2\pi}\,r\,\varphi(r),
\label{eq:vortex_theta}
\end{equation}
%
so that in Cartesian components
%
\begin{align}
(\delta u,\,\delta v)&=\delta v_\theta\,(-\sin\theta,\ \cos\theta),
\label{eq:vortex_vel}\\[2pt]
(\delta B_x,\,\delta B_y)&=\delta B_\theta\,(-\sin\theta,\ \cos\theta).
\label{eq:vortex_mag}
\end{align}
%
The amplitudes are fixed to $\tilde{v}=\sqrt{2\pi}$ and $\tilde{B}=1$. The
magnetic field derives from the vector potential
%
\begin{equation}
\delta A_z=\frac{\tilde{B}}{2\pi}\,\varphi(r),
\label{eq:vortex_Az}
\end{equation}
%
and is therefore solenoidal by construction. A stationary configuration in the
co-moving frame requires the radial pressure gradient to compensate the
centrifugal force together with the Lorentz force, the latter split into
magnetic tension and magnetic pressure:
%
\begin{equation}
\frac{\mathrm{d}\,\delta p}{\mathrm{d}r}
 =\frac{\rho_0\,\delta v_\theta^{2}}{r}
 -\frac{\delta B_\theta^{2}}{r}
 -\frac{1}{2}\,\frac{\mathrm{d}\bigl(\delta B_\theta^{2}\bigr)}{\mathrm{d}r}.
\label{eq:vortex_equil}
\end{equation}
%
Requiring $\delta p\to 0$ as $r\to\infty$, the integration
of~\eqref{eq:vortex_equil} can be carried out exactly and gives
%
\begin{equation}
\delta p=\frac{\varphi(r)^{2}}{2}
 \left[\left(\frac{\tilde{B}}{2\pi}\right)^{\!2}\bigl(1-r^{2}\bigr)
 -\rho_0\left(\frac{\tilde{v}}{2\pi}\right)^{\!2}\right].
\label{eq:vortex_dp}
\end{equation}
%
We choose $p_0=1$ and $\mathbf{u}_0=(1,1,0)$, so that the structure is
transported along the diagonal of the domain without change of shape.

The time step of the present method is restricted by a CFL condition based on
the material velocity alone, with no dependence on the acoustic or Alfv\'en
speeds. The convergence study is therefore repeated over a range of
magnetosonic regimes. Superimposing a uniform magnetic field would destroy the
exact solution, so the regime is controlled through the background density
instead: because the Alfv\'en speed scales as $\rho^{-1/2}$, reducing $\rho_0$
makes the Alfv\'en waves faster and the flow increasingly sub-Alfv\'enic. The
same reduction raises the sound speed and weakens the kinetic term
in~\eqref{eq:vortex_dp}, while the magnetic term is unaffected. Since the
unperturbed state carries no magnetic field, a global Alfv\'en Mach number
cannot be assigned; we characterise each case by the maximum acoustic Mach
number $M_c$ and by the largest Alfv\'en speed $c_a$, listed in
Table~\ref{tab:vortex_regimes} for $\rho_0=10^{-k}$, $k=0,\dots,5$, together
with the smallest Alfv\'en Mach number attained where the field does not
vanish. The sweep crosses both sonic lines: the flow is supersonic at
$\rho_0=1$ and sub-Alfv\'enic from $\rho_0=10^{-2}$ on, so that the splitting
is exercised from the fully explicit regime down to one in which both channels
are stiff. It does not, however, reach a low-$\beta$ state: the magnetic term
of~\eqref{eq:vortex_dp} and $\|\Bb\|$ are both independent of $\rho_0$, so
the plasma beta stays near $79$ throughout, and the regimes of this test are
the high-$\beta$ ones.

\begin{table}[pos=htbp]
\centering
\footnotesize
\caption{MHD traveling vortex: the magnetosonic regime of each background
density. $M_c=\max_\Omega\|\mathbf u\|/a$ and $c_a=\max_\Omega\|\Bb\|/\sqrt{\rho}$;
$M_A$ is the smallest value of $\|\mathbf u\|/c_a$ over the cells in which
$\Bb\neq\mathbf 0$, and $\beta$ the smallest plasma beta attained there. Since
$\|\Bb\|$ does not depend on $\rho_0$, the speed $c_a$ follows
$\rho_0^{-1/2}$ exactly; $M_c$ follows $\rho_0^{1/2}$ to within $0.4\%$ from
$\rho_0=10^{-1}$ on and to $4\%$ at $\rho_0=1$, the residual being the
variation of the pressure minimum through the kinetic term
of~\eqref{eq:vortex_dp}. The last two columns give the ratio of the admissible
step to that of a fully explicit discretisation,
$\max_\Omega|\lambda_{\max}|/\max_\Omega|\tilde\lambda^{c}_{\max}|$, evaluated on
the initial data for the two-parameter splitting and for the single-parameter
choice $\epsB\equiv0$.}
\label{tab:vortex_regimes}
\begin{tabular}{lrrrrrr}
\toprule
$\rho_0$ & $M_c$ & $c_a$ & $\min M_A$ & $\min\beta$ &
\multicolumn{2}{c}{$\Delta t/\Delta t_{\mathrm{MHD}}$} \\
\cmidrule(l){6-7}
 & & & & & two parameters & $\epsB\equiv0$ \\
\midrule
$10^{0}$  & $1.4732$ & $0.1592$ & $6.379$ & $72.5$ & $1.00$  & $1.00$ \\
$10^{-1}$ & $0.4459$ & $0.5033$ & $2.017$ & $78.3$ & $1.00$  & $1.00$ \\
$10^{-2}$ & $0.1405$ & $1.5915$ & $0.638$ & $78.9$ & $4.79$  & $10.2$ \\
$10^{-3}$ & $0.0444$ & $5.0329$ & $0.202$ & $78.9$ & $7.10$  & $30.5$ \\
$10^{-4}$ & $0.0141$ & $15.9155$ & $0.064$ & $78.9$ & $26.2$  & $94.6$ \\
$10^{-5}$ & $0.0044$ & $50.3292$ & $0.020$ & $79.0$ & $86.5$  & $298$ \\
\bottomrule
\end{tabular}
\end{table}

All simulations are run to $t_f=1$ on a sequence of uniformly refined meshes.
The $L^2$ errors of $\rho$, $u$, $p$, $B_x$ and $A_z$ and the corresponding
experimental orders of convergence (EOC) are collected in
Table~\ref{tab:vortex_eoc_o3}. The design third order is attained for all
variables and across all regimes considered, approached from below on the
finest grids at $\rho_0=1$ and $\rho_0=10^{-1}$ and exceeded at the smallest
background densities, where the coarsest grids do not yet resolve the
perturbation. The last two columns of Table~\ref{tab:vortex_regimes} show what the splitting
buys here, and what the second parameter costs. The admissible step exceeds
that of a fully explicit discretisation---bounded by the fastest characteristic
speed---by a factor rising to $86.5$ at $\rho_0=10^{-5}$, with the design order
retained throughout; at $\rho_0\ge10^{-1}$ the flow is still transonic and
there is nothing to gain. The gain is smaller than the four orders of magnitude
reached in the Kelvin--Helmholtz test of Section~\ref{sec:KH}
(Fig.~\ref{fig:kh_dtratio}) because this test is high-$\beta$ throughout:
$\epsB$ saturates in the annulus where the Alfv\'en speed is already a few
times the flow speed, and that annulus, not the vortex core, sets
$\max_\Omega|\tilde\lambda^{c}_{\max}|$. Pinning $\epsB$ to zero removes it and
multiplies the admissible step by between $2.1$ and $4.3$, which is the price,
in this test, of returning the magnetic flux to the upwind solver; it is the
high-$\beta$ side of the trade-off discussed in Section~\ref{sec:intro}, and by
the argument given there it is the side on which the second parameter is not
expected to pay for itself.


\begin{table}[pos=htbp]
\centering
\footnotesize
\setlength{\tabcolsep}{4pt}
\caption{MHD traveling vortex: $L^2$ errors and EOC of the proposed third-order
scheme for varying background density $\rho_0$ at $t_f=1$, computed with CFL $=0.2$.
All computations are carried out on the domain $[-5,5]^2$ with $N$ cells per
direction.}
\label{tab:vortex_eoc_o3}
\begin{tabular}{ll*{5}{rc}}
\toprule
$\rho_0$ & $N$ & \multicolumn{2}{c}{$\rho$} & \multicolumn{2}{c}{$u$} & \multicolumn{2}{c}{$p$} & \multicolumn{2}{c}{$B_x$} & \multicolumn{2}{c}{$A_z$} \\
 & & \multicolumn{1}{c}{$L^2$} & EOC & \multicolumn{1}{c}{$L^2$} & EOC & \multicolumn{1}{c}{$L^2$} & EOC & \multicolumn{1}{c}{$L^2$} & EOC & \multicolumn{1}{c}{$L^2$} & EOC \\
\midrule
$10^{0}$  & 32   & 3.06e-03 & --   & 1.98e-02 & --   & 1.01e-02 & --   & 1.68e-03 & --   & 7.94e-04 & --   \\
          & 64   & 4.67e-04 & 2.71 & 2.72e-03 & 2.86 & 1.40e-03 & 2.85 & 1.26e-04 & 3.73 & 5.43e-05 & 3.87 \\
          & 128  & 6.23e-05 & 2.91 & 3.46e-04 & 2.97 & 1.78e-04 & 2.98 & 1.46e-05 & 3.11 & 6.62e-06 & 3.04 \\
          & 256  & 7.92e-06 & 2.98 & 4.35e-05 & 2.99 & 2.23e-05 & 3.00 & 2.03e-06 & 2.85 & 8.34e-07 & 2.99 \\
$10^{-1}$ & 32   & 2.22e-04 & --   & 4.63e-02 & --   & 2.59e-03 & --   & 2.36e-03 & --   & 1.09e-03 & --   \\
          & 64   & 3.41e-05 & 2.70 & 6.74e-03 & 2.78 & 4.12e-04 & 2.65 & 3.15e-04 & 2.91 & 1.25e-04 & 3.14 \\
          & 128  & 4.45e-06 & 2.94 & 8.66e-04 & 2.96 & 5.28e-05 & 2.96 & 4.17e-05 & 2.91 & 1.55e-05 & 3.01 \\
          & 256  & 5.61e-07 & 2.99 & 1.09e-04 & 3.00 & 6.60e-06 & 3.00 & 5.37e-06 & 2.96 & 1.94e-06 & 3.00 \\
$10^{-2}$ & 32   & 9.94e-05 & --   & 2.26e-02 & --   & 4.30e-04 & --   & 2.91e-03 & --   & 1.44e-03 & --   \\
          & 64   & 8.98e-06 & 3.47 & 2.89e-03 & 2.97 & 5.58e-05 & 2.94 & 3.38e-04 & 3.10 & 1.56e-04 & 3.21 \\
          & 128  & 9.75e-07 & 3.20 & 3.67e-04 & 2.98 & 6.81e-06 & 3.04 & 4.22e-05 & 3.00 & 1.91e-05 & 3.03 \\
          & 256  & 1.20e-07 & 3.03 & 4.61e-05 & 2.99 & 8.40e-07 & 3.02 & 5.37e-06 & 2.97 & 2.38e-06 & 3.01 \\
$10^{-3}$ & 32   & 2.14e-05 & --   & 2.50e-02 & --   & 3.53e-04 & --   & 2.71e-03 & --   & 1.38e-03 & --   \\
          & 64   & 1.40e-06 & 3.94 & 4.11e-03 & 2.61 & 3.26e-05 & 3.43 & 2.32e-04 & 3.55 & 1.10e-04 & 3.65 \\
          & 128  & 8.38e-08 & 4.06 & 5.96e-04 & 2.79 & 3.63e-06 & 3.17 & 2.50e-05 & 3.21 & 1.26e-05 & 3.13 \\
          & 256  & 8.47e-09 & 3.31 & 7.77e-05 & 2.94 & 4.32e-07 & 3.07 & 3.15e-06 & 2.99 & 1.55e-06 & 3.02 \\
$10^{-4}$ & 32   & 9.13e-06 & --   & 5.90e-02 & --   & 1.29e-03 & --   & 7.20e-03 & --   & 3.72e-03 & --   \\
          & 64   & 5.64e-07 & 4.02 & 5.24e-03 & 3.49 & 9.06e-05 & 3.83 & 4.65e-04 & 3.95 & 1.86e-04 & 4.32 \\
          & 128  & 2.70e-08 & 4.38 & 7.46e-04 & 2.81 & 4.30e-06 & 4.40 & 2.21e-05 & 4.40 & 8.93e-06 & 4.38 \\
          & 256  & 1.76e-09 & 3.94 & 1.05e-04 & 2.83 & 1.95e-07 & 4.46 & 1.45e-06 & 3.93 & 6.32e-07 & 3.82 \\
$10^{-5}$ & 32   & 5.27e-06 & --   & 5.37e-01 & --   & 5.96e-03 & --   & 3.43e-02 & --   & 2.26e-02 & --   \\
          & 64   & 5.34e-07 & 3.30 & 2.15e-02 & 4.64 & 5.94e-04 & 3.33 & 2.93e-03 & 3.55 & 1.36e-03 & 4.06 \\
          & 128  & 1.81e-08 & 4.88 & 1.50e-03 & 3.84 & 2.75e-05 & 4.43 & 1.36e-04 & 4.43 & 5.25e-05 & 4.69 \\
          & 256  & 7.53e-10 & 4.59 & 1.42e-04 & 3.39 & 1.13e-06 & 4.60 & 5.63e-06 & 4.60 & 2.32e-06 & 4.50 \\
\bottomrule
\end{tabular}
\end{table}

\subsection{One-dimension MHD Riemann Tests}
\label{sec:riemann1d}
\label{sec:MHD_1d}

Having established the convergence order in Section~\ref{sec:accuracy}, we
now turn to the shock-capturing behaviour of the scheme, which is examined on a
set of classical one-dimensional Riemann problems for the ideal MHD equations.
These problems give rise to fast and slow shocks and rarefactions, contact and
rotational discontinuities and, in RP1, a compound wave. Each problem is
initialised with two uniform
states $\mathbf{W}_L=[\rho_L,\mathbf{u}_L,p_L,\mathbf{B}_L]$ and
$\mathbf{W}_R=[\rho_R,\mathbf{u}_R,p_R,\mathbf{B}_R]$ that meet at $x=x_0$.
The seven configurations RP1--RP7, with the location of the initial jump and the
output time, are listed in Table~\ref{tab:mhd_riemann}. RP1 is the shock tube
of Brio and Wu~\cite{BrioWu1988}. RP2 is the Ryu--Jones
problem~\cite{RyuJones1995}, whose solution contains all seven MHD waves. RP3
and RP4 were adopted in~\cite{Dumbser2019},~\cite{Fambri2021} and~\cite{Dematte2024} to
validate semi-implicit schemes. RP5 and RP6 are head-on collisions of
hypersonic streams, with acoustic Mach numbers of about 12 and 20 in RP5 and
29 in RP6, and RP7 is an isolated rotational discontinuity, which is
stationary since $u+c_a=0$. % CHECK: cite the sources of RP5--RP7
All computations are carried out on the
interval $[-0.5,0.5]$ using $2000$ cells and transmissive boundary conditions.
Since $B_x$ is constant in one space dimension, constrained transport is not
applied and the transverse field is updated in flux form.

\begin{table}[pos=htbp]
\centering
\caption{Initial left (L) and right (R) states for the MHD Riemann test
problems. The discontinuity is located at $x_0$. }
\label{tab:mhd_riemann}
\setlength{\tabcolsep}{3.5pt}
\renewcommand{\arraystretch}{1.15}
\resizebox{\textwidth}{!}{%
\begin{tabular}{llllllllllll}   % 12 列: ll + 8数据 + p_gas + x0 + t_out
\toprule
Test & & $\rho$ & $u$ & $v$ & $w$ & $B_x$ & $B_y$ & $B_z$ & $p$ & $x_0$ & $t_{\mathrm{out}}$ \\
\midrule
\multirow{2}{*}{RP1}
 & L: & $1.0$   & $0.0$ & $0.0$ & $0.0$ & $0.75$ & $1.0$  & $0.0$ & $1.0$ & \multirow{2}{*}{$0.0$} & \multirow{2}{*}{$0.1$} \\
 & R: & $0.125$ & $0.0$ & $0.0$ & $0.0$ & $0.75$ & $-1.0$ & $0.0$ & $0.1$ & & \\
\addlinespace
\multirow{2}{*}{RP2}
 & L: & $1.08$   & $1.2$     & $0.01$   & $0.5$      & $0.564190$ & $1.015541$ & $0.564190$ & $0.95$    & \multirow{2}{*}{$-0.1$} & \multirow{2}{*}{$0.2$} \\
 & R: & $0.9891$ & $-0.0131$ & $0.0269$ & $0.010037$ & $0.564190$ & $1.135262$ & $0.564923$ & $0.97159$ & & \\
\addlinespace
\multirow{2}{*}{RP3}
 & L: & $1.7$ & $0.0$ & $0.0$ & $0.0$       & $1.1$ & $1.0$      & $0.0$      & $1.7$ & \multirow{2}{*}{$-0.1$} & \multirow{2}{*}{$0.15$} \\
 & R: & $0.2$ & $0.0$ & $0.0$ & $-1.496891$ & $1.1$ & $0.785887$ & $0.618370$ & $0.2$ & & \\
\addlinespace
\multirow{2}{*}{RP4}
 & L: & $1.0$ & $0.0$ & $0.0$ & $0.0$ & $1.3$ & $1.0$  & $0.0$ & $1.0$ & \multirow{2}{*}{$0.0$} & \multirow{2}{*}{$0.16$} \\
 & R: & $0.4$ & $0.0$ & $0.0$ & $0.0$ & $1.3$ & $-1.0$ & $0.0$ & $0.4$ & & \\
\multirow{2}{*}{RP5}
 & L: & $0.15$ & $21.55$  & $1.0$ & $1.0$ & $0.05/\sqrt{4\pi}$ & $-2.0/\sqrt{4\pi}$ & $-1.0/\sqrt{4\pi}$ & $0.28$ & \multirow{2}{*}{$0.0$} & \multirow{2}{*}{$0.04$} \\
 & R: & $0.1$  & $-26.45$ & $0.0$ & $0.0$ & $0.05/\sqrt{4\pi}$ & $2.0/\sqrt{4\pi}$  & $1.0/\sqrt{4\pi}$  & $0.1$  & & \\
\addlinespace
\multirow{2}{*}{RP6}
 & L: & $1.0$ & $36.87$  & $-0.115$ & $-0.0386$ & $4/\sqrt{4\pi}$ & $4/\sqrt{4\pi}$ & $1/\sqrt{4\pi}$ & $1.0$ & \multirow{2}{*}{$0.0$} & \multirow{2}{*}{$0.03$} \\
 & R: & $1.0$ & $-36.87$ & $0.0$    & $0.0$     & $4/\sqrt{4\pi}$ & $4/\sqrt{4\pi}$ & $1/\sqrt{4\pi}$ & $1.0$ & & \\
\addlinespace
\multirow{2}{*}{RP7}
 & L: & $1/(4\pi)$ & $-1.0$ & $1.0$  & $-1.0$ & $1/\sqrt{4\pi}$ & $-1/\sqrt{4\pi}$ & $1/\sqrt{4\pi}$ & $1.0$ & \multirow{2}{*}{$0.0$} & \multirow{2}{*}{$0.25$} \\
 & R: & $1/(4\pi)$ & $-1.0$ & $-1.0$ & $-1.0$ & $1/\sqrt{4\pi}$ & $1/\sqrt{4\pi}$  & $1/\sqrt{4\pi}$ & $1.0$ & & \\
\bottomrule
\end{tabular}}

\vspace{2pt}
\end{table}

\begin{figure}[pos=htbp]
    \centering
    \begin{minipage}{0.32\linewidth}
        \includegraphics[width=\linewidth]{figs/RP1_rho.png}
    \end{minipage}
    \hspace{0.01\linewidth}%
    \begin{minipage}{0.32\linewidth}
        \includegraphics[width=\linewidth]{figs/RP1_p.png}
    \end{minipage}
    \hspace{0.01\linewidth}%
    \begin{minipage}{0.32\linewidth}
        \includegraphics[width=\linewidth]{figs/RP1_By.png}
    \end{minipage}
    \vspace{0.5em}

    \begin{minipage}{0.32\linewidth}
        \includegraphics[width=\linewidth]{figs/RP2_rho.png}
    \end{minipage}
    \hspace{0.01\linewidth}%
    \begin{minipage}{0.32\linewidth}
        \includegraphics[width=\linewidth]{figs/RP2_p.png}
    \end{minipage}
    \hspace{0.01\linewidth}%
    \begin{minipage}{0.32\linewidth}
        \includegraphics[width=\linewidth]{figs/RP2_By.png}
    \end{minipage}
    \vspace{0.5em}
    \begin{minipage}{0.32\linewidth}
        \includegraphics[width=\linewidth]{figs/RP3_rho.png}
    \end{minipage}
    \hspace{0.01\linewidth}%
    \begin{minipage}{0.32\linewidth}
        \includegraphics[width=\linewidth]{figs/RP3_p.png}
    \end{minipage}
    \hspace{0.01\linewidth}%
    \begin{minipage}{0.32\linewidth}
        \includegraphics[width=\linewidth]{figs/RP3_By.png}
    \end{minipage}
    
    \vspace{0.5em}

    \begin{minipage}{0.32\linewidth}
        \includegraphics[width=\linewidth]{figs/RP4_rho.png}
    \end{minipage}
    \hspace{0.01\linewidth}%
    \begin{minipage}{0.32\linewidth}
        \includegraphics[width=\linewidth]{figs/RP4_p.png}
    \end{minipage}
    \hspace{0.01\linewidth}%
    \begin{minipage}{0.32\linewidth}
        \includegraphics[width=\linewidth]{figs/RP4_By.png}
    \end{minipage}


    \caption{Numerical results for the Riemann problems RP1, RP2, RP3 and RP4 (from top to bottom) obtained using the new semi-implicit scheme compared to exact Riemann solutions (black line). The black symbols indicate results obtained by proposed IMEX scheme.
    The fluid density (left), the fluid pressure (middle) and the component of the magnetic field (right) are depicted at the chosen output time for each test.}
    \label{fig:MHDRPs1}
\end{figure}
\clearpage

\begin{figure}[pos=htbp]
    \centering
    \begin{minipage}{0.32\linewidth}
        \includegraphics[width=\linewidth]{figs/RP5_rho.png}
    \end{minipage}
    \hspace{0.01\linewidth}%
    \begin{minipage}{0.32\linewidth}
        \includegraphics[width=\linewidth]{figs/RP5_p.png}
    \end{minipage}
    \hspace{0.01\linewidth}%
    \begin{minipage}{0.32\linewidth}
        \includegraphics[width=\linewidth]{figs/RP5_By.png}
    \end{minipage}
    \vspace{0.5em}

    \begin{minipage}{0.32\linewidth}
        \includegraphics[width=\linewidth]{figs/RP6_rho.png}
    \end{minipage}
    \hspace{0.01\linewidth}%
    \begin{minipage}{0.32\linewidth}
        \includegraphics[width=\linewidth]{figs/RP6_p.png}
    \end{minipage}
    \hspace{0.01\linewidth}%
    \begin{minipage}{0.32\linewidth}
        \includegraphics[width=\linewidth]{figs/RP6_By.png}
    \end{minipage}
    \vspace{0.5em}
    \begin{minipage}{0.32\linewidth}
        \includegraphics[width=\linewidth]{figs/RP7_rho.png}
    \end{minipage}
    \hspace{0.01\linewidth}%
    \begin{minipage}{0.32\linewidth}
        \includegraphics[width=\linewidth]{figs/RP7_p.png}
    \end{minipage}
    \hspace{0.01\linewidth}%
    \begin{minipage}{0.32\linewidth}
        \includegraphics[width=\linewidth]{figs/RP7_By.png}
    \end{minipage}
    
    \caption{Numerical results for the Riemann problems RP5, RP6 and RP7 (from top to bottom) obtained using the new semi-implicit scheme compared to exact Riemann solutions (black line). The black symbols indicate results obtained by proposed IMEX scheme.
    The fluid density (left), the fluid pressure (middle) and the component of the magnetic field (right) are depicted at the chosen output time for each test.}
    \label{fig:MHDRPs2}
\end{figure}
The numerical profiles of density, pressure and transverse magnetic field
$B_y$ are shown in Figs.~\ref{fig:MHDRPs1} and~\ref{fig:MHDRPs2}. Two numerical
solutions are reported: the proposed scheme, and the standard IMEX scheme
obtained with the fixed choice $\epsM=\epsB\equiv0$, for which the explicit
sub-system is only weakly hyperbolic and is closed with the Rusanov
flux~\eqref{eq:rusanov}. % CHECK: order and reconstruction of the baseline
Both are compared with the exact solution of the Riemann problem, obtained
with the solver of Falle and Komissarov~\cite{exact1,exact2}.

In all seven cases the proposed scheme recovers the wave structure, with the
discontinuities at the correct positions and the intermediate states at the
exact levels. % CHECK against the exact solution
% CHECK (after the baseline runs): keep only what the figures show
Compared with the standard IMEX scheme, the contact and rotational
discontinuities are resolved over fewer cells, as a consequence of the
anti-diffusive terms of the HLLEM flux~\eqref{eq:hllem} and of the
reconstruction in the split characteristic variables.
For RP1 the computed solution displays the familiar compound wave, with an
accompanying overshoot in density and pressure; this feature is shared by
standard finite volume methods for ideal MHD and is found with the
HLLEM~\cite{HLLEM}, Osher-type~\cite{DumbserToro} and Rusanov fluxes
alike.

\subsection{Smooth stationary MHD vortex}\label{sec:stationary_vortex}

We now consider a stationary, isodensity variant of the vortex of Section~\ref{sec:accuracy},
run to long times following Fambri~\cite{Fambri2021}. The initial data are
\eqref{eq:vortex_profile}--\eqref{eq:vortex_dp} with $\rho_0=1$, $p_0=1$, $\tilde v=\tilde B=1$ and
$\mathbf u_0=\mathbf 0$, so that the configuration is an exact steady state of~\eqref{eq:ideal_MHD}.
The computational domain is the periodic square $[-10,10]^2$, discretised with
$200\times200$ cells, and the computation is run up to $t=1000$ with $\mathrm{CFL}=0.2$.
The flow is acoustically stiff throughout, so that the time step is set by the fluid and
Alfv\'en velocities only.

\begin{equation}\label{eq:sv_ic}
\begin{bmatrix} \rho_0 \\ u_0 \\ v_0 \\ (B_x)_0 \\ (B_y)_0 \\ p_0 \end{bmatrix}
=
\begin{bmatrix}
1 \\[2pt]
-\dfrac{\varepsilon}{2\pi}\, e^{\frac{1}{2}(1-r^2)}\, y \\[8pt]
\dfrac{\varepsilon}{2\pi}\, e^{\frac{1}{2}(1-r^2)}\, x \\[8pt]
-\dfrac{1}{2\pi}\, e^{\frac{1}{2}(1-r^2)}\, y \\[8pt]
\dfrac{1}{2\pi}\, e^{\frac{1}{2}(1-r^2)}\, x \\[8pt]
1 + \dfrac{1}{8\pi}\Bigl(\dfrac{\tilde\mu}{2\pi}\Bigr)^{2}(1-r^2)\, e^{(1-r^2)}
  - \dfrac{1}{2}\Bigl(\dfrac{\varepsilon}{2\pi}\Bigr)^{2} e^{(1-r^2)}
\end{bmatrix},
\end{equation}
with $\varepsilon=1$ and $\tilde\mu=\sqrt{4\pi}$.

Fig.~\ref{fig:MHDvortex} shows the normalised magnetic field
$\mathbf b=\mathbf B/\max\|\mathbf B_0\|$ at $t=500$ and $t=1000$, and Fig.~\ref{fig:vortex_lineout}
a cut along the $x$-axis at $t=1000$ compared with the initial (exact) profile.
The scheme preserves the stationarity of the vortex over the whole run with little
numerical dissipation, and the discrete divergence of the magnetic field remains at the
level of round-off. The results are in good agreement with those of~\cite{Fambri2021,Dematte2024}.

\begin{figure}[pos=htbp]
    \centering
    \begin{minipage}{0.48\linewidth}
        \includegraphics[width=\linewidth]{figs/MHD_500.png}
    \end{minipage}
    \hspace{0.01\linewidth}%
    \begin{minipage}{0.48\linewidth}
        \includegraphics[width=\linewidth]{figs/MHD_1000.png}
    \end{minipage}
    \caption{Long-time simulation of a smooth stationary MHD vortex. Contour plots of $\log_{10}\|\mathbf{b}\|$, overlaid with 20 equidistant contour lines between $-6.5$ and $0$, obtained with the proposed semi-implicit scheme using the unstaggered adaptation of the CT method, shown at times $t=500$ (left) and $t=1000$ (right).}
    \label{fig:MHDvortex}
\end{figure}

\begin{figure}[pos=htbp]
    \centering
    \includegraphics[width=0.6\linewidth]{figs/MHDvortex_lineout.png}
    \caption{Long-time simulation of a smooth stationary MHD vortex at time $t=1000$. One-dimensional line-out along the $x$-axis of the normalised magnetic field $\mathbf{b}$, obtained with the proposed semi-implicit algorithm using the unstaggered adaptation of the CT method described in Section~\ref{sec:CT} to enforce the divergence-free condition of the magnetic field. The solution is compared with the reference (initial) solution.}
    \label{fig:vortex_lineout}
\end{figure}

\subsection{Ideal MHD rotor problem}
\label{sec:rotor}

The rotor problem of Balsara and Spicer~\cite{BalsaraSpicer1999} combines, in
a single computation, a fast flow and a fluid at rest. A disc of heavy fluid
of radius $0.1$ spins rigidly at the centre of the domain, surrounded by a
light fluid that is initially motionless; the magnetic field threading both
is wound up by the rotation, and torsional Alfv\'en waves travel outwards and
slow the disc down. The pressure and the magnetic field are uniform at the
initial time,
%
\begin{equation}
\label{eq:rotor_ic1}
p_0 = 1, \qquad
\mathbf{B}_0 = \left(\frac{2.5}{\sqrt{4\pi}},\,0,\,0\right)^{\!\top},
\qquad w_0 = 0,
\end{equation}
%
while the density and the in-plane velocity are
%
\begin{equation}
\label{eq:rotor_ic2}
(\rho_0,\, u_0,\, v_0) =
\begin{cases}
(10,\ -10\,y,\ 10\,x), & r \le 0.1, \\[4pt]
(1,\ 0,\ 0), & r > 0.1,
\end{cases}
\qquad r = \sqrt{x^2 + y^2}.
\end{equation}
%
The computation is carried out on $[-0.5,0.5]^2$ with transmissive boundary
conditions on a uniform mesh of $1024 \times 1024$ cells, up to $t=0.25$.

The density $\rho$, the pressure $p$ and the magnetic energy $m$ at the final
time are displayed in Fig.~\ref{fig:rotor}. The flow pattern is the one
reported in the
literature~\cite{Dumbser2019,BoscheriADER,BoscheriThomann2024}.
% CHECK after the production run
The bottom row of the same figure shows the corresponding distributions of
the two splitting parameters. Inside the rotating disc the flow is supersonic
and super-Alfv\'enic and both parameters approach one, so that the whole flux
is handled by the explicit HLLEM solver; in the undisturbed outer fluid
$\epsM$ vanishes identically and $\epsB$ is small, and the acoustic flux is
implicit. In the region swept by the Alfv\'en waves the acoustic and the
Alfv\'en Mach number differ, and so do $\epsM$ and $\epsB$: this is the
situation that a single splitting parameter cannot represent. It is, moreover,
the mirror image of the blast wave of Section~\ref{sec:blast}. There the
disagreement is $\epsM\approx1$ with $\epsB\approx0$, a strongly magnetised
low-$\beta$ plasma at moderate Mach number; here it is the opposite, an
acoustically stiff but magnetically non-stiff annulus in which $\epsM$ is
near zero while $\epsB$ is close to one. The two tests therefore exercise the
two mixed regimes of Section~\ref{sec:eps_strategy} from opposite sides, and a
single parameter would misrepresent a different channel in each.
% NOTE: measured on a 256^2 run at CFL 0.2, t = 0.25: eps_M = 0.943 and
% eps_B = 0.963 on average inside r < 0.1; eps_M = 0 exactly on 80.5% of the
% cells and eps_B = 0 on 48.0%; |eps_M - eps_B| > 0.5 on 6.1% of the cells, at
% r in [0.062, 0.396], where eps_M = 0.093 and eps_B = 0.857 on average.
% Refresh from the 1024^2 production run before submission.
Throughout the computation the discrete divergence of the magnetic field stays
at the level of round-off.


\begin{figure}[pos=p]
\centering
\includegraphics[height=0.849\textheight]{figs/rotor.png}
\caption{Ideal MHD rotor problem at $t=0.25$ on a $1024\times1024$ mesh,
obtained with the proposed semi-implicit scheme. Top: density $\rho$ (left) and
gas pressure $p$ (right). Middle: squared velocity magnitude $\|\mathbf u\|^2$
(left) and magnetic energy $m=\tfrac12\|\mathbf B\|^2$ (right). Bottom: the
splitting parameters $\epsM$ (left) and $\epsB$ (right).}
\label{fig:rotor}
\end{figure}


\subsection{Ideal Orszag–Tang vortex}
\label{sec:ot}

The vortex system of Orszag and Tang~\cite{orszag1979} starts from smooth,
periodic data and develops, within a fraction of a turnover time, a set of
shocks of different strengths that cross one another and eventually leave a
field of small-scale structures typical of two-dimensional MHD turbulence; the
physics of the problem is discussed in detail
in~\cite{orszag1979,dahlburg1989,picone1991}. It therefore tests a scheme on
the formation and the interaction of shocks rather than on prescribed
discontinuities. On the periodic unit square $[0,1]^2$ the initial state is
%
\begin{align}
\label{eq:ot_ic}
\rho_0 &= \gamma^{2}, \qquad p_0 = \gamma, \qquad
\mathbf{u}_0 = \bigl(-\sin(2\pi y),\ \sin(2\pi x),\ 0\bigr)^{\!\top}, \\[4pt]
\mathbf{B}_0 &= \bigl(-\sin(2\pi y),\ \sin(4\pi x),\ 0\bigr)^{\!\top},
\end{align}
%
with $\gamma = 5/3$, and the mesh consists of $1000 \times 1000$ cells.

The gas pressure and the magnitude of the magnetic field are shown in
Fig.~\ref{fig:ot_vortex} at the times $t=0.5/2\pi$, $2/2\pi$, $3/2\pi$ and
$5/2\pi$. The position and the shape of the shocks, as well as the structures
formed after their interaction, are those found in the
literature~\cite{Dumbser2019,dumbser2008,Fambri2021,balsara2015}.
With these data the sound speed is unity and the acoustic and the Alfv\'en
Mach number are of order one over most of the domain, so that the scheme
operates mainly in its explicit regime, the splitting parameters departing
from one only in the neighbourhood of the stagnation points and of the current
sheets. The magnetic field is advanced by the unstaggered constrained
transport and its discrete divergence remains at the level of round-off.

\begin{figure}[pos=p]
\centering
\includegraphics[height=0.849\textheight]{figs/OT_grid.png}
\caption{Numerical solution for the pressure (left) and the magnetic field
magnitude (right) obtained with the proposed semi-implicit scheme for the ideal
Orszag--Tang vortex at times $t=0.5/2\pi$, $t=2.0/2\pi$, $t=3/2\pi$ and
$t=5/2\pi$ (from top to bottom).}
\label{fig:ot_vortex}
\end{figure}


\subsection{Ideal MHD blast wave problem}
\label{sec:blast}

In the blast wave problem of Balsara and Spicer~\cite{BalsaraSpicer1999} a
small region of very high pressure is released into a fluid at rest that
carries a strong uniform magnetic field. Outside the initial pocket the plasma
beta is of order $10^{-4}$, so that the internal energy is a small fraction of
the total energy, and errors in the magnetic field or in the momentum readily
produce negative pressures; the test is therefore a severe check of the
admissibility of the computed states. The data are
%
\begin{equation}
\label{eq:blast_ic1}
\rho_0 = 1, \qquad \mathbf{u}_0 = \mathbf{0},
\qquad
\mathbf{B}_0 = \frac{1}{\sqrt{4\pi}}\,(100,\,0,\,0)^{\!\top},
\end{equation}
%
\begin{equation}
\label{eq:blast_ic2}
p_0 =
\begin{cases}
1000, & r \le R, \\[2pt]
0.1, & r > R,
\end{cases}
\qquad r = \sqrt{x^2 + y^2},\quad R=0.1,
\end{equation}
%
on the square $[-0.5,0.5]^2$ with transmissive boundary conditions, and the
mesh consists of $1000 \times 1000$ cells.

Fig.~\ref{fig:blast} shows the density, the magnitude of the velocity, the gas
pressure and the magnetic pressure at $t=0.01$. The fast shock is elongated
along the direction of the field and the solution agrees with the one
of~\cite{BalsaraSpicer1999}.
% CHECK after the production run
Density and pressure remain positive throughout the computation.
The splitting parameters at the final time are reported in the bottom row of
the same figure, and this test separates them more sharply than any other
considered here. The ambient gas is at rest and strongly magnetised, so both
parameters vanish there. Inside the expanding bubble the flow is supersonic and
$\epsM$ saturates at one, while $\epsB$ stays near zero over most of the same
region, the Alfv\'en speed built on $B_0=100/\sqrt{4\pi}$ exceeding the flow
speed almost everywhere. The two therefore disagree over a wide annulus
straddling the fast shock: measured on the computed maps, $\epsM$ and $\epsB$
differ by more than one half over about two fifths of the domain, with
$\epsM\approx1$ and $\epsB\approx0$ there. A single parameter driven by either
Mach number alone would be wrong over that whole annulus --- by the acoustic
Mach number it would leave the magnetic flux explicit and the time step tied to
the Alfv\'en speed, and by the Alfv\'en Mach number it would send the
supersonic acoustic flux to the central-type implicit operator.
% NOTE: the two-fifths figure and the eps statistics were measured on a 256^2
% run at CFL 0.2 (eps_M = 0 on 46.0% and = 1 on 38.5% of the cells; eps_B = 0 on
% 77.7% and = 1 on 2.6%; |eps_M - eps_B| > 0.5 on 41.3%).  Refresh them from the
% 1000^2 production run before submission.
The magnetic field is advanced by the unstaggered constrained transport and
its discrete divergence remains at the level of round-off.

\begin{figure}[pos=p]
\centering
\includegraphics[height=0.849\textheight]{figs/blast.png}
\caption{Ideal MHD blast wave problem at $t=0.01$ on a $1000\times1000$ mesh,
obtained with the proposed semi-implicit scheme. Top: $\log_{10}\rho$ (left) and
$\log_{10}p$ (right). Middle: squared velocity magnitude $\|\mathbf u\|^2$
(left) and magnetic energy $m=\tfrac12\|\mathbf B\|^2$ (right). Bottom: the
splitting parameters $\epsM$ (left) and $\epsB$ (right).}
\label{fig:blast}
\end{figure}

\subsection{Field loop advection}
\label{sec:fieldloop}

The advection of a weak magnetic field loop, proposed by Gardiner and
Stone~\cite{GardinerStone2005}, is run here in the low Mach number setting
of~\cite{BoscheriThomann2024}. A fluid of density $\rho=1$ and pressure
$p=10^{5}$ moves with the uniform velocity $(u,v,w)=(2,1,0)$ across the
periodic domain $\Omega=[-1,1]\times[-0.5,0.5]$, which is discretised with
$512\times256$ cells. The magnetic field is the curl of the vector potential
%
\begin{equation}
\label{eq:loop_ic}
A_z =
\begin{cases}
\dfrac{A_0}{\sqrt{4\pi}}\,(R-r), & r \le R,\\[2mm]
0, & r > R,
\end{cases}
\qquad r=\sqrt{x^2+y^2},\quad R=0.3,
\end{equation}
%


so that $\|\Bb\|=A_0/\sqrt{4\pi}$ is constant inside the loop and vanishes
outside; the field is discontinuous at $r=R$ and its direction is undefined at
the centre. Four amplitudes are considered, $A_0\in\{10^{-3},10^{-2},10^{-1},1\}$, and the
final time is $t_f=1$, after one period in $x$ and in $y$.

The acoustic Mach number is $M\approx5.5\times10^{-3}$ in all four runs,
whereas the Alfv\'en Mach number decreases from $M_A\approx7.9\times10^{3}$ to
$M_A\approx7.9$ as the amplitude is increased. The flow is thus acoustically
stiff and magnetically non-stiff, one of the two mixed regimes identified in
Section~\ref{sec:eps_strategy}: the splitting returns $\epsM\approx0$ and
$\epsB=1$, so that the acoustic flux is implicit while the magnetic flux is
advanced explicitly by the HLLEM flux of Section~\ref{sec:hllem} instead of
being carried, together with the acoustics, by the central-type implicit
operator. In two dimensions this concerns the Lorentz force of the momentum
equation, the magnetic energy flux and $B_z$; the in-plane field is advanced by
the constrained transport of Section~\ref{sec:CT} at either value of $\epsB$.
The computed maps confirm the two parameters at the ends of their ranges, with
$\epsM$ at zero and $\epsB$ at one over the whole domain.

The magnitude of the magnetic field at the final time is shown in
Fig.~\ref{fig:fieldloop}. The loop keeps its shape over the whole range of
amplitudes. % CHECK after the runs
The discrete divergence of the magnetic field, measured in the maximum norm at
every time step, remains at the level of round-off in all four runs. The
quantity to compare it against is $\max\|\Bb\|/h$, since that is the size of
each term of the difference quotient: at $A_0=1$ the largest value of
$\|\nabla_4\!\cdot\!\Bb\|_\infty$ over the whole computation is
$4.7\times10^{-14}$ against $\max\|\Bb\|/h=1.8\times10^{1}$, a relative level
of $2.5\times10^{-15}$.
% NOTE: measured on a 128x64 run at A_0 = 1.  The other three amplitudes need
% the production runs; the relative level is the quantity to quote.


\begin{figure}[pos=htbp]
    \centering
    \includegraphics[width=\linewidth]{figs/Field_Bmag.png}
    \caption{Magnetic field loop advection at $t_f=1$. Magnitude of the magnetic
field $\|\mathbf{B}\|$ obtained with the proposed semi-implicit scheme for four
increasing loop amplitudes, from $A_0=10^{-3}$ (top) to $A_0=1$ (bottom).}
    \label{fig:fieldloop}
\end{figure}


\subsection{Magnetized Kelvin–Helmholtz instability}
\label{sec:KH}
The last test follows the low Mach number Kelvin--Helmholtz configuration
of~\cite{BoscheriThomann2024}, in which a weak magnetic field is aligned with a
double shear layer and the flow speed is controlled by a single parameter
$M_x$, the acoustic Mach number of the horizontal stream. On the periodic
domain $\Omega=[0,2]\times[-0.5,0.5]$, discretised with $512\times256$ cells,
the density and the pressure are uniform, $\rho=\gamma=1.4$ and $p=1$, so that
the sound speed is unity, and the velocity is
%
\begin{equation}
\label{eq:kh_vel}
u = M_x\,\bigl(1 - 2\,\eta(y)\bigr), \qquad
v = 0.1\,M_x\,\sin(2\pi x), \qquad w=0,
\end{equation}
%
with the shear profile
%
\begin{equation}
\label{eq:kh_eta}
\eta(y) =
\begin{cases}
\tfrac12\bigl(1 + \sin(16\pi(y+0.25))\bigr), & -\tfrac{9}{32} \le y < -\tfrac{7}{32}, \\[4pt]
1, & -\tfrac{7}{32} \le y < \tfrac{7}{32}, \\[4pt]
\tfrac12\bigl(1 - \sin(16\pi(y-0.25))\bigr), & \tfrac{7}{32} \le y < \tfrac{9}{32}, \\[4pt]
0, & \text{otherwise}.
\end{cases}
\end{equation}
%
The magnetic field is uniform, $\Bb=(0.1\,M_x,0,0)^{\mathsf T}$, which gives
an Alfv\'en Mach number $M_A\approx11.8$ in the free streams, independently of
$M_x$. The values $M_x=10^{-1},10^{-2},10^{-3},10^{-4}$ are considered, each
computation being stopped at $t_f=0.8/M_x$, i.e.\ at one sixth of the time
$4.8/M_x$ used in~\cite{BoscheriThomann2024}; up to this time the roll-up of
the shear layer is smooth and well resolved.

In this configuration the two splitting parameters fall on opposite sides of
the switch~\eqref{eq:switch}: since $M\le M_x\le0.1<M_0$, one has
$\epsM\approx0$ and the acoustic flux is implicit, whereas $M_A>M_1$ in the
streams gives $\epsB=1$, so that the magnetic flux is advanced by the explicit
HLLEM solver. The explicit sub-system then propagates signals at speeds
proportional to $M_x$, and so does the numerical dissipation of the HLLEM
flux. Fig.~\ref{fig:kh_mach} shows the local Mach number at the final time
divided by $M_x$; the four profiles coincide. % CHECK after the runs
This is the behaviour expected from an asymptotic-preserving scheme: an
explicit discretisation with a dissipation set by the sound speed would damp
the vortical structures increasingly as $M_x$ decreases, as reported
in~\cite{BoscheriThomann2024}.

The ratio between the time step of the scheme and that of a fully explicit
discretisation, $\Delta t_{\mathrm{MHD}}$, computed with the fast
magnetosonic speed, is shown in Fig.~\ref{fig:kh_dtratio}. It grows in
inverse proportion to $M_x$ and exceeds four orders of magnitude for
$M_x=10^{-4}$. % CHECK the values over the whole run

\begin{figure}[pos=htbp]
    \centering
    \includegraphics[width=\linewidth]{figs/KH_Mach.png}
    \caption{Magnetized Kelvin--Helmholtz instability at $t_f = t_{\max}/6$.
    Acoustic Mach number rescaled by the prescribed value $M_x$, for increasing
    $M_x$ from top left to bottom right. The rescaled profiles are almost
    identical across the whole range, showing that the proposed scheme avoids
    excessive numerical dissipation at low acoustic Mach number.}
    \label{fig:kh_mach}
\end{figure}

\begin{figure}[pos=htbp]
    \centering
    \includegraphics[width=0.6\linewidth]{figs/dt_ratio_KH.png}
    \caption{Magnetized Kelvin--Helmholtz instability. Time evolution of the
    ratio between the material time step of the proposed semi-implicit scheme and
    the magnetosonic time step of a fully explicit method, for the different
    values of $M_x$. The ratio reaches up to four orders of magnitude as $M_x$
    decreases.}
    \label{fig:kh_dtratio}
\end{figure}



\FloatBarrier % all Section 4 floats are output before the conclusions
\section{Conclusions}
\label{sec:Conclusion}

We have presented an implicit--explicit finite volume scheme for the ideal MHD
equations that is accurate and efficient across all Mach and Alfv\'en number
regimes. It is built on the three-way flux splitting of
Fambri~\cite{Fambri2021}, but departs from it in how the parts of the splitting
are assigned: the fractions of the acoustic and of the magnetic flux that are
treated explicitly are governed by two independent local parameters, $\epsM$ and
$\epsB$, driven by the acoustic and the Alfv\'en Mach number. Since the two
stiffness mechanisms of compressible MHD are controlled by distinct
dimensionless numbers, each is thereby transferred to the implicit solver
independently of the other.

The first contribution is the characteristic structure of this splitting, and it
has two halves that serve different parts of the scheme. The Alfv\'en speeds
follow from a quadratic in $\epsB$ alone and the magnetosonic speeds from a
quartic in both parameters, with real roots at every value of the two
parameters; that is what makes the explicit CFL condition and the wave-speed
envelope of the Riemann solver well defined. Whenever in addition $\epsM>0$ the
explicit sub-system is strongly hyperbolic, and the resulting complete basis is
what allows the characteristic reconstruction to be carried out in the split
basis---bounded away from $\epsM=0$, where the basis degenerates while the flux
does not. Independently of both, and at every value of the parameters, an exact block decomposition confines the singular limits of the
eigenvectors to sub-blocks treated in closed form: the divergence mode is
decoupled, and the Alfv\'en fields span a two-dimensional invariant subspace
whose dual functionals annihilate every other wave family. It is this second half
that makes a wave-resolving Riemann solver possible in the stiff limit, where the
complete basis is lost and where the standard three-way splitting, whose explicit
spectrum collapses onto $\{0,u\}$, retains no trace of the rotational
discontinuities at all. We use HLLEM in a $\delta_k$-driven form whose wave-speed
envelope is built from the split spectrum itself.

The second contribution concerns the implicit step. The kinetic and magnetic
energies are linearised about the explicit argument of the partitioned IMEX
Runge--Kutta scheme, and the unstaggered, potential-based constrained transport
is written as a linear map of the momentum; both are embedded in the pressure
operator. Each stage then solves for a magnetic predictor and then for the
pressure, without an outer iteration between the two, the latter being balanced
with a divergence-free field throughout its solution. The total energy is
conserved without a second pressure solve, and the third order of the time
integrator is retained. Positivity of the pressure is enforced with the
mechanism of~\cite{MHDIMEX}, whose modified pressure relation makes the system
nonlinear and is solved by a nested Newton iteration; since that mechanism needs
an M-matrix, it is applied to a second-order operator and the fourth-order
solution is recovered by defect correction, so that the nested Newton iteration
runs inside every defect step. Iteration is thus not avoided; what is avoided
is the outer iteration between the magnetic field and the pressure.

% CHECK every sentence of this paragraph against Section 4 and delete what the
% experiments do not show.
The scheme has been assessed on one- and two-dimensional benchmarks ranging from
the near-incompressible regime to supersonic and strongly magnetised flows.
Convergence studies on a travelling MHD vortex confirm the design third order
for all variables, in the low-Mach and in the low-Alfv\'en limit alike; the
measured orders approach three from below at unit background density and
exceed it in the stiffest cases, where the coarse grids are still
pre-asymptotic.
The time step is governed by the material velocity wherever the corresponding
Mach number is small, independently for the two channels: in the
Kelvin--Helmholtz test the measured ratio to the magnetosonic step of a fully
explicit discretisation grows in inverse proportion to $M_x$ and exceeds four
orders of magnitude at $M_x=10^{-4}$ (Fig.~\ref{fig:kh_dtratio}). Where
neither Mach number is small the gain disappears by construction, as in the
blast wave, whose global step is set by the supersonic interior of the bubble.
The gain of the second parameter over a single one is likewise
regime-dependent, and lies with the strongly magnetised rather than with the
high-$\beta$ mixed regime.
Shock-dominated problems are computed without tuning of the parameters, and the
magnetic field remains divergence-free to machine precision in the discrete
sense of Section~\ref{sec:CT}. The identity
$\nabla_4\!\cdot\!\mathrm{curl}_4\equiv0$ needs the full five-point stencil, so
the two cells next to a non-periodic boundary are excluded from the norms; this
excludes nothing of substance, since the maximum of
$\|\nabla_4\!\cdot\!\Bb\|_\infty$ is attained at the magnetic structure itself
and not at the boundary. Relative to $\max\|\Bb\|/h$ it stays at
$1.1\times10^{-14}$ in the blast wave and at $1.6\times10^{-15}$ in the
field-loop advection.

Several directions follow. The argument advanced here, that independent
stiffness mechanisms call for independent splitting parameters, is not specific
to two: Hall MHD supplies a natural third instance, its whistler branch
introducing a time-step restriction that scales as $\Delta x^{2}$ and is governed
by its own dimensionless number. A second direction concerns the switching
function. Both $\eps$ and the shock indicator are driven by the velocity field
and are therefore blind to a large pressure jump in a fluid at rest, as at the
initial interface of the blast wave; coupling the parameters to a measure of
departure from equilibrium is a question we intend to pursue. On the analytical
side, the convergence of the defect correction and the stability of the scheme
at large Alfv\'en CFL numbers are supported here by numerical evidence only.
Finally, the extension to three space dimensions is not a transcription: the
scalar potential $A_z$ of Section~\ref{sec:CT} becomes a vector potential, with
the gauge freedom that entails, and the linear map of the momentum has to be
rebuilt accordingly. Combining the discrete divergence-free property with the
positivity framework of Wu and Shu~\cite{WuShu2018,WuShu2019} may allow the constructive
admissibility mechanisms of Section~\ref{sec:zslimiter} to be replaced by a
provable bound on the evolved cell averages.
% Numbered list
% Use the style of numbering in square brackets.
% If nothing is used, default style will be taken.
%\begin{enumerate}[a)]
%\item 
%\item 
%\item 
%\end{enumerate}  

% Unnumbered list
%\begin{itemize}
%\item 
%\item 
%\item 
%\end{itemize}  

% Description list
%\begin{description}
%\item[]
%\item[] 
%\item[] 
%\end{description}  



% Uncomment and use as the case may be
%\begin{theorem} 
%\end{theorem}

% Uncomment and use as the case may be
%\begin{lemma} 
%\end{lemma}

%% The Appendices part is started with the command \appendix;
%% appendix sections are then done as normal sections
%% \appendix



% To print the credit authorship contribution details


%% Loading bibliography style file
%\bibliographystyle{model1-num-names}

% Loading bibliography database
\bibliographystyle{elsarticle-num}
% Loading bibliography database
\bibliography{cas-refs-capitalised}

% Biography
%\bio{}
% Here goes the biography details.
%\endbio

%\bio{pic1}
% Here goes the biography details.
%\endbio

\end{document}

